\documentclass[11pt,a4paper]{report}
\usepackage[margin=2.6cm]{geometry}
\usepackage{amsmath,amssymb,amsthm,mathtools,bm}
\usepackage{enumitem}
\usepackage{booktabs,array}
\usepackage{tikz}
\usetikzlibrary{arrows.meta,positioning,calc,shapes.geometric,decorations.pathreplacing,fit,backgrounds}
\usepackage[most]{tcolorbox}
\usepackage[colorlinks=true,linkcolor=blue!50!black,citecolor=green!40!black,urlcolor=blue!50!black]{hyperref}
\usepackage{cleveref}
\usepackage[expansion=false]{microtype}
\setlength{\parskip}{0.4em}
\setlist{itemsep=0.15em,topsep=0.3em}

% ---------- theorem environments ----------
\theoremstyle{plain}
\newtheorem{theorem}{Theorem}[chapter]
\newtheorem{lemma}[theorem]{Lemma}
\newtheorem{proposition}[theorem]{Proposition}
\newtheorem{corollary}[theorem]{Corollary}
\newtheorem{claim}[theorem]{Claim}
\theoremstyle{definition}
\newtheorem{definition}[theorem]{Definition}
\newtheorem{example}[theorem]{Example}
\newtheorem{exercise}[theorem]{Exercise}
\theoremstyle{remark}
\newtheorem{remark}[theorem]{Remark}
\crefname{theorem}{Theorem}{Theorems}
\crefname{lemma}{Lemma}{Lemmas}
\crefname{definition}{Definition}{Definitions}
\crefname{example}{Example}{Examples}
\crefname{claim}{Claim}{Claims}

% ---------- boxes ----------
\tcbset{protostyle/.style={enhanced,breakable,colback=blue!3,colframe=blue!50!black,fonttitle=\bfseries}}
\newtcolorbox[auto counter,number within=chapter]{protocol}[2][]{protostyle,title={Protocol~\thetcbcounter: #2},#1}
\newtcolorbox[use counter from=protocol]{algorithm}[2][]{protostyle,title={Algorithm~\thetcbcounter: #2},#1}
\newtcolorbox{intuition}[1][Intuition]{enhanced,breakable,colback=yellow!8,colframe=orange!70!black,
  fonttitle=\bfseries,title={#1},left=4pt,right=4pt}
\newtcolorbox{warning}[1][Pitfall]{enhanced,breakable,colback=red!4,colframe=red!60!black,
  fonttitle=\bfseries,title={#1}}

% ---------- macros ----------
\newcommand{\NN}{\mathbb{N}}\newcommand{\ZZ}{\mathbb{Z}}\newcommand{\FF}{\mathbb{F}}
\newcommand{\GG}{\mathbb{G}}
\newcommand{\getsr}{\xleftarrow{\$}}
\newcommand{\negl}{\mathsf{negl}}
\newcommand{\poly}{\mathrm{poly}}
\newcommand{\polylog}{\mathrm{polylog}}
\newcommand{\PPT}{\mathrm{PPT}}
\renewcommand{\Pr}{\mathrm{Pr}}
\newcommand{\Adv}{\mathsf{Adv}}
\newcommand{\Prv}{\mathsf{P}}\newcommand{\Ver}{\mathsf{V}}
\newcommand{\Sim}{\mathsf{Sim}}\newcommand{\Ext}{\mathsf{Ext}}
\newcommand{\view}{\mathsf{View}}
\newcommand{\out}{\mathsf{out}}
\newcommand{\NP}{\mathsf{NP}}\newcommand{\BPP}{\mathsf{BPP}}\newcommand{\IP}{\mathsf{IP}}
\newcommand{\PSPACE}{\mathsf{PSPACE}}\newcommand{\AM}{\mathsf{AM}}\newcommand{\coAM}{\mathsf{coAM}}
\newcommand{\SZK}{\mathsf{SZK}}\newcommand{\CZK}{\mathsf{CZK}}\newcommand{\PZK}{\mathsf{PZK}}
\newcommand{\HVSZK}{\mathsf{HVSZK}}
\newcommand{\coNP}{\mathsf{coNP}}\newcommand{\PCP}{\mathsf{PCP}}
\newcommand{\Setup}{\mathsf{Setup}}\newcommand{\Com}{\mathsf{Com}}
\newcommand{\Open}{\mathsf{Open}}
\newcommand{\SD}{\mathsf{SD}}
\newcommand{\pair}[2]{\left\langle #1,#2\right\rangle}
\newcommand{\ra}[1]{\xrightarrow{\makebox[2.5cm]{$#1$}}}
\newcommand{\la}[1]{\xleftarrow{\makebox[2.5cm]{$#1$}}}
\newcommand{\sample}{\leftarrow}
\newcommand{\inv}{^{-1}}
\newcommand{\vc}[1]{\mathbf{#1}}
\newcommand{\ip}[2]{\langle #1, #2\rangle}
\DeclareMathOperator{\GI}{GI}\DeclareMathOperator{\GNI}{GNI}
\DeclareMathOperator{\QR}{QR}\DeclareMathOperator{\QNR}{QNR}
\DeclareMathOperator{\dlog}{dlog}
\DeclareMathOperator{\ord}{ord}
\DeclareMathOperator{\Aut}{Aut}
\newcommand{\Z}[1]{\ZZ_{#1}}
\newcommand{\Zs}[1]{\ZZ_{#1}^{*}}

% a two-party transcript table: three columns P | channel | V
\newenvironment{flow}[2]{%
  \par\smallskip\noindent
  \begin{tabular}{@{}p{0.34\linewidth} c p{0.34\linewidth}@{}}
  \textbf{#1} & & \textbf{#2}\\ \hline}{\end{tabular}\par\smallskip}

\title{\Huge\bfseries Zero-Knowledge Proofs\\[4pt]\Large From Foundations to zk-SNARKs and STARKs\\[10pt]
\large Graduate Lecture Notes (MPhil / PhD level)}
\author{}
\date{\today}

\begin{document}
\maketitle
\begin{abstract}\noindent
These notes give a self-contained, proof-oriented treatment of zero-knowledge proofs.
We start from the definitions of interactive proofs, simulation-based security and proofs of knowledge;
develop the classical protocols (graph isomorphism, quadratic residuosity, graph 3-colouring) with full
proofs; study $\Sigma$-protocols and their algebra (AND/OR/equality composition, Pedersen commitments);
treat the Fiat--Shamir transform and the forking lemma; and then move on to the modern
\emph{succinct} proof systems: sum-check, GKR, arithmetisation (R1CS/QAP), polynomial commitments
(KZG, FRI, inner-product arguments), Groth16, PLONK and STARKs. Every algorithm is accompanied by a
worked numerical example (small enough to be verified by hand) and by a diagram.
Chapters end with exercises. Starred sections are optional on a first reading.
\end{abstract}
\tableofcontents

\chapter{Introduction and Motivation}
\section{What is a zero-knowledge proof?}
A \emph{proof} in the classical sense (a written derivation, or an $\NP$ witness) convinces the reader of a statement
but also \emph{teaches} the reader something: the reader obtains the derivation and can now convince others.
Goldwasser, Micali and Rackoff (1985) asked a startling question:
\begin{quote}\itshape
Can a prover convince a verifier that a statement is true while revealing \emph{nothing} except its truth?
\end{quote}
The answer, perhaps surprisingly, is \emph{yes}, provided we relax the notion of proof in two ways:
the proof is \textbf{interactive} (the verifier talks to the prover) and \textbf{probabilistic} (the verifier may
err with negligible probability). Three properties characterise a zero-knowledge proof for a statement $x$:
\begin{description}
\item[Completeness.] If $x$ is true and both parties follow the protocol, the verifier accepts.
\item[Soundness.] If $x$ is false, no (possibly cheating, possibly unbounded) prover can make the verifier accept, except with small probability.
\item[Zero-knowledge.] If $x$ is true, whatever a (possibly cheating) verifier can compute after the interaction, it could have computed \emph{without} the interaction, from $x$ alone. Formalised by a \emph{simulator}.
\end{description}

\subsection{The Ali Baba cave}
Quisquater et al.\ gave the following popular story. A circular cave has an entrance and, at the far end, a magic door that opens with a secret word. Peggy (prover) claims to know the word; Victor (verifier) wants to be convinced without learning the word.
\begin{center}
\begin{tikzpicture}[scale=1.15,every node/.style={font=\small}]
  % ring-shaped corridor
  \draw[thick,fill=gray!12] (0,0) circle (1.7); \draw[thick,fill=white] (0,0) circle (0.8);
  % entrance corridor at the bottom
  \draw[thick,fill=gray!12] (-0.3,-1.7) rectangle (0.3,-2.3);
  \node at (0,-2.65) {entrance};
  % magic door blocks the ring at the top
  \draw[line width=3pt,red] (0,0.8)--(0,1.7); \node[red] at (0,2.0) {magic door};
  \node at (-1.25,0) {path $A$}; \node at (1.25,0) {path $B$};
  \draw[->,thick,blue] (0,-2.3)--(0,-1.75);
  \draw[->,dashed,blue] (-0.1,-1.45) to[bend left=35] (-1.25,-0.5);
  \draw[->,dashed,blue] (0.1,-1.45) to[bend right=35] (1.25,-0.5);
\end{tikzpicture}
\end{center}
\begin{enumerate}
\item Peggy enters and walks randomly down path $A$ or $B$ while Victor waits outside (he cannot see which).
\item Victor walks to the entrance and shouts ``come out via $A$'' or ``via $B$'' (a fair coin).
\item Peggy complies, using the door if necessary.
\end{enumerate}
If Peggy knows the secret she always succeeds. If she does not, she succeeds only if she guessed Victor's coin, probability $1/2$. After $k$ independent rounds a cheater survives with probability $2^{-k}$. Zero-knowledge: Victor sees only ``Peggy came out of the side I asked for'', which he could have staged with a friend (the \emph{simulator}) who simply hides in the correct side after seeing his request (rewinding!). Victor cannot even convince third parties, because a video of the exchange is equally producible by a fraudster.

\begin{intuition}[The three ideas behind every ZK protocol]
\textbf{(1) Commit-challenge-response:} the prover commits to something random; the verifier asks one of several questions; the prover can answer \emph{every} question only if she knows the witness, but any \emph{single} answer is independent of the witness.
\textbf{(2) Special soundness:} two answers to two different questions under the \emph{same} commitment reveal the witness. Hence a prover who could answer both ``knows'' the witness.
\textbf{(3) Simulation by rewinding/programming:} a simulator that knows the question in advance can fabricate a convincing transcript without the witness.
\end{intuition}

\section{Why do we care?}
\begin{itemize}
\item \textbf{Authentication}: prove knowledge of a password/secret key without revealing it (Schnorr, Fiat--Shamir identification).
\item \textbf{Signatures}: Fiat--Shamir converts ZK identification protocols into signature schemes.
\item \textbf{Cryptographic protocol design}: GMW compiler---any protocol secure against honest-but-curious parties becomes secure against malicious ones by having each party prove in ZK that it followed the protocol.
\item \textbf{Blockchains}: Zcash (shielded transactions), zk-rollups (verify $10^6$ transactions with one short proof), anonymous credentials, verifiable computation.
\item \textbf{Complexity theory}: ZK led to $\IP=\PSPACE$, PCP theorem, hardness of approximation.
\end{itemize}

\section{A map of the territory}\label{sec:map}
\begin{center}
\begin{tikzpicture}[node distance=1.0cm and 1.3cm, every node/.style={font=\small},
 box/.style={draw,rounded corners,align=center,minimum height=0.9cm,fill=blue!6,inner sep=4pt}]
 \node[box] (ip) {Interactive proofs\\(Ch.~\ref{ch:ip})};
 \node[box,right=of ip] (zk) {Zero-knowledge\\definitions (Ch.~\ref{ch:zk})};
 \node[box,right=of zk] (cl) {Classical protocols\\GI, QR, 3-COL (Ch.~\ref{ch:classical}--\ref{ch:np})};
 \node[box,below=of zk] (sig) {$\Sigma$-protocols\\(Ch.~\ref{ch:sigma})};
 \node[box,left=of sig] (fs) {Fiat--Shamir,\\NIZK (Ch.~\ref{ch:fs})};
 \node[box,below=of sig] (sc) {Sum-check, GKR\\(Ch.~\ref{ch:sumcheck})};
 \node[box,left=of sc] (ar) {Arithmetisation\\R1CS/QAP (Ch.~\ref{ch:arith})};
 \node[box,below=of sc] (pc) {Polynomial commitments\\KZG, FRI, IPA (Ch.~\ref{ch:pcs})};
 \node[box,left=of pc] (sn) {SNARKs/STARKs\\Groth16, PLONK (Ch.~\ref{ch:snark})};
 \node[box,right=of pc] (th) {Limits \& complexity\\(Ch.~\ref{ch:theory})};
 \draw[-{Stealth}] (ip)--(zk); \draw[-{Stealth}] (zk)--(cl); \draw[-{Stealth}] (zk)--(sig); \draw[-{Stealth}] (sig)--(fs);
 \draw[-{Stealth}] (sig)--(sc); \draw[-{Stealth}] (sc)--(ar);\draw[-{Stealth}] (sc)--(pc); \draw[-{Stealth}] (ar)--(sn); \draw[-{Stealth}] (pc)--(sn);
 \draw[-{Stealth}] (cl.south) -- ++(0,-0.4) -| (th.north);
\end{tikzpicture}
\end{center}

\section{Conventions}
Security parameter $\lambda$; $1^\lambda$ is its unary encoding. $x\getsr S$ means sampling uniformly from the finite set $S$.
$[n]=\{1,\dots,n\}$. For a finite group $\GG$ of prime order $q$ with generator $g$ we write exponents in $\ZZ_q$.
Algorithms are probabilistic; ``PPT'' means probabilistic polynomial time in $\lambda$ (or in $|x|$ when $x$ is the input).

\chapter{Mathematical Preliminaries}\label{ch:prelim}
\section{Asymptotics, negligible functions and indistinguishability}
\begin{definition}[Negligible]
$\mu:\NN\to\mathbb{R}_{\ge0}$ is \emph{negligible} if for every polynomial $p$ there is $N$ such that $\mu(\lambda)<1/p(\lambda)$ for all $\lambda>N$. We write $\mu=\negl(\lambda)$. A function is \emph{overwhelming} if $1-\mu$ is negligible.
\end{definition}
Negligible functions are closed under addition and under multiplication by polynomials: $\poly\cdot\negl=\negl$. This is exactly what is needed so that an event of negligible probability cannot be amplified by polynomially many repetitions.

\begin{definition}[Ensembles, statistical distance]
A \emph{probability ensemble} $X=\{X_{a}\}_{a\in S}$ indexed by $S\subseteq\{0,1\}^*$ is a family of random variables. For distributions $D_0,D_1$ on a finite set $\Omega$,
\[\SD(D_0,D_1)=\max_{T\subseteq\Omega}\big|D_0(T)-D_1(T)\big|=\tfrac12\sum_{\omega}|D_0(\omega)-D_1(\omega)|.\]
\end{definition}
\begin{definition}[Indistinguishability]
Two ensembles $X=\{X_a\}$, $Y=\{Y_a\}$ are
\begin{itemize}
\item \emph{perfectly indistinguishable} ($X\equiv Y$) if $X_a$ and $Y_a$ have identical distributions for every $a$;
\item \emph{statistically indistinguishable} ($X\approx_s Y$) if $\SD(X_a,Y_a)=\negl(|a|)$;
\item \emph{computationally indistinguishable} ($X\approx_c Y$) if for every non-uniform PPT distinguisher $D$ the advantage $\big|\Pr[D(a,X_a)=1]-\Pr[D(a,Y_a)=1]\big|=\negl(|a|)$.
\end{itemize}
\end{definition}
\begin{lemma}[Hybrid argument]\label{lem:hybrid}
Let $X^{(0)},\dots,X^{(m)}$ be distributions, $m=\poly$. If $D$ distinguishes $X^{(0)}$ from $X^{(m)}$ with advantage $\varepsilon$, then for some $i<m$ it distinguishes $X^{(i)}$ from $X^{(i+1)}$ with advantage at least $\varepsilon/m$.
\end{lemma}
\begin{proof}
Let $p_i=\Pr[D(X^{(i)})=1]$. Then $\varepsilon\le|p_0-p_m|=\big|\sum_{i<m}(p_i-p_{i+1})\big|\le\sum_{i<m}|p_i-p_{i+1}|$, so some term is at least $\varepsilon/m$.
\end{proof}
The hybrid argument is the workhorse for every ZK proof with polynomially many commitments (e.g.\ Chapter~\ref{ch:np}). Its constructive form: a distinguisher for the endpoints yields a distinguisher for an adjacent pair that you can \emph{embed a challenge into}, by guessing $i$ with probability $1/m$.

\section{Languages, relations and $\NP$}
\begin{definition}
A binary relation $R\subseteq\{0,1\}^*\times\{0,1\}^*$ is an \emph{$\NP$-relation} if membership $(x,w)\in R$ is decidable in time $\poly(|x|)$ and there is a polynomial $p$ with $(x,w)\in R\Rightarrow|w|\le p(|x|)$. Its language is $L_R=\{x:\exists w,(x,w)\in R\}$; $w$ is a \emph{witness} for $x$. We write $R(x)$ for the set of witnesses.
\end{definition}
\begin{example}\label{ex:relations}
\begin{itemize}
\item \textbf{Discrete log}: $R_{\mathrm{DL}}=\{((\GG,g,y),x): g^x=y\}$.
\item \textbf{Graph isomorphism}: $R_{\mathrm{GI}}=\{((G_0,G_1),\pi):\pi(G_0)=G_1\}$.
\item \textbf{Quadratic residuosity}: $R_{\mathrm{QR}}=\{((N,v),s):s^2\equiv v\pmod N\}$.
\item \textbf{Circuit satisfiability}: $R_{\mathrm{SAT}}=\{(C,w):C(w)=1\}$, $\NP$-complete.
\item \textbf{Hamiltonian cycle}, \textbf{3-colourability}: both $\NP$-complete.
\end{itemize}
\end{example}
Proofs \emph{of membership} ($x\in L$) and proofs \emph{of knowledge} ($\Prv$ knows $w$ with $(x,w)\in R$) are different goals; Chapter~\ref{ch:zk} treats the difference.

\section{Algebraic background}
\subsection{Cyclic groups and discrete logarithm assumptions}
Let $\GG=\langle g\rangle$ be cyclic of prime order $q$. Throughout we use a toy group for numerical examples:
\begin{center}
\fbox{\parbox{0.85\linewidth}{\centering\textbf{Toy group.} $p=23$, $q=11$, $\GG=\langle 2\rangle\subset\Zs{23}$, $g=2$. \\
Indeed $2^{11}=2048=89\cdot23+1\equiv 1$, so $\ord(2)=11$. $\GG$ is the set of squares mod $23$: $\{1,2,3,4,6,8,9,12,13,16,18\}$.}}
\end{center}
Powers of $2$ mod $23$ (the \emph{discrete-log table}, handy for all later examples):
\begin{center}\small
\begin{tabular}{c|ccccccccccc}
$k$ & 0&1&2&3&4&5&6&7&8&9&10\\\hline
$2^k \bmod 23$ & 1&2&4&8&16&9&18&13&3&6&12
\end{tabular}
\end{center}
\begin{definition}[Discrete logarithm, CDH, DDH]
For a group $\GG$ of prime order $q$ with generator $g$:
\begin{itemize}
\item \textbf{DL}: given $g^a$ ($a\getsr\ZZ_q$) find $a$.
\item \textbf{CDH}: given $(g^a,g^b)$ find $g^{ab}$.
\item \textbf{DDH}: distinguish $(g^a,g^b,g^{ab})$ from $(g^a,g^b,g^c)$ with $a,b,c\getsr\ZZ_q$.
\end{itemize}
Each assumption states that every PPT algorithm succeeds with negligible advantage. The assumptions get stronger from DL to DDH: if DDH is hard then CDH is hard, and if CDH is hard then DL is hard (an algorithm solving DL immediately solves CDH, and one solving CDH immediately distinguishes DDH tuples).
\end{definition}
DDH is \emph{false} in the toy group and in $\Zs p$ itself (Legendre symbol), but holds conjecturally in prime-order subgroups of $\Zs p$ and in elliptic curve groups.

\subsection{Quadratic residues and the Blum integer setting}
Let $N=PQ$ with $P,Q$ distinct odd primes. $\QR_N=\{a\in\Zs N:\exists s,\,s^2\equiv a\}$. The Jacobi symbol $\left(\frac aN\right)=\left(\frac aP\right)\left(\frac aQ\right)$ is efficiently computable without the factorisation. Let $J_N^{+1}=\{a:(\frac aN)=1\}$; then $\QR_N\subsetneq J_N^{+1}$ and $|J_N^{+1}|=2|\QR_N|$.
\begin{definition}[Quadratic residuosity assumption (QRA)]
Given $N$ (of unknown factorisation) and a uniform $a\in J_N^{+1}$, deciding whether $a\in\QR_N$ is hard for PPT algorithms.
\end{definition}
\begin{lemma}[Square roots $\Leftrightarrow$ factoring]\label{lem:sqrtfactor}
Given $N=PQ$ and two square roots $s,s'$ of $a\in\QR_N$ with $s'\not\equiv\pm s$, one can factor $N$.
\end{lemma}
\begin{proof}
$s^2\equiv s'^2$ gives $N\mid(s-s')(s+s')$ but $N\nmid(s\pm s')$; hence $\gcd(s-s',N)\in\{P,Q\}$.
\end{proof}
Every $a\in\QR_N$ has exactly four roots $\pm s_1,\pm s_2$ (CRT). \textbf{Toy modulus:} $N=77=7\cdot11$.

\subsection{Pairings}
A (type-3) \emph{bilinear pairing} is an efficiently computable map $e:\GG_1\times\GG_2\to\GG_T$ between groups of prime order $q$, with generators $g_1,g_2$ and $e(g_1,g_2)\ne1$, satisfying $e(g_1^a,g_2^b)=e(g_1,g_2)^{ab}$. We use additive-bracket notation $[a]_1=g_1^a$, $[a]_2=g_2^a$, $[a]_T=e(g_1,g_2)^a$ so that $e([a]_1,[b]_2)=[ab]_T$. Pairings let us verify \emph{one} multiplicative relation between hidden exponents---the key to SNARKs (Ch.~\ref{ch:pcs}--\ref{ch:snark}). Instantiated by BN254 or BLS12-381.

\subsection{Polynomials over finite fields}
\begin{lemma}[Schwartz--Zippel]\label{lem:sz}
Let $f\in\FF[X_1,\dots,X_n]$ be a nonzero polynomial of total degree $d$ and $S\subseteq\FF$ finite. For $r_1,\dots,r_n\getsr S$, $\Pr[f(r_1,\dots,r_n)=0]\le d/|S|$.
\end{lemma}
\begin{proof}
Induction on $n$. For $n=1$, a nonzero degree-$d$ polynomial has at most $d$ roots. For $n>1$ write $f=\sum_{i=0}^{k}X_n^i f_i(X_1,\dots,X_{n-1})$ with $f_k\ne0$, $\deg f_k\le d-k$. Let $E$ be the event $f_k(r_1..r_{n-1})=0$; $\Pr[E]\le(d-k)/|S|$ by induction. If $\neg E$, then $f(r_1,..,r_{n-1},X_n)$ is a nonzero univariate polynomial of degree $k$ so vanishes at $r_n$ with probability $\le k/|S|$. Hence $\Pr[f(\vec r)=0]\le\Pr[E]+\Pr[\neg E\wedge f=0]\le(d-k)/|S|+k/|S|=d/|S|$.
\end{proof}
\begin{intuition}
Schwartz--Zippel says: \emph{a random point is a perfect polynomial-identity test}. To check two huge polynomials $f=g$, evaluate $f-g$ at a random point. This single lemma underlies sum-check, QAP verification, KZG and FRI.
\end{intuition}
\begin{lemma}[Lagrange interpolation]
For distinct $x_1,\dots,x_n\in\FF$ and any $y_1..y_n$, the unique polynomial of degree $<n$ with $f(x_i)=y_i$ is $f(X)=\sum_i y_i L_i(X)$ with $L_i(X)=\prod_{j\ne i}\frac{X-x_j}{x_i-x_j}$.
\end{lemma}
\begin{example}[Lagrange interpolation over $\FF_{17}$]
Interpolate $(1,1),(2,0),(3,1)$. $L_1(X)=\frac{(X-2)(X-3)}{(1-2)(1-3)}=\frac{(X-2)(X-3)}2$, $L_3(X)=\frac{(X-1)(X-2)}{(3-1)(3-2)}=\frac{(X-1)(X-2)}2$, and $L_2$ is irrelevant since $y_2=0$. So $f=L_1+L_3=\frac{(X-2)\big[(X-3)+(X-1)\big]}{2}=(X-2)^2=X^2-4X+4$. Check: $f(1)=1$, $f(2)=0$, $f(3)=1$. \checkmark
\end{example}
\begin{example}[Schwartz--Zippel in action]
Let $f(X,Y)=XY-X=X(Y-1)$, total degree $2$, over $\FF_{17}$. It vanishes iff $X=0$ or $Y=1$, so for uniform $(x,y)$ the exact probability is $\frac1{17}+\frac1{17}-\frac1{289}=\frac{33}{289}\approx0.114$, which indeed is at most $d/|\FF|=2/17=\frac{34}{289}$. The bound is nearly tight here.
\end{example}
\begin{lemma}[Factor theorem / division]\label{lem:factor}
$f(z)=y$ if and only if $(X-z)$ divides $f(X)-y$ in $\FF[X]$.
\end{lemma}
This last fact powers KZG openings (Chapter~\ref{ch:pcs}).

\section{One-way functions, pseudorandomness and trapdoor permutations}
\begin{definition}
A polynomial-time computable $f:\{0,1\}^*\to\{0,1\}^*$ is a \emph{one-way function} (OWF) if for every PPT $\mathcal A$, $\Pr_{x\getsr\{0,1\}^n}\big[f(\mathcal A(f(x)))=f(x)\big]=\negl(n)$. A \emph{pseudorandom generator} (PRG) is a deterministic polynomial-time $G:\{0,1\}^n\to\{0,1\}^{\ell(n)}$ with $\ell(n)>n$ whose output on a uniform seed is computationally indistinguishable from a uniform $\ell(n)$-bit string. A \emph{trapdoor permutation} (TDP) is a family of permutations $f_{pk}$ that are easy to evaluate and one-way to invert, but easy to invert given a trapdoor $sk$ (e.g.\ RSA: $f_{N,e}(x)=x^e\bmod N$).
\end{definition}
H\aa stad--Impagliazzo--Levin--Luby: OWF $\Leftrightarrow$ PRG. These are the \emph{minimal} assumptions for computational ZK proofs of $\NP$ (Theorem~\ref{thm:gmw} and its converse in Chapter~\ref{ch:theory}).
\begin{example}[Naor's commitment from a PRG]\label{ex:naor}
Let $G:\{0,1\}^n\to\{0,1\}^{3n}$ be a PRG. \emph{Commit to a bit $b$:} the receiver sends a random $R\in\{0,1\}^{3n}$; the committer picks a seed $s$ and sends $c=G(s)$ if $b=0$ and $c=G(s)\oplus R$ if $b=1$. \emph{Hiding:} $G(s)$ and $G(s)\oplus R$ are both pseudorandom, hence indistinguishable. \emph{Binding:} a cheating committer needs $s,s'$ with $G(s)=G(s')\oplus R$, i.e.\ $R\in\{G(s)\oplus G(s'):s,s'\}$, a set of size at most $2^{2n}$; a uniform $R$ lies in it with probability at most $2^{2n}/2^{3n}=2^{-n}$. So binding is \emph{statistical}, hiding is computational, and only OWFs are needed.
\end{example}

\section{Commitment schemes}
\begin{definition}[Commitment scheme]
A commitment scheme consists of PPT algorithms $\Setup(1^\lambda)\to ck$, $\Com_{ck}(m;r)\to c$, with opening $(m,r)$ to $c$. It must be
\begin{itemize}
\item \textbf{Hiding}: for all $m_0,m_1$ the distributions $\Com(m_0;r)$ and $\Com(m_1;r)$ are indistinguishable (perfectly / statistically / computationally);
\item \textbf{Binding}: no PPT adversary outputs $c$ and two openings $(m,r)\neq(m',r')$, $m\ne m'$, with $\Com(m;r)=\Com(m';r')=c$ (except negligible probability; perfect binding if impossible).
\end{itemize}
\end{definition}
\begin{theorem}[No scheme is both perfectly hiding and perfectly binding]
Perfect hiding means every $c$ has openings to every message; so a computationally unbounded committer can equivocate: binding is at best computational. The reverse holds as well.
\end{theorem}
\begin{example}[Pedersen commitment (perfectly hiding, computationally binding)]\label{ex:pedersen}
Let $g,h\in\GG$ generate a group of prime order $q$ with $\log_gh$ \emph{unknown}. Define $\Com(m;r)=g^mh^r$ for $m,r\in\ZZ_q$.
\emph{Hiding:} for each $m$, $g^mh^r$ is uniform in $\GG$ when $r$ is uniform, so $c$ is independent of $m$.
\emph{Binding:} an opening pair $g^mh^r=g^{m'}h^{r'}$ with $m\ne m'$ gives $\log_gh=(m-m')(r'-r)^{-1}\bmod q$, breaking DL.
\emph{Homomorphic:} $\Com(m;r)\cdot\Com(m';r')=\Com(m+m';r+r')$.

\textbf{Numbers.} Toy group, $g=2$, $h=3=2^8$ (so $\log_gh=8$, which in reality must be unknown). Commit $m=4,r=6$: $g^4h^6=16\cdot3^6=16\cdot16=256\equiv3$. Commit $m'=2,r'=3$: $4\cdot3^3=4\cdot4=16$. Product $3\cdot16=48\equiv2$, and $\Com(6;9)=2^6\cdot3^9=18\cdot18=324\equiv2$. \checkmark\\
\emph{Equivocation with the trapdoor $t=8$:} to open $c=3$ (committed to $4$ with $r=6$) as $m^*=7$ we need $4+8\cdot6=7+8r^*\pmod{11}$, i.e.\ $52\equiv8\equiv7+8r^*$, so $8r^*\equiv1$, $r^*=7$ (as $8\cdot7=56\equiv1$). Check: $2^7\cdot3^7=13\cdot2=26\equiv3$. \checkmark
\end{example}
\begin{example}[Hash-based commitment]
$\Com(m;r)=H(m\|r)$ with $H$ modelled as a random oracle (or collision resistant + hiding via randomness extraction). Computationally hiding and binding; the basis of Merkle trees used in FRI/STARKs. A \emph{Merkle tree} commits to a vector $(v_1..v_n)$ with a single root and opens position $i$ with a log-size authentication path.
\begin{center}
\begin{tikzpicture}[every node/.style={font=\small},level distance=1.0cm,
 level 1/.style={sibling distance=4cm},level 2/.style={sibling distance=2cm}]
 \node[draw,rounded corners,fill=yellow!20] {root $=H(h_{12}\|h_{34})$}
  child {node[draw,rounded corners,fill=green!15] {$h_{12}=H(v_1\|v_2)$}
    child {node {$v_1$}} child {node {$v_2$}}}
  child {node[draw,rounded corners] {$h_{34}=H(v_3\|v_4)$}
    child {node[draw,rounded corners,fill=blue!10] {$v_3$}} child {node[draw,rounded corners,fill=green!15] {$v_4$}}};
\end{tikzpicture}
\end{center}
\emph{Example (opening $v_3$).} The prover sends $v_3$ with the authentication path $(v_4,h_{12})$ (green nodes). The verifier computes $h_{34}=H(v_3\|v_4)$ and then $H(h_{12}\|h_{34})$ and compares it with the root. Any other value $v_3'$ accepted for the same root yields a collision of $H$, so the commitment is binding; the path has $\log_2n$ hashes.
\end{example}
\begin{example}[Perfectly binding: Blum / Goldwasser--Micali style]
Given $N=PQ$, commit to a bit $b$ by $c=z^2\cdot y^b\bmod N$ for a fixed non-residue $y\in J_N^{+1}$ and random $z$. Computationally hiding under QRA and perfectly binding (a QR cannot also be a non-residue). Hence for each bit the commitment determines it uniquely.
\end{example}

\section{The random oracle model}
A \emph{random oracle} $H:\{0,1\}^*\to\{0,1\}^\ell$ is a function chosen uniformly at random; all parties (including the adversary) access it as a black box. In proofs the simulator \emph{programs} $H$ (answers queries lazily and consistently, possibly choosing outputs in advance) and \emph{observes} the adversary's queries. The ROM is a heuristic: instantiating $H$ with SHA-3 voids the formal guarantee, but no natural attack is known for natural protocols; contrived counterexamples exist \cite{CGH04}.

\chapter{Interactive Proof Systems}\label{ch:ip}
\section{Definitions}
An \emph{interactive protocol} is a pair $(\Prv,\Ver)$ of interactive probabilistic algorithms that alternately send messages. $\pair{\Prv(w)}{\Ver}(x)$ denotes the (random) output of $\Ver$ (accept $=1$ / reject $=0$) when $\Prv$ has private input $w$ and both have common input $x$. $\Ver$ is PPT; $\Prv$ is computationally unbounded in the \emph{information-theoretic} setting.
\begin{definition}[Interactive proof system]\label{def:ip}
$(\Prv,\Ver)$ is an interactive proof system for $L$ with completeness error $c$ and soundness error $s$ ($c+s<1$ with a gap, typically $c=0$, $s=1/2$) if
\begin{description}
\item[Completeness] $\forall x\in L$: $\Pr[\pair{\Prv}{\Ver}(x)=1]\ge1-c(|x|)$.
\item[Soundness] $\forall x\notin L$, $\forall$ (unbounded) $\Prv^*$: $\Pr[\pair{\Prv^*}{\Ver}(x)=1]\le s(|x|)$.
\end{description}
$\IP$ is the class of languages with an interactive proof system with $c,s$ at most $1/3$ (resp.\ $\negl$ after amplification) and polynomially many rounds. An \emph{argument system} restricts soundness to $\PPT$ provers.
\end{definition}
Why interaction and randomness? Without randomness and without interaction one gets exactly $\NP$ ($\Prv$ sends $w$). Allowing only interaction but a deterministic verifier also gives $\NP$ (the prover can predict every verifier message). Randomness \emph{and} interaction jointly give $\PSPACE$.

\begin{theorem}[Error reduction]\label{thm:errred}
If $(\Prv,\Ver)$ has $c=0$ and soundness $s$, then the $k$-fold \emph{sequential} repetition, accepting iff all $k$ runs accept, has soundness $s^k$. The same bound holds for \emph{parallel} repetition when the protocol is public-coin or has only three messages.
\end{theorem}
\begin{proof}[Proof (sequential case)]
Fix $x\notin L$ and any $\Prv^*$. In each run, conditioned on any history, $\Prv^*$ with the fresh randomness of $\Ver$ is a prover for a single run, so it wins that run with probability $\le s$. Multiply conditional probabilities. (Parallel repetition is subtler: $s^k$ is proved by a different argument for public-coin and three-message protocols, while for private-coin protocols and for computationally sound arguments with more than three rounds parallel repetition may fail to reduce the error at the expected rate \cite{BIN97}.)
\end{proof}

\section{Example: Graph Non-Isomorphism}
\label{sec:gni}
$\GNI=\{(G_0,G_1): G_0\not\cong G_1\}$ is not known to be in $\NP$ (it lies in $\coNP$, but a certificate of non-isomorphism would have to rule out all $n!$ permutations), yet it has a very short interactive proof.
\begin{protocol}[label=pr:gni]{$\GNI\in\IP$ (Goldreich--Micali--Wigderson)}
\textbf{Common input:} graphs $G_0,G_1$ on vertex set $[n]$.
\begin{flow}{Prover (unbounded)}{Verifier}
 & & $b\getsr\{0,1\}$, $\sigma\getsr S_n$\\
 & & $H=\sigma(G_b)$\\
 & $\la{H}$ & \\
find $b'$ with $H\cong G_{b'}$ & & \\
 & $\ra{b'}$ & accept iff $b'=b$
\end{flow}
\end{protocol}
\begin{theorem}
The protocol has completeness $1$ and soundness error $1/2$.
\end{theorem}
\begin{proof}
\emph{Completeness.} If $G_0\not\cong G_1$, then $H$ is isomorphic to exactly one of $G_0,G_1$, namely $G_b$. The unbounded prover finds it and answers $b'=b$.\\
\emph{Soundness.} If $G_0\cong G_1$: the distribution of $H=\sigma(G_b)$ is uniform on the isomorphism class of $G_0$ (for fixed $b$, uniform $\sigma$ yields uniform output, because each graph in the class has the same number $|\Aut(G_0)|$ of permutations mapping $G_b$ to it) and does not depend on $b$. So $H$ carries no information about $b$ and any prover guesses it with probability exactly $1/2$.
\end{proof}
\begin{example}
$G_0$: path $1\!-\!2\!-\!3\!-\!4$, edges $\{12,23,34\}$. $G_1$: star centred at $1$, edges $\{12,13,14\}$. Degree sequences $(1,2,2,1)$ vs.\ $(3,1,1,1)$ so non-isomorphic. Verifier picks $b=1$, $\sigma=(1\mapsto2,2\mapsto4,3\mapsto1,4\mapsto3)$: edges $12\mapsto24$, $13\mapsto21$, $14\mapsto23$, so $H=\{12,23,24\}$, a star centred at $2$. The prover sees a vertex of degree $3$, so $H\cong G_1$ and answers $b'=1$. If the graphs had been isomorphic, a vertex of degree 3 would be there with the same chance whichever $b$ was used.
\end{example}
\begin{remark}
This protocol is \emph{honest-verifier} zero-knowledge (a simulator picks $b$ itself and outputs $(H,b)$), but not zero-knowledge against a cheating $\Ver^*$: $\Ver^*$ could send an $H$ of his own choosing (e.g.\ a graph isomorphic to a third graph $G_2$) and learn from the prover whether $H\cong G_0$ or $G_1$. A fix by the verifier first proving he knows $b,\sigma$ makes it full ZK \cite{GMW91}. This illustrates why ZK must be defined against \emph{arbitrary} (malicious) verifiers.
\end{remark}

\section{Public coins and the landscape of $\IP$}
A protocol is \emph{public-coin} (Arthur--Merlin) if the verifier's messages are just his coin tosses. The $\GNI$ protocol above uses private coins ($\sigma$ is hidden). Goldwasser--Sipser showed that private coins can be converted into public coins at the cost of two extra rounds (set-lower-bound protocol), so $\IP[k]\subseteq\AM[k+2]$. Public-coin protocols are important for us: \emph{only} public-coin protocols can be made non-interactive by Fiat--Shamir.

\begin{theorem}[Shamir 1990 \cite{Sha92}; LFKN 1990]
$\IP=\PSPACE$.
\end{theorem}
The proof arithmetises a $\PSPACE$ computation (quantified Boolean formula, TQBF) into a low-degree polynomial identity and checks it with a variant of the sum-check protocol. We prove sum-check in Chapter~\ref{ch:sumcheck} and discuss the arithmetisation there.

\begin{center}
\begin{tikzpicture}[scale=0.9,every node/.style={font=\small}]
 \draw[rounded corners,fill=gray!8] (0,0) rectangle (11,5);\node at (5.5,4.7) {$\PSPACE=\IP$};
 \draw[rounded corners,fill=blue!6] (0.5,0.4) rectangle (7.5,4.0);\node at (4,3.7) {$\AM$ (public-coin, const rounds)};
 \draw[rounded corners,fill=green!8] (1,0.7) rectangle (4.5,3.2);\node at (2.7,2.9) {$\NP$};
 \node[align=center] at (2.7,1.5) {$\GI,\QR,$\\3-COL, SAT};
 \node[align=center] at (6.1,1.7) {$\GNI$};
 \draw[rounded corners,dashed] (8,0.7) rectangle (10.7,4.0); \node[align=center] at (9.35,2.3) {all of\\\textsf{TQBF}, $\#$SAT,\\permanent};
\end{tikzpicture}
\end{center}

\section*{Exercises}
\begin{exercise}\label{ex:detv}
Show that if $\Ver$ is deterministic then $L\in\IP\Rightarrow L\in\NP$.
\end{exercise}
\begin{exercise}
Prove that in Protocol~\ref{pr:gni} the verifier may use fresh randomness $\sigma$ and $b$ per round, and compute the soundness after $k$ parallel repetitions. Why is private-coin parallel repetition fine here?
\end{exercise}

\chapter{Zero-Knowledge: Definitions}\label{ch:zk}
\section{The simulation paradigm}
The informal slogan ``the verifier learns nothing'' is made precise through a thought experiment of two worlds.
\begin{center}
\begin{tikzpicture}[every node/.style={font=\small},box/.style={draw,rounded corners,minimum width=2.2cm,minimum height=1cm,align=center}]
 \node[box,fill=red!8] (v1) at (0,1.2) {$\Ver^*$\\(cheating)};
 \node[box,fill=blue!8] (p) at (5,1.2) {$\Prv(x,w)$\\honest, has witness};
 \draw[<->,thick] (v1)--(p) node[midway,above] {real interaction};
 \node[box,fill=red!8] (v2) at (0,-1.2) {$\Ver^*$\\(same code)};
 \node[box,fill=green!8] (s) at (5,-1.2) {$\Sim(x)$\\NO witness};
 \draw[<->,thick] (v2)--(s) node[midway,above] {ideal interaction};
 \node[draw,dashed,rounded corners,fit=(v1)(p),inner sep=6pt,label={[left,xshift=-0.6cm]:\textbf{Real world}}] {};
 \node[draw,dashed,rounded corners,fit=(v2)(s),inner sep=6pt,label={[left,xshift=-0.6cm]:\textbf{Ideal world}}] {};
 \node at (8.2,0) [align=center] {the two \emph{views}\\must be\\indistinguishable};
 \draw[<->,dotted,thick] (7.1,1.2)--(7.1,-1.2);
\end{tikzpicture}
\end{center}
If $\Ver^*$ can compute some function $f(x)$ after talking to $\Prv$, the same $\Ver^*$ talking to $\Sim$ (who has no witness) computes $f(x)$ with about the same probability: so the interaction contributed nothing. The simulator may \emph{interact with $\Ver^*$ as a black box with rewinding}, and may run in \emph{expected} polynomial time.

\begin{definition}[View]
For an interactive pair and $\Ver^*$ with auxiliary input $z$, $\view^{\Prv(w)}_{\Ver^*(z)}(x)$ is the random variable $(r,m_1,m_2,\dots,m_t)$ where $r$ is the content of $\Ver^*$'s random tape and $m_i$ are the messages that $\Ver^*$ received from $\Prv$. (The messages $\Ver^*$ sent are determined by $x,z,r$ and $m_1..m_{i-1}$.)
\end{definition}

\begin{definition}[Zero-knowledge, GMR/GMW]\label{def:zk}
Let $(\Prv,\Ver)$ be an interactive proof for the $\NP$-relation $R$. It is
\emph{perfect / statistical / computational zero-knowledge} if for every $\PPT$ $\Ver^*$ there is a (expected) $\PPT$ simulator $\Sim_{\Ver^*}$ such that
\[
\Big\{\view^{\Prv(w)}_{\Ver^*(z)}(x)\Big\}_{(x,w)\in R,\;z\in\{0,1\}^*}\;\;\equiv\;/\approx_s/\approx_c\;\;\Big\{\Sim_{\Ver^*}(x,z)\Big\}_{(x,w)\in R,\;z\in\{0,1\}^*}.
\]
If this only holds for the honest $\Ver$ (with $z=\bot$) it is \emph{honest-verifier} zero-knowledge (HVZK).
\end{definition}
Points to note.
\begin{enumerate}[label=(\alph*)]
\item \textbf{Quantifier order.} The simulator may depend on $\Ver^*$ (it ``knows'' its code) but not on $(x,w,z)$ other than as input.
\item \textbf{Auxiliary input $z$.} It represents whatever the adversary knew before (earlier protocol runs, partial information on $w$). Without it, ZK is not closed under sequential composition \cite{GO94}; with it, it is (Theorem~\ref{thm:seqcomp}).
\item \textbf{Perfect vs computational.} Perfect ZK: distributions identical; no assumption needed; typical for GI, QR. Computational ZK: needs a commitment scheme, typical for 3-COL / all of $\NP$.
\item \textbf{Black-box simulation.} Almost all known simulators only use $\Ver^*$ as an oracle $\Sim^{\Ver^*}(x)$ and rewind it; \emph{non-black-box} simulators (Barak 2001) exist but are exotic.
\item \textbf{Expected polynomial time.} A rewinding simulator may need a random number of attempts. (Strict polynomial-time simulation is possible for fewer protocols.)
\end{enumerate}

\begin{intuition}
Why does the definition use \emph{views} and not only outputs? Because $\Ver^*$ could output its entire view; so view-indistinguishability is the strongest notion and equivalent to output-indistinguishability for all $\Ver^*$.
\end{intuition}

\section{Basic properties}
\begin{proposition}[ZK for $\BPP$ is trivial]
If $L\in\BPP$, the trivial protocol in which $\Prv$ sends nothing and $\Ver$ decides $x\in L$ is perfect ZK with the simulator outputting $(r)$ only. So ZK is only interesting for $L\notin\BPP$ (e.g.\ assuming $\GI\notin\BPP$).
\end{proposition}
\begin{theorem}[Sequential composition]\label{thm:seqcomp}
If $(\Prv,\Ver)$ is (auxiliary-input) computational ZK, then the protocol obtained by running it $m=\poly$ times sequentially (with independent randomness) is also computational ZK.
\end{theorem}
\begin{proof}[Proof sketch]
Let $\Ver^*$ interact with $\Prv$ for $m$ sessions. Define hybrids $H_i$: the first $i$ sessions are simulated and the remaining $m-i$ are real. $H_0$ is the real view, $H_m$ the simulated one. If a distinguisher separates $H_{i-1}$ from $H_i$, construct a $\Ver'$ which plays the first $i-1$ sessions via the simulator, session $i$ with the real prover (the challenge), and the rest honestly with $\Prv$; its auxiliary input encodes the state of $\Ver^*$ after session $i-1$ --- \emph{this is where auxiliary input is needed}. The two hybrids differ only in whether session $i$ is real or simulated, which contradicts single-session ZK for $\Ver'$. Finish by Lemma~\ref{lem:hybrid} (loss factor $m$).
\end{proof}
\begin{warning}[Parallel composition may break ZK]
Running a ZK protocol $k$ times \emph{in parallel} need not preserve ZK: a rewinding simulator must guess all $k$ challenge bits at once and succeeds with probability $2^{-k}$ only. Indeed, Goldreich and Krawczyk \cite{GK96} proved that no $3$-round public-coin protocol with negligible soundness error can be \emph{black-box} ZK unless $L\in\BPP$. Remedies: (i) accept constant soundness error per round and repeat sequentially; (ii) use more rounds; (iii) use HVZK with parallel repetition (this \emph{is} safe: HVZK is preserved because the simulator generates the challenge itself); (iv) switch to non-interactive proofs.
\end{warning}

\section{Witness indistinguishability, proofs of knowledge, arguments}
\begin{definition}[Witness indistinguishability, WI]
$(\Prv,\Ver)$ is WI if for every PPT $\Ver^*$, all $x$, any two witnesses $w_1,w_2\in R(x)$ and every $z$: $\view^{\Prv(w_1)}_{\Ver^*(z)}(x)\approx_c\view^{\Prv(w_2)}_{\Ver^*(z)}(x)$.
\end{definition}
ZK $\Rightarrow$ WI (both views are indistinguishable from $\Sim(x,z)$). WI is weaker but is preserved under parallel and concurrent composition and suffices when the statement has multiple witnesses (OR-proofs, ring signatures, Chapter~\ref{ch:sigma}).

\begin{definition}[Proof of knowledge (Bellare--Goldreich)]\label{def:pok}
$(\Prv,\Ver)$ is a \emph{proof of knowledge} for $R$ with knowledge error $\kappa$ if there is a probabilistic \emph{knowledge extractor} $\Ext$ such that for every prover $\Prv^*$ and every $x$, with $\varepsilon(x)=\Pr[\pair{\Prv^*}{\Ver}(x)=1]$: if $\varepsilon(x)>\kappa(x)$ then $\Ext^{\Prv^*}(x)$ outputs $w\in R(x)$ in expected time at most $\poly(|x|)/(\varepsilon(x)-\kappa(x))$.
\end{definition}
\begin{intuition}
``$\Prv^*$ knows a witness'' means: a witness can be efficiently \emph{computed from $\Prv^*$} (given rewinding access to it) with success tied to how often $\Prv^*$ convinces $\Ver$. A prover who convinces the verifier only $1\%$ of the time ``knows'' it at that level, requiring $\approx100\times$ more work to extract.
\end{intuition}
\begin{example}[Why PoK $\neq$ membership proofs]
The statement ``$y\in\GG$ has a discrete logarithm'' is true for \emph{every} $y\in\GG$, so the language is trivial. But the \emph{knowledge} of $x$ with $g^x=y$ is a meaningful (indeed identification-enabling) claim. Only the PoK definition can express it.
\end{example}
\begin{remark}[Arguments]
If soundness (or knowledge soundness) only holds against PPT provers, we speak of an \emph{argument (of knowledge)}. Pairing-based SNARKs (Ch.~\ref{ch:snark}) are \emph{arguments of knowledge} (AoK). Statistical ZK arguments for $\NP$ exist from collision-resistant hashing; perfect-hiding commitments yield them, whereas statistical ZK \emph{proofs} for all of $\NP$ are unlikely (Theorem~\ref{thm:szkam}).
\end{remark}

\section{Non-interactive variants (preview)}
A \emph{non-interactive zero-knowledge (NIZK)} proof consists of $\Setup\to\mathsf{crs}$, $\mathsf{Prove}(\mathsf{crs},x,w)\to\pi$, $\mathsf{Verify}(\mathsf{crs},x,\pi)\to\{0,1\}$ with completeness, soundness and ZK defined via a simulator $(\Sim_1,\Sim_2)$ that produces the CRS \emph{together with a trapdoor} and then proofs without a witness. In the \emph{random oracle model} the simulator programs the oracle instead. See Chapter~\ref{ch:fs}.

\section*{Exercises}
\begin{exercise}
Show that if a simulator is allowed to output any distribution at all (no indistinguishability) the definition is trivially satisfied. What stops a protocol with soundness error $1$ from being ZK?
\end{exercise}
\begin{exercise}
Prove that ZK implies WI. Give an example of a WI protocol that is not ZK.
\end{exercise}
\begin{exercise}\label{ex:extractor}
Let $\epsilon=\Pr[\pair{\Prv^*}{\Ver}(x)=1]$. For a protocol with challenge space of size $2$ and special soundness (Chapter~\ref{ch:sigma}), design an extractor with expected running time $O(1/\epsilon)$.
\end{exercise}

\chapter{Classical Zero-Knowledge Protocols}
\label{ch:classical}
In this chapter we study the first three zero-knowledge proofs, each exhibiting the same three-move pattern \emph{commit $\to$ challenge $\to$ response}, and we prove every property in full.

\section{Graph isomorphism (Goldreich--Micali--Wigderson)}
\label{sec:gi}
\textbf{Statement:} $G_0,G_1$ are isomorphic graphs on $[n]$. \textbf{Witness:} $\pi\in S_n$ with $G_1=\pi(G_0)$ (meaning: $\{u,v\}\in E_0\iff\{\pi(u),\pi(v)\}\in E_1$).

\begin{protocol}[label=pr:gi]{Zero-knowledge proof of graph isomorphism}
\textbf{Common input} $(G_0,G_1)$. \textbf{Prover's witness} $\pi$ with $\pi(G_0)=G_1$.
\begin{flow}{Prover $(G_0,G_1,\pi)$}{Verifier $(G_0,G_1)$}
$\sigma\getsr S_n$; $H:=\sigma(G_0)$ & $\ra{H}$ & \\
 & $\la{b}$ & $b\getsr\{0,1\}$\\
$\rho:=\sigma$ if $b=0$ & & \\
$\rho:=\sigma\circ\pi\inv$ if $b=1$ & $\ra{\rho}$ & accept iff $H=\rho(G_b)$
\end{flow}
\end{protocol}
\begin{intuition}
$H$ is a freshly shuffled copy of the graph. The verifier asks either ``show me that $H$ is a relabelling of $G_0$'' or ``\ldots of $G_1$''. Either answer alone is a random permutation (it reveals nothing about $\pi$), but if the prover could answer both, then $\rho_1\inv\circ\rho_0$ would be an isomorphism $G_0\to G_1$.
\end{intuition}

\begin{example}\label{ex:gi}
$G_0$: edges $\{12,23,34\}$ (path $1\!-\!2\!-\!3\!-\!4$). The secret $\pi=(1\mapsto3,2\mapsto1,3\mapsto4,4\mapsto2)$ gives $G_1=\{13,14,24\}$ (path $3\!-\!1\!-\!4\!-\!2$). Prover picks $\sigma=(1\mapsto2,2\mapsto3,3\mapsto4,4\mapsto1)$ and sends $H=\sigma(G_0)=\{23,34,41\}$.
\begin{itemize}
\item $b=0$: sends $\rho=\sigma$. Check: $\rho(G_0)=H$. \checkmark
\item $b=1$: $\pi\inv=(1\mapsto2,2\mapsto4,3\mapsto1,4\mapsto3)$ so $\rho=\sigma\circ\pi\inv=(1\mapsto3,2\mapsto1,3\mapsto2,4\mapsto4)$. Then $\rho(G_1)$: $13\mapsto32$, $14\mapsto34$, $24\mapsto14$, i.e.\ $\{23,34,14\}=H$. \checkmark
\end{itemize}
\end{example}
\begin{center}
\begin{tikzpicture}[every node/.style={circle,draw,inner sep=1.5pt,font=\small},scale=0.9]
 \begin{scope}
 \node (a1) at (0,0) {1};\node (a2) at (1,0) {2};\node (a3) at (2,0) {3};\node (a4) at (3,0) {4};
 \draw (a1)--(a2)--(a3)--(a4); \node[draw=none] at (1.5,-0.7) {$G_0$};
 \end{scope}
 \begin{scope}[xshift=5cm]
 \node (b3) at (0,0) {3};\node (b1) at (1,0) {1};\node (b4) at (2,0) {4};\node (b2) at (3,0) {2};
 \draw (b3)--(b1)--(b4)--(b2); \node[draw=none] at (1.5,-0.7) {$G_1=\pi(G_0)$};
 \end{scope}
 \begin{scope}[xshift=2.5cm,yshift=-2.2cm]
 \node (h2) at (0,0) {2};\node (h3) at (1,0) {3};\node (h4) at (2,0) {4};\node (h1) at (3,0) {1};
 \draw (h2)--(h3)--(h4)--(h1); \node[draw=none] at (1.5,-0.7) {$H=\sigma(G_0)$};
 \end{scope}
 \draw[->,thick,blue] (1.5,-0.9)--(2.7,-1.7) node[draw=none,midway,left,font=\scriptsize] {$\sigma$};
 \draw[->,thick,red] (6.5,-0.9)--(5.3,-1.7) node[draw=none,midway,right,font=\scriptsize] {$\rho=\sigma\pi\inv$};
\end{tikzpicture}
\end{center}

\begin{theorem}\label{thm:gi}
Protocol~\ref{pr:gi} is a perfect zero-knowledge proof system for $\GI$ with completeness $1$ and soundness error $1/2$; it is also a proof of knowledge of $\pi$ with knowledge error $1/2$.
\end{theorem}
\begin{proof}
\textbf{Completeness.} $b=0$: $\rho(G_0)=\sigma(G_0)=H$. $b=1$: $\rho(G_1)=\sigma(\pi\inv(\pi(G_0)))=\sigma(G_0)=H$.

\textbf{Soundness.} Suppose $G_0\not\cong G_1$. Whatever $H$ the (unbounded) prover sends, $H$ is isomorphic to at most one of $G_0,G_1$. So at most one challenge value $b$ can be answered; $\Ver$ picks $b$ after seeing $H$ and independently, hence accepts with probability $\le1/2$.

\textbf{Special soundness / extraction.} If the prover answers both challenges for the same $H$ with $\rho_0,\rho_1$, then $\rho_0(G_0)=H=\rho_1(G_1)$ and $\pi'=\rho_1\inv\circ\rho_0$ satisfies $\pi'(G_0)=G_1$: a witness. The extractor runs the prover, then rewinds it to just after $H$ and asks the other challenge; if $\varepsilon$ is the success probability this takes $O(1/\varepsilon)$ expected attempts.

\textbf{Zero-knowledge.} Let $\Ver^*$ be any (efficient, possibly randomized) verifier with randomness $r$ and auxiliary input $z$; The simulator samples $r$ once, uniformly, and keeps it fixed across all rewinds (the output view contains $r$), so we may regard $\Ver^*$ as a deterministic function of the messages it has received. We define a simulator $\Sim$ (using $\Ver^*$ as a black box) that does \emph{not} use $\pi$:
\begin{enumerate}
\item Pick $b'\getsr\{0,1\}$ and $\rho'\getsr S_n$; set $H':=\rho'(G_{b'})$.
\item Run $\Ver^*(H')$ to get its challenge $b$.
\item If $b=b'$: output $(H',\rho')$ together with $\Ver^*$'s tape. Otherwise rewind $\Ver^*$ to the start and go to step 1.
\end{enumerate}
\emph{Claim A.} $H'$ is uniform over the isomorphism class $\mathcal{C}$ of $G_0$ irrespective of $b'$ (since $G_0\cong G_1$, $\rho'(G_{b'})$ is uniform over $\mathcal C$ for either value of $b'$: each $H\in\mathcal C$ has exactly $|\Aut(G_0)|$ preimages $\rho'$). Consequently the challenge $b=\Ver^*(H')$ is independent of $b'$ and $\Pr[b=b']=1/2$ in each attempt. The expected number of attempts is $2$.

\emph{Claim B.} The output has exactly the real distribution. In the real interaction, $(H,\rho)$ is distributed as: $H$ uniform on $\mathcal C$; given $H$, $b=\Ver^*(H)$ is determined; $\rho$ is uniform over the $|\Aut(G_0)|$ isomorphisms $G_b\to H$. In the simulation, by the above, a successful attempt outputs $(H',\rho')$ with
\[
\Pr[(H',\rho')=(H,\rho)\mid \text{success}]
=\frac{\tfrac12\cdot\tfrac{1}{|\mathcal C|\,|\Aut|}\cdot\mathbf 1[\rho(G_{b(H)})=H]}{\tfrac12}=\frac{1}{|\mathcal C||\Aut|},
\]
because for $(H,\rho)$ with $\rho(G_{b(H)})=H$ exactly one value $b'=b(H)$ succeeds, and $\Pr[\text{success}]=\sum_H\frac{1}{|\mathcal C|}\cdot\frac12=\frac12$. This is the real probability. Hence the distributions are \emph{identical}: perfect ZK.
\end{proof}
\begin{example}[Simulating on Example~\ref{ex:gi}]
The simulator never sees $\pi$. It picks $b'=1$, $\rho'=(1\mapsto2,2\mapsto4,3\mapsto1,4\mapsto3)$, so $H'=\rho'(G_1)=\{13,14,24\}^{\rho'}=\{21,23,43\}=\{12,23,34\}$ (which happens to be $G_0$ itself --- a legal element of $\mathcal C$). If $\Ver^*$ answers $b=1$ the simulator outputs $(H',\rho')$; if $b=0$ it rewinds and retries with fresh randomness.
\end{example}
\begin{center}
\begin{tikzpicture}[>=Stealth,every node/.style={font=\small}]
 \draw[thick] (0,0)--(0,-4.4); \node at (0,0.3) {$\Sim$}; \draw[thick] (6,0)--(6,-4.4); \node at (6,0.3) {$\Ver^*$};
 \draw[->] (0,-0.7)--(6,-0.7) node[midway,above] {$H'=\rho'(G_{b'})$};
 \draw[<-] (0,-1.4)--(6,-1.4) node[midway,above] {$b\ne b'$};
 \draw[red,thick,<-,dashed,bend left=25] (6.15,-1.45) to node[right] {rewind} (6.15,-0.6);
 \draw[->] (0,-2.3)--(6,-2.3) node[midway,above] {fresh $H''$};
 \draw[<-] (0,-3.0)--(6,-3.0) node[midway,above] {$b=b''$ \ \checkmark};
 \draw[->] (0,-3.7)--(6,-3.7) node[midway,above] {$\rho''$};
 \draw[decorate,decoration={brace,amplitude=4pt}] (-0.4,-0.5)--(-0.4,-1.6) node[midway,left] {failed};
 \draw[decorate,decoration={brace,amplitude=4pt}] (-0.4,-2.1)--(-0.4,-3.9) node[midway,left] {output};
\end{tikzpicture}
\end{center}

\section{Quadratic residuosity}
\label{sec:qr}
Let $N=PQ$ and $v\in\Zs N$. The prover wants to show $v\in\QR_N$, i.e.\ knows $s$ with $s^2\equiv v$.
\begin{protocol}[label=pr:qr]{ZK proof of quadratic residuosity (GMR 1985)}
\begin{flow}{Prover $(N,v,s)$, $s^2=v$}{Verifier $(N,v)$}
$r\getsr\Zs N$, $x:=r^2$ & $\ra{x}$ & \\
 & $\la{b}$ & $b\getsr\{0,1\}$\\
$y:=r\cdot s^b$ & $\ra{y}$ & accept iff $y^2\equiv x\cdot v^b \pmod N$
\end{flow}
\end{protocol}
\begin{theorem}
Protocol~\ref{pr:qr} is a perfect ZK proof (with soundness $1/2$) of $v\in\QR_N$, and a proof of knowledge of a square root.
\end{theorem}
\begin{proof}
\emph{Completeness:} $y^2=r^2s^{2b}=xv^b$.
\emph{Soundness:} If $v\notin\QR_N$ and both challenges can be answered for $x$, then $x=y_0^2\in\QR_N$ and $xv=y_1^2$, so $v=y_1^2/y_0^2\in\QR_N$ (the quadratic residues form a group): contradiction. So at most one challenge is answerable.
\emph{Extraction:} with answers $y_0,y_1$: $s=y_1/y_0$ since $(y_1/y_0)^2=xv/x=v$.
\emph{ZK:} The simulator guesses $b'\getsr\{0,1\}$, picks $y\getsr\Zs N$ and sets $x:=y^2v^{-b'}$. Because $v\in\QR_N$ and squaring is a 4-to-1 map onto $\QR_N$, $x$ is uniform over $\QR_N$ for both values of $b'$, exactly as the real $x=r^2$. Therefore the challenge $b=\Ver^*(x)$ equals $b'$ with probability exactly $1/2$ and, given a match, $(x,y)$ is distributed as in the real interaction: $x$ uniform in $\QR_N$ and $y$ uniform among the 4 square roots of $xv^b$.
\end{proof}
\begin{example}\label{ex:qr}
$N=77=7\cdot11$, $s=10$, $v=s^2=100\equiv23\pmod{77}$.\\
Prover: $r=3$, $x=9$. 
\begin{itemize}
\item $b=0$: $y=3$; check $y^2=9=x$. \checkmark
\item $b=1$: $y=3\cdot10=30$; $y^2=900\equiv53$ ($900-11\cdot77=53$), and $x\cdot v=9\cdot23=207\equiv53$. \checkmark
\end{itemize}
Simulator (no $s$): guess $b'=1$, pick $y=5$; $v\inv\equiv67\pmod{77}$ (as $23\cdot67=1541=20\cdot77+1$); $x=y^2v\inv=25\cdot67=1675\equiv58$. Then if the verifier asks $b=1$, answer $y=5$: $y^2=25$ and $xv=58\cdot23=1334\equiv25$. \checkmark
\end{example}

\section{Fiat--Shamir and Feige--Fiat--Shamir identification}
In an identification scheme, Alice publishes a public key and then repeatedly proves knowledge of the secret key. Feige--Fiat--Shamir uses $k$ secrets simultaneously to shrink the soundness error to $2^{-k}$ per round.
\begin{protocol}[label=pr:ffs]{Feige--Fiat--Shamir identification}
\textbf{Keys.} Secret $s_1..s_k\in\Zs N$; public $v_i=s_i^2$.
\begin{flow}{Prover}{Verifier}
$r\getsr\Zs N$, $x=r^2$ & $\ra{x}$ & \\
 & $\la{e_1..e_k}$ & $e\getsr\{0,1\}^k$\\
$y=r\prod_i s_i^{e_i}$ & $\ra{y}$ & accept iff $y^2=x\prod_iv_i^{e_i}$
\end{flow}
\end{protocol}
\begin{example}
$N=77$, $s_1=10,s_2=4$, $v_1=23$, $v_2=16$. Prover: $r=3$, $x=9$. Verifier: $e=(1,1)$. $y=3\cdot10\cdot4=120\equiv43$. Check: $43^2=1849=24\cdot77+1\equiv1$ and $x v_1v_2=9\cdot23\cdot16=3312=43\cdot77+1\equiv1$. \checkmark
\end{example}
\begin{remark}[Soundness/ZK trade-off]
\emph{Extraction.} From two accepting answers $y,y'$ to the same $x$ with challenges $e,e'$ we get $(y/y')^2=\prod_jv_j^{e_j-e'_j}$, so $y/y'=\pm\prod_js_j^{e_j-e'_j}$. If $e,e'$ differ in exactly one coordinate $i$, this is $\pm s_i^{\pm1}$, a square root of $v_i^{\pm1}$; a prover that answers all $2^k$ challenges for one $x$ therefore yields all $s_i$ with $k$ such pairs. \emph{HVZK} holds for any $k$: the simulator picks $e,y$ and sets $x=y^2\prod_iv_i^{-e_i}$. Full ZK against malicious verifiers requires the simulator to guess $e$, probability $2^{-k}$, so $k=O(\log\lambda)$ is the limit for one round; for bigger $k$ one passes to HVZK + the techniques of Chapter~\ref{ch:sigma} (e.g.\ committing to $e$ in advance).
\end{remark}

\section*{Exercises}
\begin{exercise}
Write out the knowledge extractor for Theorem~\ref{thm:gi} in detail: given a prover $\Prv^*$ that makes $\Ver$ accept with probability $\varepsilon>1/2$, show that $\Ext$ outputs an isomorphism in expected time $O(1/(\varepsilon-1/2))$.
\end{exercise}
\begin{exercise}
Why must $v$ be required to lie in $\Zs N$ and the prover's $x$ also checked to be in $\Zs N$? (Hint: consider $r$ a multiple of $P$.)
\end{exercise}
\begin{exercise}
Show that Protocol~\ref{pr:qr} is \emph{not} sound for the language $\{v:(\frac vN)=1\}$ and conclude why the Jacobi symbol does not help.
\end{exercise}
\begin{exercise}
Implement the GI simulator for $n=5$ and check empirically that the distribution of transcripts is the same as the real one.
\end{exercise}

\chapter{Zero-Knowledge Proofs for all of NP}
\label{ch:np}
We now show: \emph{every language in $\NP$ has a (computational) zero-knowledge proof} (Goldreich--Micali--Wigderson 1986/1991), assuming one-way functions exist. The strategy: give a ZK proof for an $\NP$-complete problem (graph 3-colouring or Hamiltonicity) and compose it with the Karp reduction.

\section{Commitments as ``digital envelopes''}
Recall (Chapter~\ref{ch:prelim}) that a commitment hides a value yet binds the committer. In the following protocols we need a \emph{non-interactive, perfectly (or statistically) binding, computationally hiding} commitment $\Com$; Naor's construction from a PRG \cite{Nao91} requires only one-way functions (one round of interaction), and the Goldwasser--Micali commitment gives one non-interactively under QRA. Think of it as locking a value in an envelope.

\section{3-colouring}
\begin{protocol}[label=pr:col3]{ZK proof of 3-colourability (one round)}
\textbf{Input} $G=(V,E)$. \textbf{Prover's witness} a proper colouring $\varphi:V\to\{1,2,3\}$.
\begin{flow}{Prover}{Verifier}
$\pi\getsr S_3$ (colour permutation); & & \\
$c_v:=\Com(\pi(\varphi(v));r_v)\ \forall v$ & $\ra{(c_v)_{v\in V}}$ & \\
 & $\la{\{u,w\}}$ & $\{u,w\}\getsr E$\\
open $c_u,c_w$: & $\ra{(a,r_u),(a',r_w)}$ & accept iff openings valid, \\
 & & $a,a'\in\{1,2,3\}$ and $a\ne a'$
\end{flow}
\end{protocol}
\begin{center}
\begin{tikzpicture}[scale=1.1,every node/.style={font=\small}]
 \node[circle,draw,fill=red!30,minimum size=0.8cm] (a) at (0,1.2) {$a$};
 \node[circle,draw,fill=blue!30,minimum size=0.8cm] (b) at (2,2) {$b$};
 \node[circle,draw,fill=green!30,minimum size=0.8cm] (c) at (2,0.2) {$c$};
 \node[circle,draw,fill=red!30,minimum size=0.8cm] (d) at (4,1.2) {$d$};
 \draw (a)--(b)--(d)--(c)--(a) (b)--(c);
 \node at (2,-0.6) {Graph with a proper 3-colouring $\varphi$};
 \begin{scope}[xshift=6.5cm]
 \foreach \n/\x/\y in {a/0/1.2,b/2/2,c/2/0.2,d/4/1.2}{
   \node[draw,rounded corners,fill=gray!20,minimum size=0.8cm] (\n2) at (\x,\y) {\large $\bullet$};}
 \draw (a2)--(b2)--(d2)--(c2)--(a2) (b2)--(c2);
 \draw[red,very thick] (b2)--(c2);
 \node at (2,-0.6) {Commitments; $\Ver$ picks edge $bc$ and sees two different colours};
 \end{scope}
\end{tikzpicture}
\end{center}
\begin{example}\label{ex:3col}
$G$: vertices $a,b,c,d$, edges $ab,ac,bc,bd,cd$. Witness $\varphi=(a,b,c,d)\mapsto(1,2,3,1)$. Prover picks $\pi:1\mapsto3,2\mapsto1,3\mapsto2$, so commits to $(3,1,2,3)$. $\Ver$ asks $bc$; prover opens $(1,2)$; $1\neq2$: accept. Had $\Ver$ asked $ab$ he would have seen $(3,1)$; for $bd$ he would have seen $(1,3)$. In each round the pair of opened colours is a uniformly random ordered pair of \emph{distinct} colours, whatever $\varphi$ is. 

\emph{Cheating prover:} for $K_4$ (6 edges, not 3-colourable) any assignment has $\ge1$ monochromatic edge, so a round detects the cheat w.p.\ $\ge1/6$. After $6k$ sequential rounds the error is $\le(5/6)^{6k}<e^{-k}$.
\end{example}
\begin{theorem}\label{thm:3col}
Repeating Protocol~\ref{pr:col3} $|E|\cdot k$ times sequentially (independent $\pi$, fresh commitments) yields a computational zero-knowledge proof system for 3-COL with completeness $1$ and soundness error $\le e^{-k}$, provided $\Com$ is perfectly binding and computationally hiding. (With a computationally binding commitment one gets an \emph{argument}.)
\end{theorem}
\begin{proof}
\emph{Completeness.} $\pi\circ\varphi$ is a proper colouring so the opened pair always differs.

\emph{Soundness.} If $G$ is not 3-colourable, then perfect binding ensures that the commitments $(c_v)$ determine a unique assignment $\psi:V\to\{1,2,3,\perp\}$ (where $\perp$ means ``not a valid commitment to a colour''). Since $G$ has no proper colouring, some edge $e^*$ has equal or invalid endpoint values. $\Ver$ asks for $e^*$ with probability $\ge1/|E|$, and then the prover cannot open validly. So a single round accepts with probability $\le1-1/|E|$; by Theorem~\ref{thm:errred} the $|E|k$-fold repetition gives $(1-1/|E|)^{|E|k}\le e^{-k}$.

\emph{Zero-knowledge (single round, then sequential composition by Thm.~\ref{thm:seqcomp}).} Given $\Ver^*$, define $\Sim$:
\begin{enumerate}
\item\label{st:simstart} Pick a random edge $\{u,w\}\in E$ and a uniform pair of distinct colours $(a,a')$.
\item Set $c_u=\Com(a)$, $c_w=\Com(a')$ and $c_v=\Com(1)$ for every other $v$ (all with fresh randomness).
\item Send $(c_v)$ to $\Ver^*$ and obtain its edge $e^*$. If $e^*=\{u,w\}$, open $c_u,c_w$ and output $\Ver^*$'s view. Otherwise rewind to step~\ref{st:simstart} (at most $|E|\cdot n$ times, then output $\perp$).
\end{enumerate}
Define distributions over commitment vectors: $\mathcal R$ (the real one, to $\pi\circ\varphi$) and, for each edge $e$, $\mathcal D_e$ (the simulated one when $e$ is guessed). In both, the plaintext on $e$ is a uniform pair of distinct colours; they differ in the plaintexts of the remaining $\le n-2$ vertices. By the hybrid argument (Lemma~\ref{lem:hybrid}), replacing the commitments of vertices one at a time and using hiding,
\[ \mathcal R\approx_c\mathcal D_e\quad\text{for each }e\in E. \]
Let $p_e=\Pr_{\mathcal R}[\Ver^*\text{ asks }e]$, so $\sum_ep_e=1$. The test ``$\Ver^*$ asks $e$'' is an efficient test so $|\Pr_{\mathcal D_e}[\Ver^*\text{ asks }e]-p_e|=\negl$. Hence in one iteration
\[
\Pr[\text{success}]=\frac1{|E|}\sum_e\Pr_{\mathcal D_e}[\Ver^*\text{ asks }e]=\frac1{|E|}\pm\negl,
\]
so the expected number of iterations is $|E|(1\pm\negl)$ and the simulator is expected polynomial time. Conditioned on success, the output is a view with edge $e$ with probability $\frac{(1/|E|)\Pr_{\mathcal D_e}[e]}{\Pr[\text{success}]}\approx p_e$, matching the real world; and conditioned on $e$, the view (commitments to $\mathcal D_e$ vs.\ $\mathcal R$, plus the opened colours, uniformly random distinct in both worlds) are computationally indistinguishable by the same hybrid argument. Summing over the polynomially many $e$ gives $\view_{\Ver^*}\approx_c\Sim(G)$.
\end{proof}
\begin{intuition}
The verifier tests one edge, sees two distinct random colours, and nothing else: all other envelopes remain closed. The permutation $\pi$ guarantees that the two colours are just ``two different random colours'' no matter what. If a coloured graph is invalid, the verifier will, with probability at least $1/|E|$, ask for a violated edge. The need for $|E|$ repetitions is a weakness addressed by Blum's protocol.
\end{intuition}

\section{Blum's protocol for Hamiltonicity}
\textbf{Statement:} $G$ on $[n]$ has a Hamiltonian cycle $C$. Soundness error $1/2$ per round.
\begin{protocol}[label=pr:blum]{Blum's ZK proof for Hamiltonicity}
\begin{flow}{Prover $(G,C)$}{Verifier}
$\pi\getsr S_n$; $A:=$ adjacency matrix of $\pi(G)$ & & \\
commit entrywise: $c_{ij}=\Com(A_{ij})$ & $\ra{(c_{ij})}$ & \\
 & $\la{b}$ & $b\getsr\{0,1\}$\\
$b=0$: reveal $\pi$ and all openings & $\ra{}$ & check $A=\mathrm{adj}(\pi(G))$\\
$b=1$: open only the $n$ entries of $\pi(C)$ & $\ra{}$ & check these are all $1$ and form an $n$-cycle
\end{flow}
\end{protocol}
\emph{Why it works.} For $b=0$ the prover shows the commitments are a faithful relabelled copy of $G$; for $b=1$ she shows a Hamiltonian cycle exists in the committed graph; she cannot do both for a non-Hamiltonian $G$ (otherwise $\pi\inv(\pi(C))$ would be a Hamiltonian cycle of $G$). Revealed information: for $b=0$ a random permutation of $G$; for $b=1$ a random $n$-cycle on random labels. Simulation: guess $b'$; if $b'=0$ commit to $\pi(G)$, if $b'=1$ commit to a random $n$-cycle plus zeros; rewind on mismatch (expected 2 attempts). The soundness error $1/2$ per round, not $1-1/|E|$, is the main advantage over Protocol~\ref{pr:col3}.
\begin{example}
$G=K_4$ minus edge $\{13\}$; Hamiltonian cycle $C:1\!-\!2\!-\!3\!-\!4\!-\!1$. With $\pi=(1\mapsto3,2\mapsto4,3\mapsto1,4\mapsto2)$, $\pi(C)$ is the cycle $3\!-\!4\!-\!1\!-\!2\!-\!3$. For $b=1$ she opens the four commitments $c_{34},c_{41},c_{12},c_{23}$, all containing $1$.
\end{example}

\section{The main theorem}
\begin{theorem}[Goldreich--Micali--Wigderson \cite{GMW91}]\label{thm:gmw}
If one-way functions exist, every language in $\NP$ has a computational zero-knowledge proof system. Moreover the prover can be implemented in probabilistic polynomial time given a witness.
\end{theorem}
\begin{proof}[Proof sketch]
(i) OWF $\Rightarrow$ PRG \cite{HILL99} $\Rightarrow$ perfectly binding, computationally hiding commitment \cite{Nao91}. (ii) Theorem~\ref{thm:3col} gives a ZK proof for 3-COL, $\NP$-complete. (iii) For $L\in\NP$ let $f$ be a Karp reduction to 3-COL and map witnesses by the (efficient) witness-translation of the reduction; run Protocol~\ref{pr:col3} on $f(x)$. Soundness and ZK transfer because $x\in L\iff f(x)\in$ 3-COL.
\end{proof}
\begin{center}
\begin{tikzpicture}[node distance=1.2cm,every node/.style={font=\small}]
 \node[draw,rounded corners,fill=blue!6] (x) {$(x,w)\in R_L$};
 \node[draw,rounded corners,fill=blue!6,right=of x] (g) {$(G_x,\varphi_w)$};
 \node[draw,rounded corners,fill=green!8,right=of g,text width=3.7cm,align=center] (zk) {ZK proof of 3-colourability\\(Protocol~\ref{pr:col3}$\times|E|k$)};
 \draw[->] (x)--(g) node[midway,above] {Karp};
 \draw[->] (g)--(zk);
\end{tikzpicture}
\end{center}
\begin{remark}[Efficiency]
Via Karp reductions the statement size blows up polynomially (and with a big constant): this is irrelevant for feasibility but is terrible in practice---a motivation for the direct arithmetisation techniques of Chapters~\ref{ch:sumcheck}--\ref{ch:snark}, where one gets \emph{succinct} proofs.
\end{remark}

\section{Round complexity and composition limits}
\begin{itemize}
\item The protocols above use $\omega(\log\lambda)$ rounds when repeated sequentially. \textbf{Constant-round} ZK proofs for $\NP$ exist (Goldreich--Kahan 1996; Feige--Shamir) at the cost of non-trivial constructions; by Goldreich--Krawczyk, \emph{black-box} ZK with $3$ rounds is impossible outside $\BPP$ \cite{GK96}.
\item \textbf{Concurrent ZK}: if the verifier interleaves many sessions, rewinding-based simulators blow up. Richardson--Kilian showed that concurrent ZK is possible but needs $\omega(\log\lambda)$ rounds for black-box simulation (Canetti--Kilian--Petrank--Rosen). Barak's \emph{non-black-box} simulation circumvents black-box lower bounds, giving constant-round bounded-concurrent ZK \cite{Bar01}.
\item \textbf{Non-interactive ZK} (Chapter~\ref{ch:fs}) removes the issue entirely, at the price of a CRS or the random oracle.
\end{itemize}

\section*{Exercises}
\begin{exercise}
Show formally that if the commitment in Protocol~\ref{pr:col3} is only computationally binding, the protocol is an argument but not a proof, and explain what a computationally unbounded prover can do.
\end{exercise}
\begin{exercise}
The simulator in the proof of Theorem~\ref{thm:3col} fixes all other colours to $1$. Why would the simulator break if it also revealed those colours? What if the commitment were not hiding?
\end{exercise}
\begin{exercise}
Prove the completeness, soundness and perfect ZK (with a perfectly hiding commitment!) of Protocol~\ref{pr:blum} formally. Which type of binding do you need, and why does a perfectly hiding commitment only give an \emph{argument}?
\end{exercise}
\begin{exercise}\label{ex:rounds3col}
Compute exactly the number of rounds of Protocol~\ref{pr:col3} needed for soundness $2^{-80}$ on a graph with $|E|=1000$.
\end{exercise}

\chapter{$\Sigma$-Protocols}
\label{ch:sigma}
$\Sigma$-protocols are three-move public-coin proofs of knowledge with a very clean algebraic structure. They are the ``assembly language'' of practical zero-knowledge: Schnorr identification and signatures, Pedersen commitment proofs, the ring signatures of CDS, confidential transactions, and the building blocks of Bulletproofs/Halo2 are all $\Sigma$-protocols.

\section{Definition}
\begin{definition}[$\Sigma$-protocol]\label{def:sigma}
A $\Sigma$-protocol for an $\NP$-relation $R$ is a 3-move protocol: $\Prv(x,w)$ sends a \emph{commitment} $a$; $\Ver(x)$ replies with a uniformly random \emph{challenge} $e\getsr\mathcal C$ ($\mathcal C$ finite, often $\ZZ_q$); $\Prv$ sends a \emph{response} $z$; $\Ver$ accepts iff $\phi(x,a,e,z)=1$ for a deterministic predicate $\phi$. It must satisfy:
\begin{description}
\item[Completeness] if $(x,w)\in R$ and both follow the protocol, $\Ver$ accepts.
\item[Special soundness] there is an efficient extractor $\Ext$ which, given $x$ and two accepting transcripts $(a,e,z)$, $(a,e',z')$ with the \emph{same} first message and $e\ne e'$, outputs $w$ with $(x,w)\in R$.
\item[Special honest-verifier ZK (SHVZK)] there is an efficient simulator $\Sim$ which on input $(x,e)$ outputs $(a,z)$ such that $(a,e,z)$ is distributed identically (or indistinguishably) to a real accepting transcript with challenge $e$.
\end{description}
\end{definition}
\begin{center}
\begin{tikzpicture}[>=Stealth,every node/.style={font=\small}]
 \draw[thick] (0,0)--(0,-3.2);\draw[thick] (6,0)--(6,-3.2);
 \node at (0,0.35) {Prover$(x,w)$};\node at (6,0.35) {Verifier$(x)$};
 \draw[->] (0,-0.8)--(6,-0.8) node[midway,above] {$a$ (commitment)};
 \draw[<-] (0,-1.7)--(6,-1.7) node[midway,above] {$e\getsr\mathcal C$ (challenge)};
 \draw[->] (0,-2.6)--(6,-2.6) node[midway,above] {$z$ (response)};
 \node[right] at (6.2,-2.6) {$\phi(x,a,e,z)\stackrel?=1$};
\end{tikzpicture}
\end{center}
\begin{intuition}[Why special soundness gives knowledge]
Two accepting answers to \emph{different} questions about the \emph{same} commitment can only be produced by someone who could ``open the commitment in two ways'', which reveals the witness. The prover cannot change her commitment after seeing the challenge, so a prover that succeeds for many challenges for a given $a$ must know $w$ (rewinding formalises this).
\end{intuition}

\begin{theorem}[Special soundness $\Rightarrow$ proof of knowledge]\label{thm:ss2pok}
Let $(\Prv,\Ver)$ be a $\Sigma$-protocol with $|\mathcal C|=N$. Then it is a proof of knowledge with knowledge error $1/N$. In the simple form proved below: if a prover $\Prv^*$ makes $\Ver$ accept with probability $\varepsilon\ge4/N$, there is an extractor running in expected time $O(1/\varepsilon)$ (calls to $\Prv^*$) that outputs a witness. (Constants are not optimised; the sharp bound $\poly/(\varepsilon-1/N)$ follows from a finer analysis.)
\end{theorem}
\begin{proof}
Fix the internal randomness $\omega$ of $\Prv^*$; then its first message $a=a(\omega)$ is determined. Let $\varepsilon_\omega=\Pr_e[\Ver\text{ accepts}\mid\omega]$ so that $\mathbb E_\omega[\varepsilon_\omega]=\varepsilon$. Call $\omega$ \emph{good} if $\varepsilon_\omega\ge\varepsilon/2$. Bad $\omega$'s contribute less than $\varepsilon/2$ to the acceptance probability, hence
\[\Pr[\omega\text{ good}\mid\text{accept}]\ \ge\ \tfrac12.\]
\emph{Extractor.}
\begin{enumerate}
\item Sample $(\omega,e)$ and run $\Prv^*$ until $\Ver$ accepts; this takes $1/\varepsilon$ attempts in expectation. Let $(a,e,z)$ be the accepting transcript.
\item\label{st:resample} Keeping $\omega$ (so $a$) fixed, sample fresh challenges $e'$ until $\Ver$ accepts with $e'\ne e$, giving up after $8/\varepsilon$ samples.
\item If a second transcript $(a,e',z')$ was found, output $\Ext(x,(a,e,z),(a,e',z'))$ (special soundness).
\end{enumerate}
\emph{Analysis.} If $\omega$ is good, at least $\varepsilon N/2$ challenges are accepted; removing $e$ leaves at least $\varepsilon N/2-1\ge\varepsilon N/4$ (using $\varepsilon N\ge4$). Thus each sample in step~\ref{st:resample} succeeds with probability at least $\varepsilon/4$, and $8/\varepsilon$ samples fail with probability at most $(1-\varepsilon/4)^{8/\varepsilon}\le e^{-2}<0.14$. One pass therefore succeeds with probability at least $\tfrac12\cdot0.86>0.4$ and costs at most $1/\varepsilon+8/\varepsilon$ calls; repeating passes until success takes expected $O(1/\varepsilon)$ calls.
\end{proof}

\begin{theorem}[Parallel repetition and AND]
Running $\Sigma$-protocols for $R_1,R_2$ in parallel with a \emph{single} common challenge (when $\mathcal C_1=\mathcal C_2$) yields a $\Sigma$-protocol for $R_1\wedge R_2=\{((x_1,x_2),(w_1,w_2))\}$. Special soundness follows by applying the two extractors; SHVZK by running the two simulators with the same $e$.
\end{theorem}

\section{Schnorr's protocol}
\label{sec:schnorr}
Let $\GG=\langle g\rangle$ of prime order $q$, $y=g^x$. The prover claims to know $x$.
\begin{protocol}[label=pr:schnorr]{Schnorr identification / proof of knowledge of discrete log}
\begin{flow}{Prover $(g,y;x)$}{Verifier $(g,y)$}
$r\getsr\ZZ_q$, $t:=g^r$ & $\ra{t}$ & \\
 & $\la{c}$ & $c\getsr\ZZ_q$\\
$s:=r+cx\bmod q$ & $\ra{s}$ & accept iff $g^s=t\cdot y^c$
\end{flow}
\end{protocol}
\begin{theorem}\label{thm:schnorr}
Schnorr's protocol is a $\Sigma$-protocol for $R_{\mathrm{DL}}$ with perfect SHVZK; hence a proof of knowledge with knowledge error $1/q$.
\end{theorem}
\begin{proof}
\emph{Completeness:} $g^s=g^{r+cx}=g^r(g^x)^c=ty^c$.\\
\emph{Special soundness:} from $g^s=ty^c$ and $g^{s'}=ty^{c'}$, $c\ne c'$, divide: $g^{s-s'}=y^{c-c'}$, so $y=g^{(s-s')(c-c')\inv}$ and $x:=(s-s')(c-c')\inv\bmod q$ is a witness.\\
\emph{SHVZK:} Given $c$, $\Sim$ picks $s\getsr\ZZ_q$ and sets $t:=g^sy^{-c}$. In a real transcript with challenge $c$, $s=r+cx$ is uniform in $\ZZ_q$ (as $r$ is), and $t$ is then determined by $t=g^sy^{-c}$ (verification equation). In the simulation, $s$ is uniform and $t$ is determined the same way. The two distributions on $(t,s)$ are identical.
\end{proof}
\begin{example}[Toy group: $q=11$, $g=2$ in $\Zs{23}$]\label{ex:schnorr}
Secret $x=7$, public $y=2^7=13$.
\begin{itemize}
\item \textbf{Real run.} $r=8$: $t=2^8=3$. Challenge $c=4$. $s=8+4\cdot7=36\equiv3\pmod{11}$. Verify: $g^s=2^3=8$; $t\,y^c=3\cdot13^4=3\cdot18=54\equiv8$. \checkmark ($13^2=169\equiv8$, $13^4\equiv64\equiv18$.)
\item \textbf{Extraction.} Rewind, challenge $c'=9$: $s'=8+63=71\equiv5$. Then $x=(3-5)(4-9)\inv=(-2)(-5)\inv=2\cdot5\inv\equiv2\cdot9=18\equiv7$ (since $5\cdot9=45\equiv1$). \checkmark
\item \textbf{Simulation} for $c=4$ without $x$: pick $s=6$, $t=2^6\cdot13^{-4}=18\cdot18\inv=1$. Transcript $(t,c,s)=(1,4,6)$; verify: $2^6=18$, $t\,y^4=1\cdot18$. \checkmark
\end{itemize}
Note that the simulated $t=1$ happens to be the identity; in a larger group the simulated commitment looks uniform.
\end{example}
\begin{center}
\begin{tikzpicture}[>=Stealth,scale=1.0,every node/.style={font=\small}]
 % line s = r + c x in the plane (c, s): two accepted transcripts determine x as slope
 \draw[->] (0,0)--(6.2,0) node[right] {$c$}; \draw[->] (0,0)--(0,4) node[above] {$s$};
 \draw[thick,blue] (0,0.9)--(5.6,3.7) node[right] {$s=r+xc$};
 \filldraw (1.4,1.6) circle (2pt) node[above left] {$(c,s)$}; \filldraw (4.2,3.0) circle (2pt) node[below right] {$(c',s')$};
 \draw[dashed] (1.4,1.6)--(4.2,1.6)--(4.2,3.0); \node at (3.0,1.35) {$\Delta c$}; \node[right] at (4.2,2.3) {$\Delta s=x\,\Delta c$};
 \node[align=left] at (3,-0.9) {Two points on the line determine the slope $x$ ($\Rightarrow$ special soundness);\\ one point reveals nothing, as $r$ is a uniform intercept ($\Rightarrow$ ZK).};
\end{tikzpicture}
\end{center}

\subsection{Zero-knowledge against malicious verifiers}
Schnorr is only HVZK as stated. A malicious $\Ver^*$ may choose $c$ as a function of $t$, and the simulator cannot program $t$ \emph{after} seeing $c$.
\begin{itemize}
\item \textbf{Small challenge space.} If $|\mathcal C|=\poly$ the simulator guesses $c'$, builds $(t,s)$ by SHVZK, and rewinds if $\Ver^*(t)\ne c'$; expected $|\mathcal C|$ attempts. So $\Sigma$-protocols with polynomial-size challenge space are \emph{perfect ZK}, but need superlogarithmically many sequential repetitions for negligible soundness. 
\item \textbf{Large challenge space, Damgård's trick.} Add a first move: $\Ver$ commits to his challenge, $\mathsf{cm}=\Com(c;\rho)$, using a perfectly hiding, computationally binding scheme (e.g.\ Pedersen). Then $\Prv$ sends $t$, $\Ver$ opens $(c,\rho)$, and $\Prv$ checks the opening and answers $s$. \emph{Soundness} is unaffected because $\mathsf{cm}$ perfectly hides $c$ from the prover. \emph{Simulation:} run $\Ver^*$ once with a random $t_0$ to learn the challenge $c$ that it opens; rewind $\Ver^*$ to just after it sent $\mathsf{cm}$ and send $t$ produced by the SHVZK simulator for this $c$. By binding, $\Ver^*$ must open $\mathsf{cm}$ to the same $c$ (otherwise it breaks the commitment), so the simulated transcript is accepting. The usual argument shows the expected number of rewinds is $O(1)$.
\item \textbf{Fiat--Shamir} (Ch.~\ref{ch:fs}) removes the malicious verifier by making the challenge a hash.
\end{itemize}

\section{A general view: proofs about group homomorphisms}
\begin{theorem}[Maurer 2009]\label{thm:maurer}
Let $\varphi:(\ZZ_q^m,+)\to(\GG,\cdot)$ be a group homomorphism (with $\GG$ of prime order $q$), $y\in\GG$, and the relation $R_\varphi=\{(y,\vec w):\varphi(\vec w)=y\}$. Then the protocol
\[
t=\varphi(\vec r),\ \vec r\getsr\ZZ_q^m;\qquad c\getsr\ZZ_q;\qquad \vec s=\vec r+c\vec w;\qquad\text{check }\varphi(\vec s)=t\cdot y^c
\]
is a $\Sigma$-protocol (special soundness: $\vec w=(\vec s-\vec s')/(c-c')$; SHVZK: $\vec s\getsr\ZZ_q^m$, $t=\varphi(\vec s)y^{-c}$).
\end{theorem}
\begin{proof}
Verbatim the proof of Theorem~\ref{thm:schnorr} using $\varphi(\vec s-\vec s')=y^{c-c'}$ and the homomorphism property $\varphi(\lambda\vec v)=\varphi(\vec v)^\lambda$.
\end{proof}
Different choices of $\varphi$ give a whole zoo:
\begin{center}\small
\begin{tabular}{@{}lll@{}}\toprule
Name & $\varphi$ & Proves knowledge of \\\midrule
Schnorr & $x\mapsto g^x$ & discrete log $x$ \\
Okamoto/Pedersen opening & $(x_1,x_2)\mapsto g^{x_1}h^{x_2}$ & representation / opening of a Pedersen commitment\\
Chaum--Pedersen & $x\mapsto(g^x,h^x)\in\GG^2$ & $x$ with $\log_gu=\log_hv$ (DDH tuple)\\
Multi-exponentiation & $\vec x\mapsto\prod g_i^{x_i}$ & vector commitment opening\\
Guillou--Quisquater (RSA) & $x\mapsto x^e$ in $\Zs N$ & $e$-th root (needs $c<e$, $\gcd$ arguments)\\\bottomrule
\end{tabular}
\end{center}

\subsection{Guillou--Quisquater: knowledge of an $e$-th root}
Let $N=PQ$, $e$ a prime with $\gcd(e,\varphi(N))=1$, public $Y=s^e\bmod N$ and secret $s$. Here the homomorphism is $\varphi(s)=s^e$ on the \emph{multiplicative} group $\Zs N$ and the challenge space is $\mathcal C=\{0,\dots,e-1\}$.
\begin{protocol}[label=pr:gq]{Guillou--Quisquater}
\begin{flow}{Prover $(N,e,Y;s)$}{Verifier}
$r\getsr\Zs N$, $t:=r^e$ & $\ra{t}$ & \\
 & $\la{c}$ & $c\getsr\{0,..,e-1\}$\\
$z:=r\cdot s^c$ & $\ra{z}$ & accept iff $z^e\equiv t\cdot Y^c\pmod N$
\end{flow}
\end{protocol}
\emph{Completeness:} $z^e=r^e(s^e)^c=tY^c$. \emph{Special soundness:} from $(t,c,z),(t,c',z')$ with $c\ne c'$: $(z/z')^e=Y^{c-c'}$. Since $e$ is prime and $0<|c-c'|<e$, $\gcd(c-c',e)=1$; with Bezout $u(c-c')+ve=1$ we get $s=(z/z')^u\,Y^{v}$, because $\big((z/z')^uY^v\big)^e=Y^{u(c-c')}Y^{ve}=Y$ and $e$-th roots in $\Zs N$ are unique. \emph{SHVZK:} pick $z\getsr\Zs N$, $t:=z^eY^{-c}$.
\begin{example}
$N=77$, $e=7$ (coprime to $\varphi(77)=60$), $s=2$, so $Y=2^7=128\equiv51$. Prover: $r=3$, $t=3^7=2187\equiv31$ ($77\cdot28=2156$). Challenge $c=5$: $z=3\cdot2^5=96\equiv19$. Verify: $z^7=19^7\equiv68$ (as $19^2\equiv53$, $19^4\equiv37$, $19^6\equiv36$, $19^7\equiv68$) and $tY^5=31\cdot51^5\equiv31\cdot32=992\equiv68$ (as $51^2\equiv60$, $51^4\equiv58$, $51^5\equiv32$). \checkmark \emph{Extraction} with a second challenge: only $c-c'$ invertible modulo $e$ matters, which is why $\mathcal C=\{0,..,e-1\}$ with $e$ prime.
\end{example}

\subsection{Okamoto's protocol: knowledge of a representation}
Given $y=g^{x_1}h^{x_2}$ prove knowledge of $(x_1,x_2)$---the opening of a Pedersen commitment.
\begin{protocol}[label=pr:okamoto]{Okamoto / opening of a Pedersen commitment}
\begin{flow}{Prover $(x_1,x_2)$}{Verifier}
$k_1,k_2\getsr\ZZ_q$; $t=g^{k_1}h^{k_2}$ & $\ra{t}$ & \\
 & $\la{c}$ & $c\getsr\ZZ_q$\\
$s_i=k_i+cx_i$ & $\ra{s_1,s_2}$ & accept iff $g^{s_1}h^{s_2}=t\,y^c$
\end{flow}
\end{protocol}
\begin{example}
Toy group, $g=2$, $h=3$, $y=g^4h^6=3$ (the commitment of Example~\ref{ex:pedersen}), witness $(4,6)$. Pick $k=(2,5)$: $t=2^2\cdot3^5=4\cdot13=52\equiv6$. Challenge $c=3$: $s_1=2+12=14\equiv3$, $s_2=5+18=23\equiv1$. Verify: $g^3h^1=8\cdot3=24\equiv1$; $t\,y^3=6\cdot27\equiv6\cdot4=24\equiv1$. \checkmark
\end{example}
Special soundness: two transcripts with $c\ne c'$ give $(x_1,x_2)$. Since Pedersen commitments are \emph{binding}, a prover who knows two different representations of the same $y$ would break DL, hence the \emph{unique} witness.

\subsection{Chaum--Pedersen: equality of discrete logarithms}
\label{sec:cp}
Statement: $(g,h,u,v)$ with $u=g^x$, $v=h^x$ (equivalently $(g,h,u,v)$ is a DH-tuple).
\begin{protocol}[label=pr:cp]{Chaum--Pedersen}
\begin{flow}{Prover $(x)$}{Verifier}
$r\getsr\ZZ_q$; $a=g^r$, $b=h^r$ & $\ra{a,b}$ & \\
 & $\la{c}$ & $c\getsr\ZZ_q$\\
$s=r+cx$ & $\ra{s}$ & accept iff $g^s=a\,u^c$ and $h^s=b\,v^c$
\end{flow}
\end{protocol}
\begin{example}
Toy group; $g=2$, $h=3$, $x=7$: $u=2^7=13$; $v=3^7=2187=95\cdot23+2\equiv2$. Prover $r=5$: $a=2^5=9$, $b=3^5=243=10\cdot23+13\equiv13$. Challenge $c=4$: $s=5+28=33\equiv0$. Check: $g^0=1$ and $a\,u^4=9\cdot18=162=7\cdot23+1\equiv1$ \checkmark; $h^0=1$ and $b\,v^4=13\cdot16=208=9\cdot23+1\equiv1$ \checkmark.
\end{example}
Proof as in Theorem~\ref{thm:maurer} with $\varphi(x)=(g^x,h^x)$. \emph{Applications:} proving correct decryption of ElGamal ($\log_gpk=\log_{c_1}(c_1^{sk})$), correct re-randomisation in mixnets, verifiable random functions (VRFs), threshold cryptography (proof that a partial decryption share is honest).

\section{Compound statements: AND, equality, OR}
\subsection{AND and equality}
To prove $(x_1,x_2)\in R_1\wedge R_2$, run the two protocols with the same challenge. To prove two statements with the \emph{same} witness (e.g.\ $y_1=g_1^x$, $y_2=g_2^x$) use the same response $s$---this is exactly Chaum--Pedersen.

\subsection{OR proofs (Cramer--Damgård--Schoenmakers 1994)}
Prove: ``I know $w_0$ with $(x_0,w_0)\in R_0$ \emph{or} $w_1$ with $(x_1,w_1)\in R_1$'' without revealing which. The trick: the prover simulates the branch whose witness she lacks, and splits the verifier's challenge.
\begin{protocol}[label=pr:or]{Schnorr OR-proof (knows $\log_gy_0$ or $\log_gy_1$)}
Suppose the prover knows $x_b$, $b\in\{0,1\}$, and let $\bar b=1-b$.
\begin{enumerate}
\item \textbf{Commit.} Real branch: $r\getsr\ZZ_q$, $t_b=g^r$. Simulated branch: choose $c_{\bar b},s_{\bar b}\getsr\ZZ_q$, set $t_{\bar b}=g^{s_{\bar b}}y_{\bar b}^{-c_{\bar b}}$. Send $(t_0,t_1)$.
\item \textbf{Challenge.} Verifier sends $c\getsr\ZZ_q$.
\item \textbf{Respond.} $c_b=c-c_{\bar b}$, $s_b=r+c_bx_b$. Send $(c_0,c_1,s_0,s_1)$.
\item \textbf{Verify.} $c_0+c_1\equiv c$, and $g^{s_i}=t_iy_i^{c_i}$ for $i=0,1$.
\end{enumerate}
\end{protocol}
\begin{center}
\begin{tikzpicture}[>=Stealth,every node/.style={font=\small}]
 \node[draw,rounded corners,fill=green!10,text width=3.4cm,align=center] (r) at (0,1.2) {\textbf{real branch} $b$\\ $t_b=g^r$\\ $c_b=c-c_{\bar b}$ (forced)\\ $s_b=r+c_bx_b$};
 \node[draw,rounded corners,fill=red!10,text width=3.4cm,align=center] (s) at (0,-1.4) {\textbf{simulated branch} $\bar b$\\ pick $c_{\bar b},s_{\bar b}$\\ $t_{\bar b}=g^{s_{\bar b}}y_{\bar b}^{-c_{\bar b}}$};
 \node[draw,circle,fill=yellow!20] (c) at (5.5,0) {$c$};
 \draw[->] (r) -- (c) node[midway,above,sloped] {$c_b+c_{\bar b}=c$};
 \draw[->] (s) -- (c);
 \node[align=left] at (9.2,0) {the verifier's $c$ is \emph{free},\\ but the prover has only\\ one degree of freedom:\\ whichever $c_{\bar b}$ she fixed in\\ advance, $c_b$ is determined.};
\end{tikzpicture}
\end{center}
\begin{theorem}
Protocol~\ref{pr:or} is a $\Sigma$-protocol for $R_0\vee R_1$ (with $\mathcal C=\ZZ_q$) and is \emph{witness indistinguishable}: the verifier's view is identically distributed for $b=0$ and $b=1$.
\end{theorem}
\begin{proof}
\emph{Completeness:} for the real branch Schnorr completeness; for the simulated branch $g^{s_{\bar b}}=t_{\bar b}y_{\bar b}^{c_{\bar b}}$ by construction; $c_0+c_1=c$.
\emph{Special soundness:} two accepting transcripts with the same $(t_0,t_1)$ and challenges $c\ne c'$ yield $(c_0,c_1)$, $(c_0',c_1')$ with $c_0+c_1\ne c_0'+c_1'$, hence $c_i\ne c_i'$ for some $i$. For this $i$, the transcripts $(t_i,c_i,s_i)$, $(t_i,c_i',s_i')$ satisfy Schnorr's special soundness and give $x_i$. (The prover ``cannot have cheated in both branches'', since $c_0+c_1$ varies with $c$.)
\emph{WI:} In both cases the distribution of $(t_0,t_1,c_0,c_1,s_0,s_1)$ is: $c_0$ uniform in $\ZZ_q$, $c_1=c-c_0$, $s_i$ uniform, $t_i=g^{s_i}y_i^{-c_i}$. In the real branch $c_b=c-c_{\bar b}$ is uniform as $c_{\bar b}$ is; $s_b=r+c_bx_b$ is uniform since $r$ is. The distribution does not depend on $b$.
\end{proof}
\begin{example}[Numerical OR-proof]
Toy group. Statement: $y_0=13$ (prover knows $x_0=7$) OR $y_1=3$ (prover pretends not to know $\log_23$). So $b=0$.
\begin{itemize}
\item Simulated branch 1: choose $c_1=3$, $s_1=4$; $t_1=g^4y_1^{-3}=16\cdot3^{-3}=16\cdot27\inv=16\cdot4\inv=16\cdot6=96\equiv4$ (since $4\cdot6=24\equiv1$).
\item Real branch 0: $r=5$, $t_0=2^5=9$. Send $(t_0,t_1)=(9,4)$.
\item Verifier's $c=8$. Then $c_0=c-c_1=5$, $s_0=r+c_0x_0=5+35=40\equiv7$.
\item Verify: $c_0+c_1=8$ \checkmark; $g^{s_0}=2^7=13$ and $t_0y_0^{c_0}=9\cdot13^5=9\cdot4=36\equiv13$ \checkmark ($13^5=13^4\cdot13=18\cdot13=234\equiv4$); $g^{s_1}=2^4=16$ and $t_1y_1^{c_1}=4\cdot27\equiv4\cdot4=16$ \checkmark.
\end{itemize}
\end{example}
\begin{example}[Proving a committed bit; building block for range proofs]
Let $C=g^bh^r$ be a Pedersen commitment. To show $b\in\{0,1\}$ prove ``$C=h^r$ or $C/g=h^r$''---an OR of two Schnorr proofs w.r.t.\ base $h$ for statements $C$ and $Cg^{-1}$. Applying it to bits $b_0..b_{n-1}$ and showing $C=\prod C_i^{2^i}$ (using the homomorphism) proves $C$ commits to a value in $[0,2^n)$: a range proof of size $O(n)$ (Bulletproofs improve this to $O(\log n)$, Chapter~\ref{ch:pcs}).
\end{example}
\begin{remark}[Threshold / general access structures]
CDS generalise OR to ``know witnesses for $k$ of $n$ statements'' by sharing the verifier's challenge via Shamir secret sharing: the simulated branches get $n-k$ free shares which determine a degree-$(n-k)$ polynomial $f$ with $f(0)=c$. This yields ring signatures (Rivest--Shamir--Tauman; Abe--Ohkubo--Suzuki) and privacy-preserving credentials.
\end{remark}

\section{From identification to signatures: preview}
A Schnorr transcript $(t,c,s)$ where $c=H(t,m)$ is a signature on $m$. Security and the forking lemma are the subject of Chapter~\ref{ch:fs}.

\section*{Exercises}
\begin{exercise}
Show that Schnorr's protocol is not special-sound if $\mathcal C$ is allowed to contain $c,c'$ with $c-c'$ not invertible mod the group order (e.g.\ over $\ZZ_{pq}$).
\end{exercise}
\begin{exercise}\label{ex:nonce}
Prove that if the prover reuses the nonce $r$ across two proofs with different challenges, the secret key is revealed. (This is how Sony's PS3 ECDSA key was extracted.)
\end{exercise}
\begin{exercise}
Design a $\Sigma$-protocol for ``$y_1=g^{x}$ and $y_2=h^{x^{2}}$'' (hint: Pedersen-commit to $x$, prove a multiplicative relation).
\end{exercise}
\begin{exercise}
Write the 1-out-of-3 OR proof for Schnorr statements and prove WI. How does the challenge splitting generalise?
\end{exercise}
\begin{exercise}
Prove that the Okamoto protocol is perfectly HVZK and that, for a given $y$, a prover who can produce two different representations of $y$ can compute $\log_gh$.
\end{exercise}

\chapter{Fiat--Shamir, Random Oracles and Non-Interactive Zero-Knowledge}
\label{ch:fs}
Interaction is costly: a verifier must be online, and a transcript convinces only the verifier who took part. Can the challenge be generated without a verifier? Fiat and Shamir (1986) suggested: let the prover compute the challenge herself as a hash of everything said so far.

\section{The Fiat--Shamir transform}
Let $(\Prv,\Ver)$ be a public-coin $\Sigma$-protocol with transcript $(a,e,z)$ and $H:\{0,1\}^*\to\mathcal C$ a hash function.
\begin{algorithm}[label=pr:fs]{Fiat--Shamir transform (non-interactive proof of knowledge)}
\textbf{Prove}$(x,w;\,m)$: compute $a\leftarrow\Prv(x,w)$; set $e:=H(x,a,m)$; compute $z\leftarrow\Prv(x,w,a,e)$; output $\pi=(a,z)$ (or $(e,z)$).\\
\textbf{Verify}$(x,\pi=(a,z);\,m)$: $e:=H(x,a,m)$; accept iff $\phi(x,a,e,z)=1$.
\end{algorithm}
The optional message $m$ turns the proof into a \emph{signature of knowledge} / \emph{signature on $m$}.
\begin{center}
\begin{tikzpicture}[>=Stealth,every node/.style={font=\small}]
 \node[draw,rounded corners,fill=blue!6] (p) at (0,0) {Prover};
 \node[draw,rounded corners,fill=yellow!15] (h) at (4,0) {$H(x,a,m)$};
 \node[draw,rounded corners,fill=green!8] (v) at (8.3,0) {Verifier (offline)};
 \draw[->] (p) to[bend left=25] node[above] {$a$} (h);
 \draw[->] (h) to[bend left=25] node[below] {$e$} (p);
 \draw[->,thick] (p) -- ++(0,-1.4) -| node[pos=0.25,below] {$\pi=(a,z)$} (v);
\end{tikzpicture}
\end{center}
\begin{example}[Schnorr signature, toy group]\label{ex:schnorrsig}
Key: $x=7$, $y=13$. Message $m$. Sign: $r=8$, $t=3$; suppose $H(t\|m)=c=4$ (mod 11). Then $s=r+cx=3$; signature $(t,s)=(3,3)$ (or $(c,s)$). Verify: recompute $c=H(3\|m)=4$; check $g^3=8$ equals $t\,y^c=3\cdot18=54\equiv8$. \checkmark In the compact form $(c,s)=(4,3)$ the verifier recomputes $t'=g^sy^{-c}=8\cdot18^{-1}=8\cdot9=72\equiv3$ (as $18\cdot9=162\equiv1$) and checks $H(t'\|m)=c$. Anyone can verify; nobody can forge without $x$ (see below).
\end{example}

\section{Security in the random oracle model}
\begin{theorem}[FS preserves ZK and soundness in the ROM]\label{thm:fs}
Let $\Pi$ be a $\Sigma$-protocol with SHVZK and special soundness and with $\mathcal C$ super-polynomially large. Then FS$(\Pi)$ in the ROM is
\begin{enumerate}
\item \textbf{Non-interactive zero-knowledge}: there is a simulator that, with the power to program $H$, outputs proofs indistinguishable from real;
\item \textbf{Sound / proof of knowledge}: for every PPT adversary $\mathcal A$ making $Q$ random-oracle queries and outputting an accepting $(x,\pi)$ with probability $\varepsilon$, there is an extractor which outputs a witness with probability roughly $\varepsilon^2/Q-\varepsilon/|\mathcal C|$ (forking lemma).
\end{enumerate}
\end{theorem}
\begin{proof}[Proof of ZK part]
Simulator: on a request to prove $x$ (for $x\in L$), sample $e\getsr\mathcal C$, run the SHVZK simulator to get $(a,z)$, and \emph{program} $H(x,a,m):=e$. If $H$ was already queried on that input (probability $\le Q/\text{(entropy of }a)$, negligible if $a$ has superlogarithmic min-entropy) abort. Because $e$ is uniform, $H(x,a,m)=e$ is a perfectly valid lazily-sampled answer, and $(a,e,z)$ has the real distribution by SHVZK.
\end{proof}

\subsection{The forking lemma}
The extractor must obtain two accepting transcripts with the same $a$ but different $e$ from a non-interactive adversary: it \emph{reruns the adversary with the oracle answers changed from the critical query onward}.
\begin{lemma}[General forking lemma, Bellare--Neven]\label{lem:fork}
Let $\mathcal A$ be a randomized algorithm that takes random coins $\rho$ and oracle answers $h_1,\dots,h_Q\in\mathcal C$ ($|\mathcal C|=N$) and outputs either $\perp$ or a pair $(I,\text{out})$ with $I\in[Q]$. Let $\mathrm{acc}=\Pr[\mathcal A\ne\perp]$ over uniform $\rho,\vec h$. The forking algorithm $\mathcal F_{\mathcal A}$: pick $\rho$, $h_1..h_Q$; run $\mathcal A$; if $\perp$ return $\perp$; else let $I$ be its index; pick fresh $h'_I..h'_Q$ and rerun $\mathcal A(\rho;h_1..h_{I-1},h'_I..h'_Q)$; if it returns the same index $I$ and $h_I\ne h'_I$ output success. Then
\[\mathrm{frk}:=\Pr[\mathcal F_{\mathcal A}\text{ succeeds}]\ \ge\ \mathrm{acc}\Big(\frac{\mathrm{acc}}{Q}-\frac1N\Big).\]
\end{lemma}
\begin{proof}
Let $S_i$ be the set of $(\rho,\vec h)$ on which $\mathcal A$ succeeds with index $i$, so $\Pr[S_i]$ sums over $i$ to $\mathrm{acc}$. Fix $i$ and a prefix $(\rho,h_1,\dots,h_{i-1})$ and let
\[\alpha_i(\rho,h_{<i})=\Pr_{h_i,\dots,h_Q}\big[(\rho,\vec h)\in S_i\big].\]
Conditioned on the prefix, the two runs of $\mathcal F_{\mathcal A}$ use independent suffixes, so both land in $S_i$ with probability $\alpha_i^2$. The probability that additionally $h_i=h_i'$ is at most $\alpha_i/N$ (given the first run, the fresh $h'_i$ equals $h_i$ with probability $1/N$). Hence the probability of a successful fork at index $i$ is at least $\alpha_i^2-\alpha_i/N$. Averaging over prefixes and using Jensen's inequality $\mathbb E[\alpha_i^2]\ge\mathbb E[\alpha_i]^2$ with $\beta_i:=\mathbb E[\alpha_i]=\Pr[S_i]$:
\[
\mathrm{frk}\ \ge\ \sum_{i=1}^Q\Big(\beta_i^2-\frac{\beta_i}{N}\Big)\ \ge\ \frac{(\sum_i\beta_i)^2}{Q}-\frac{\sum_i\beta_i}{N}=\frac{\mathrm{acc}^2}{Q}-\frac{\mathrm{acc}}N,
\]
where we used Cauchy--Schwarz, $\sum_i\beta_i^2\ge(\sum_i\beta_i)^2/Q$.
\end{proof}
\begin{intuition}
Run the cheater once. Find the oracle query that determined its commitment $a$. Wind time back to just before that query and give a \emph{different} random answer. If the cheater still produces a valid proof (about $1/Q$ of the time it is ``the same query'' again), you hold two proofs of the same $a$ under different challenges---special soundness extracts $w$. The price is the quadratic loss $\varepsilon^2/Q$: forking proofs are \emph{not tight}.
\end{intuition}
\begin{figure}[h]
\centering
\begin{tikzpicture}[>=Stealth,every node/.style={font=\small}]
 \draw[->,thick] (0,0)--(10,0) node[right] {time};
 \filldraw (3,0) circle (2pt) node[above left] {query $I$: $H(x,a,m)$};
 \draw[->,blue,thick] (3,0) .. controls (5,1) .. (9,1.2) node[right] {$h_I$, accepts $(a,z)$};
 \draw[->,red,thick] (3,0) .. controls (5,-1) .. (9,-1.2) node[right] {$h'_I$, accepts $(a,z')$};
 \draw[dashed] (3,0)--(3,-1.8);
 \node[align=center] at (3,-2.3) {same coins $\rho$ and same answers up to this point};
\end{tikzpicture}
\caption{Forking: two executions diverge at the critical oracle query.}
\end{figure}
\begin{corollary}[EUF-CMA security of Schnorr signatures in the ROM]
If DL is hard in $\GG$, Schnorr signatures are existentially unforgeable under chosen message attacks in the random oracle model (Pointcheval--Stern 1996, with signing-oracle queries answered by the ZK simulator that programs $H$).
\end{corollary}
\begin{proof}[Idea]
Given a DL instance $y$, set the public key to $y$. Answer signing queries via the simulator (no secret needed); answer $H$ queries lazily and consistently. When the forger outputs a forgery $(t,s)$ on a fresh $m^*$, fork on the query $H(t\|m^*)$ to obtain another accepting transcript with the same $t$ and a different $c$; special soundness yields $x=\dlog_gy$.
\end{proof}

\begin{warning}[Pitfalls: ``weak'' Fiat--Shamir and the Frozen Heart family]
\begin{itemize}
\item \textbf{Hash the statement!} The challenge must bind the \emph{statement} $x$ (and all public parameters) as well as the commitment. Omitting $x$ ("weak FS") allows a prover to pick $(a,e,z)$ first and then solve for a statement $x$ consistent with it (e.g.\ Schnorr: pick $s,c$, set $y=(g^s/t)^{1/c}$ for any $t$). This is exactly the class of bugs known as \emph{Frozen Heart} in several zk-SNARK/Bulletproofs implementations (2022).
\item \textbf{Nonce reuse/bias.} The randomness $r$ of Schnorr/ECDSA must be fresh (or derived deterministically: RFC 6979/EdDSA). Reuse leaks $x$.
\item \textbf{Multi-round protocols.} For $(2\mu+1)$-move protocols, \emph{every} challenge must hash the full transcript so far (``strong FS''); security then requires \emph{round-by-round soundness} or $(k_1,\dots,k_\mu)$-special soundness; the knowledge error grows with the number of rounds (Attema--Fehr--Klooß).
\item \textbf{Grinding.} If the proof has soundness error $2^{-\kappa}$ per round and the adversary can retry, repeat/grind; parameters must absorb the adversary's number of hash queries ($\kappa+\log Q$ bits).
\item \textbf{Not secure for all protocols.} There exist (contrived) protocols secure interactively whose FS compilation is insecure for every instantiation of $H$ \cite{GK03}; for natural protocols, FS from correlation-intractable hash functions (Canetti et al.\ 2019, Peikert--Shiehian) is secure from standard assumptions such as LWE.
\end{itemize}
\end{warning}

\section{NIZK in the common reference string model}
The CRS model assumes a trusted setup that samples a string $\mathsf{crs}$ from some distribution known to all.
\begin{definition}[NIZK]\label{def:nizk}
A tuple $(\Setup,\mathsf{Prove},\mathsf{Verify})$ is a NIZK argument/proof for $R$ if:
\begin{description}
\item[Completeness] $\Pr[\mathsf{Verify}(\mathsf{crs},x,\mathsf{Prove}(\mathsf{crs},x,w))=1]=1$ for $(x,w)\in R$.
\item[Soundness] For all (PPT) $\mathcal A$, $\Pr[\mathsf{Verify}(\mathsf{crs},x,\pi)=1\wedge x\notin L]\le\negl$.
\item[Zero-knowledge] There is a PPT $\Sim=(\Sim_1,\Sim_2)$ with $(\mathsf{crs}_{sim},\tau)\leftarrow\Sim_1(1^\lambda)$, $\pi\leftarrow\Sim_2(\mathsf{crs}_{sim},\tau,x)$ so that $(\mathsf{crs},\mathsf{Prove}(\mathsf{crs},x,w))\approx_c(\mathsf{crs}_{sim},\Sim_2(\mathsf{crs}_{sim},\tau,x))$.
\item[Knowledge soundness (optional)] An extractor $\Ext_{\mathsf{crs}}$ with a trapdoor outputs $w$ from any accepting proof.
\end{description}
\end{definition}
\begin{remark}
The simulator \emph{chooses the CRS}: it knows a trapdoor $\tau$ that the real prover does not. Whoever knows $\tau$ can forge proofs for false statements; hence the setup ceremony must be trusted or run as a multi-party computation (Zcash Powers of Tau). NIZK is impossible in the plain model for nontrivial languages (Goldreich--Oren): a deterministic non-interactive verifier lets the simulator output a proof, then a decider for $L$ exists.
\end{remark}

\subsection{The hidden-bits paradigm (Feige--Lapidot--Shamir)}
The classical construction (Blum--Feldman--Micali 1988; Feige--Lapidot--Shamir 1990) proceeds in two steps.

\textbf{Step 1: the hidden-bits model (HBM).} The CRS is a uniformly random string $r\in\{0,1\}^m$ which the \emph{prover} sees but the verifier does not. The prover sends a proof $\pi$ together with a set $I\subseteq[m]$ of positions to \emph{reveal}; the verifier learns only $(r_i)_{i\in I}$ and checks $(x,\pi,I,r_I)$. In this model one can mimic Blum's Hamiltonicity protocol without any commitments: the hidden random bits play the role of commitments that are \emph{automatically} binding (the prover cannot change $r$) and hiding (the verifier cannot see the unrevealed bits). The prover uses blocks of $r$ as adjacency matrices of random graphs, selects a block that happens to be the adjacency matrix of a Hamiltonian cycle on a relabelled copy of $G$, and reveals exactly the positions of the cycle edges. Zero-knowledge holds because the revealed bits are all $1$s, and a simulator can choose $(I,r_I)$ first and fill the unrevealed bits at random. (The technical details of how blocks are generated are in \cite{FLS99}, and Goldreich, \emph{Foundations of Cryptography Vol.\ 1}, \S4.10.)

\textbf{Step 2: from hidden bits to a real CRS.} Use a trapdoor permutation $f$ with hard-core predicate $B$. The CRS is a random string $\sigma=(\sigma_1,\dots,\sigma_m)$ with $\sigma_i$ in the domain of $f$. The prover generates $(f,f^{-1})$ herself and defines the hidden bits $r_i:=B(f^{-1}(\sigma_i))$. To reveal position $i\in I$, she sends the preimage $\rho_i=f^{-1}(\sigma_i)$; the verifier checks $f(\rho_i)=\sigma_i$ and computes $r_i=B(\rho_i)$. The other bits $B(f^{-1}(\sigma_j))$, $j\notin I$, stay computationally hidden by the hard-core property, and they are fixed by $\sigma$, so the prover cannot change them (binding). The verifier must also be convinced that $f$ is a permutation (for RSA this needs an extra argument, e.g.\ a certified key).
\begin{theorem}[FLS90]
If trapdoor permutations exist (e.g.\ RSA), every $L\in\NP$ has a (statistical) NIZK argument with perfect completeness and the prover efficient given $w$.
\end{theorem}
\begin{theorem}[Groth--Ostrovsky--Sahai 2006; Groth--Sahai 2008]
Pairing-based NIZK for satisfiability of pairing-product equations exists under DLIN/SXDH: practical for proving statements such as ``this ciphertext decrypts to a valid signature'', without going through $\NP$-reductions. Proofs are linear in the number of variables.
\end{theorem}
\begin{theorem}[Peikert--Shiehian 2019; Canetti et al.\ 2019]
NIZK for $\NP$ exists from LWE, in the CRS (or even the common random string) model, by instantiating FS with a correlation-intractable hash function.
\end{theorem}

\section{Interactive vs.\ non-interactive: a summary}
\begin{center}\small
\begin{tabular}{@{}p{3.3cm}p{5.5cm}p{5.5cm}@{}}\toprule
 & Interactive ZK & NIZK \\\midrule
Prover/verifier & online both & prover offline; verifier any time\\
Transferability & transcript deniable (simulatable) & proof can be forwarded/verified publicly \\
Setup & none & CRS or random oracle\\
Soundness against & unbounded or PPT provers & PPT (arguments) typically\\
Example & Schnorr ID, GMW 3-COL & Schnorr signature, Groth16, STARK \\\bottomrule
\end{tabular}
\end{center}

\section*{Exercises}
\begin{exercise}
Show that Schnorr signatures with hash $H(t\|m)$ are insecure if $m$ is replaced by a hash that omits $t$ (i.e.\ $c=H(m)$).
\end{exercise}
\begin{exercise}\label{ex:weakfs}
Prove that the weak Fiat--Shamir variant $c=H(t)$ (statement omitted) allows forging proofs for \emph{some} statement of the adversary's choice. Give the attack for Schnorr.
\end{exercise}
\begin{exercise}
Work through the forking lemma with $Q=1$, and conclude that in this special case $\mathrm{frk}\ge\mathrm{acc}^2-\mathrm{acc}/N$. Compare with the extraction in Theorem~\ref{thm:ss2pok}.
\end{exercise}
\begin{exercise}\label{ex:forkloss}
Using your answer to the previous exercise, estimate the security loss: if the best attacker against the DL problem in the group runs in time $2^{128}$ and $Q=2^{60}$, what bit security can the forking-lemma reduction guarantee for Schnorr signatures?
\end{exercise}

\chapter{The Sum-Check Protocol and Interactive Proofs for Computation}
\label{ch:sumcheck}
From here on our goal changes: we want proofs that are \emph{succinct}---much shorter and faster to verify than the computation they certify. The first tool is the sum-check protocol of Lund, Fortnow, Karloff and Nisan (1990), the engine behind $\IP=\PSPACE$, GKR, and many modern SNARKs (Spartan, HyperPlonk, Lasso/Jolt).

\section{Multivariate polynomials and multilinear extensions}
\begin{definition}[Multilinear extension]
For $f:\{0,1\}^v\to\FF$ the \emph{multilinear extension} (MLE) is the unique polynomial $\tilde f\in\FF[X_1..X_v]$ of degree $\le1$ in each variable agreeing with $f$ on $\{0,1\}^v$:
\[\tilde f(\vec X)=\sum_{\vec b\in\{0,1\}^v}f(\vec b)\cdot\chi_{\vec b}(\vec X),\qquad\chi_{\vec b}(\vec X)=\prod_{i=1}^v\big(b_iX_i+(1-b_i)(1-X_i)\big).\]
\end{definition}
\begin{example}
$f:\{0,1\}^2\to\FF$ with $f(00)=6$, $f(01)=9$, $f(10)=1$, $f(11)=4$. Then $\tilde f(X_1,X_2)=6(1-X_1)(1-X_2)+9(1-X_1)X_2+X_1(1-X_2)+4X_1X_2$. At $(X_1,X_2)=(2,3)$: $6(-1)(-2)+9(-1)(3)+2(-2)+4\cdot6=12-27-4+24=5$.
\end{example}
Uniqueness: a multilinear polynomial vanishing on $\{0,1\}^v$ is zero (induction on $v$). An MLE of a vector of length $2^v$ is evaluable at any point in $O(2^v)$ operations. Key advantage over univariate encodings: \emph{no FFT is required}, and evaluations on the Boolean hypercube correspond to the original data.

\section{The protocol}
\textbf{Claim:} $H=\sum_{\vec b\in\{0,1\}^v}g(\vec b)$ where $g\in\FF[X_1..X_v]$ has degree $\le d$ in each variable. The verifier can evaluate $g$ at a single random point (via an oracle, or through a commitment, or via recursion).
\begin{protocol}[label=pr:sumcheck]{Sum-check}
\textbf{Round 1.} $\Prv$ sends univariate $g_1(X)=\sum_{\vec b\in\{0,1\}^{v-1}}g(X,\vec b)$. $\Ver$ checks $g_1(0)+g_1(1)=H$, picks $r_1\getsr\FF$, sends it.\\
\textbf{Round $i$ ($2\le i\le v$).} $\Prv$ sends $g_i(X)=\sum_{\vec b\in\{0,1\}^{v-i}}g(r_1,\dots,r_{i-1},X,\vec b)$. $\Ver$ checks $g_i(0)+g_i(1)=g_{i-1}(r_{i-1})$, picks $r_i$.\\
\textbf{Final check.} $\Ver$ checks $g_v(r_v)=g(r_1,\dots,r_v)$ (one evaluation of $g$).\\
(Each $g_i$ has degree $\le d$, so is sent as $d+1$ field elements.)
\end{protocol}
\begin{center}
\begin{tikzpicture}[>=Stealth,every node/.style={font=\small}]
 \foreach \i/\w in {0/3.6,1/2.6,2/1.6,3/0.6}{
   \draw[fill=blue!\the\numexpr 10+\i*10\relax!white] (0,-\i*1.1) rectangle (\w*1.6,-\i*1.1-0.7);}
 \node[right] at (6.0,-0.35) {$H=\sum_{b_1b_2b_3b_4}g(b)$ \quad (claim)};
 \node[right] at (4.3,-1.45) {$g_1(X_1)$: $2^{3}$ terms, fix $r_1$};
 \node[right] at (3.0,-2.55) {$g_2(X_2)$: $2^{2}$ terms, fix $r_2$};
 \node[right] at (1.8,-3.65) {$g_3,g_4$: remaining; final check $g(r_1..r_4)$};
 \node[left] at (-0.2,-0.35) {$2^4$};
\end{tikzpicture}
\end{center}
\begin{theorem}[Completeness and soundness]\label{thm:sumcheck}
The protocol has perfect completeness. If $H\ne\sum_{\vec b}g(\vec b)$, then for every prover $\Prv^*$, $\Ver$ accepts with probability at most $vd/|\FF|$.
\end{theorem}
\begin{proof}
\emph{Completeness} is immediate from $g_i(0)+g_i(1)=\sum_{b_i}\sum_{\vec b'}g(r_{<i},b_i,\vec b')=g_{i-1}(r_{i-1})$.

\emph{Soundness} by induction on $v$. For $v=1$: $\Prv^*$ sends some polynomial $g_1^*$ of degree $\le d$ with $g_1^*(0)+g_1^*(1)=H$ (otherwise reject). Since $H\ne g(0)+g(1)$, $g_1^*\ne g$ as polynomials, so they agree on at most $d$ points: $\Pr_{r_1}[g_1^*(r_1)=g(r_1)]\le d/|\FF|$, and the final check passes only then. For $v>1$: let $g_1^*$ be the first message. If $g_1^*=g_1$ (the true polynomial) then $g_1(0)+g_1(1)=\sum_{\vec b}g(\vec b)\ne H$ and the verifier rejects at once. So $g^*_1\ne g_1$, and with probability $\ge1-d/|\FF|$ over $r_1$, $g_1^*(r_1)\ne g_1(r_1)=\sum_{\vec b'}g(r_1,\vec b')$. Then the remaining $(v-1)$-round protocol is on a \emph{false} claim about $g(r_1,\cdot)$, which by the inductive hypothesis is accepted with probability $\le(v-1)d/|\FF|$. A union bound gives $vd/|\FF|$.
\end{proof}
\begin{example}\label{ex:sumcheck}
Let $g(x_1,x_2,x_3)=2x_1^3+x_1x_3+x_2x_3$ over a large field (here integers suffice for the arithmetic). The true sum over $\{0,1\}^3$ is $H=12$: on Boolean points $x_1^3=x_1$ so the $2x_1$ term contributes $2\cdot4=8$, $x_1x_3$ contributes $2$ and $x_2x_3$ contributes $2$.
\begin{enumerate}
\item $\Prv$: $g_1(X)=\sum_{x_2,x_3}g(X,x_2,x_3)=8X^3+2X+1$. $\Ver$: $g_1(0)+g_1(1)=1+11=12=H$ \checkmark. $r_1=2$; $g_1(2)=64+4+1=69$.
\item $\Prv$: $g_2(X)=\sum_{x_3}g(2,X,x_3)=g(2,X,0)+g(2,X,1)=16+(16+2+X)=34+X$. $\Ver$: $g_2(0)+g_2(1)=34+35=69=g_1(2)$ \checkmark. $r_2=3$; $g_2(3)=37$.
\item $\Prv$: $g_3(X)=g(2,3,X)=16+2X+3X=16+5X$. $\Ver$: $g_3(0)+g_3(1)=16+21=37$ \checkmark. $r_3=5$; $g_3(5)=41$.
\item Final: $\Ver$ computes $g(2,3,5)=16+10+15=41$. \checkmark
\end{enumerate}
A cheating prover claiming $H'=13$ must send a $g_1^*$ with $g^*_1(0)+g^*_1(1)=13$, hence $g_1^*\ne g_1$; it is caught unless $r_1$ is one of the $\le3$ roots of $g_1^*-g_1$.
\end{example}
\begin{theorem}[Efficiency]
The verifier does $O(vd)$ field operations plus one evaluation of $g$ and communicates $O(vd)$ elements. If $g$ is given as a table of $2^v$ values (e.g.\ products of MLEs of tables), the prover can be implemented in time $O(2^v\cdot\poly(d))$ by halving the table in each round (Thaler/Libra/Virgo).
\end{theorem}

\section{Application 1: $\#$SAT, so $\mathsf{coNP}\subseteq\IP$}\label{sec:sharpsat}
Arithmetise a 3-CNF $\varphi$ with polynomial $p_\varphi$: variable $x\mapsto x$; $\neg x\mapsto1-x$; clause $\ell_1\vee\ell_2\vee\ell_3\mapsto1-\prod_j(1-\ell_j)$; $\varphi=\bigwedge_k C_k\mapsto\prod_kp_{C_k}$. On Boolean inputs $p_\varphi(\vec b)=\varphi(\vec b)\in\{0,1\}$, hence $\#\mathrm{SAT}(\varphi)=\sum_{\vec b}p_\varphi(\vec b)$, with $\deg\le m$ (the number of clauses) in total and per variable $\le$ clause occurrences.
\begin{example}
$\varphi=(x_1\vee x_2)\wedge(\neg x_1\vee x_2)$. $p=\big(x_1+x_2-x_1x_2\big)\big(1-x_1+x_1x_2\big)$, per-variable degree $2$. Boolean evaluations: $p(00)=0$, $p(01)=1$, $p(10)=0$, $p(11)=1$, so $\#\mathrm{SAT}=2$ ($x_2=1$).\\
Sum-check over $\FF_{17}$: $g_1(X)=p(X,0)+p(X,1)=X(1-X)+1=1+X-X^2$; $g_1(0)+g_1(1)=1+1=2$ \checkmark. $r_1=5$: $g_1(5)=1+5-25=-19\equiv15$. $g_2(X)=p(5,X)=(5-4X)(-4+5X)=-20+41X-20X^2$; $g_2(0)+g_2(1)=-20+1=-19\equiv15$ \checkmark. $r_2=3$: $g_2(3)=-20+123-180=-77\equiv8$. Final: $p(5,3)=(5+3-15)(1-5+15)=(-7)(11)=-77\equiv8$ \checkmark.
\end{example}
Since $\coNP$-complete \textsf{UNSAT} $=\{\varphi:\#\mathrm{SAT}(\varphi)=0\}$, we get $\coNP\subseteq\IP$---a result that astonished the community in 1989--90 because $\coNP$ proofs were thought to need exponential length. For $\IP=\PSPACE$ one must handle the degree blow-up in the quantifier arithmetisation (LFKN's ``degree-reduction'' linearisation operators).

\section{Application 2: the proof of $\IP=\PSPACE$}
\label{sec:ippspace}
Since \textsf{TQBF} (true quantified Boolean formulas) is $\PSPACE$-complete, it suffices to give an interactive proof for $\Phi=\exists x_1\forall x_2\exists x_3\cdots\,\psi(x_1,\dots,x_n)$ with $\psi$ a 3-CNF.

\textbf{Arithmetisation.} Replace $\psi$ by its polynomial $p_\psi$ (Section~\ref{sec:sharpsat}) and quantifiers by operators acting on polynomials:
\[\exists x\,p\ \mapsto\ \textstyle\sum_{x\in\{0,1\}}p=p|_{x=0}+p|_{x=1},\qquad \forall x\,p\ \mapsto\ \prod_{x\in\{0,1\}}p=p|_{x=0}\cdot p|_{x=1}.\]
On Boolean inputs $p_\psi\in\{0,1\}$, so $\Phi$ is true iff the resulting integer $V=\sum_{x_1}\prod_{x_2}\sum_{x_3}\cdots p_\psi(\vec x)$ is nonzero.
\begin{example}
$\Phi=\exists x\,\forall y\,(x\vee y)$ with $p_\psi=x+y-xy$. The $\forall y$ step gives $p_\psi(x,0)p_\psi(x,1)=x\cdot1=x$; the $\exists x$ step gives $0+1=1\ne0$. So $\Phi$ is true ($x=1$).
\end{example}
\textbf{The obstacle.} Each $\prod$ \emph{squares the degree} in the remaining variables, so after $n$ quantifiers the degree is exponential and the prover's univariate polynomials cannot be sent. \textbf{The fix (Shamir):} insert \emph{linearisation operators}
\[L_xp=(1-x)\,p|_{x=0}+x\,p|_{x=1},\]
which return the multilinear polynomial that agrees with $p$ on $\{0,1\}$ in $x$ and thus does not change any Boolean value. Applying $L_{x_j}$ for all $j$ after each quantifier keeps every variable at degree $\le 1$ in the intermediate polynomials (e.g.\ $p=x^2\Rightarrow L_xp=x$, which agrees with $x^2$ on $\{0,1\}$).
\begin{protocol}[label=pr:shamir]{Shamir's protocol (outline)}
The claim is $v=O_1O_2\cdots O_N\,p_\psi$ where each $O_k\in\{\sum_{x},\prod_{x},L_{x}\}$ ($N=O(n^2)$ operators). Round $k$ peels off the leftmost remaining operator $O_k$ acting on variable $x$:
\begin{enumerate}
\item $\Prv$ sends the univariate polynomial $g(X)$ obtained by evaluating the remaining expression with $x:=X$ (and previous variables fixed to the verifier's random points $r_1,\dots$).
\item $\Ver$ checks consistency with the current claim $v$: $g(0)+g(1)=v$ for $\sum$; $g(0)g(1)=v$ for $\prod$; $(1-r)g(0)+r\,g(1)=v$ for $L_x$ where $r$ is the current value of $x$.
\item $\Ver$ picks a fresh random $r'\in\FF$, sets the new claim $v':=g(r')$ for the remaining operators with $x:=r'$ (for $L_x$ the verifier first sends $r'$ and the prover answers with $g$ as above).
\end{enumerate}
At the end the verifier evaluates $p_\psi(r_1,..,r_n)$ itself and compares.
\end{protocol}
\begin{theorem}[LFKN/Shamir]
Over a field $|\FF|\ge2^n$ the protocol is complete and has soundness error $O(n^2m/|\FF|)$, with $m$ the number of clauses, since each of the $O(n^2)$ rounds is a sum-check step with polynomial degree $\le\max(2,\deg_x p_\psi)\le m$. Hence $\PSPACE\subseteq\IP$; with $\IP\subseteq\PSPACE$ (the honest prover's optimal strategy can be computed in polynomial space) we obtain $\IP=\PSPACE$.
\end{theorem}
\begin{proof}[Proof of soundness]
Exactly as in Theorem~\ref{thm:sumcheck}: if the current claim $v$ is false, the polynomial $g^*$ sent by a cheating prover differs from the true $g$ (otherwise the consistency check fails at once); two distinct polynomials of degree $\le m$ agree on at most $m$ points, so with probability $\ge1-m/|\FF|$ the new claim $g^*(r')$ is again false. A union bound over the $N=O(n^2)$ rounds, and the fact that a false initial claim implies a false last claim, gives the bound.
\end{proof}

\section{Application 3: GKR---delegating computation}
\label{sec:gkr}
Goldwasser, Kalai and Rothblum (2008) verify a layered arithmetic circuit $C$ of size $S$, depth $D$, fan-in 2 with a prover of time $\tilde O(S)$ and verifier of time $\tilde O(n+D\log S)$---far less than $S$ when $C$ is low-depth and ``regular''.
\begin{center}
\begin{tikzpicture}[>=Stealth,every node/.style={font=\small},gate/.style={circle,draw,minimum size=0.7cm}]
 \node at (-1.7,2.2) {layer 0 (output)};\node at (-1.7,1.1) {layer 1};\node at (-1.7,0) {layer 2 (input)};
 \node[gate] (o) at (2,2.2) {$\times$};
 \node[gate] (g1) at (1,1.1) {$\times$};\node[gate] (g2) at (3,1.1) {$+$};
 \foreach \n/\x/\v in {i1/0.4/2,i2/1.6/3,i3/2.4/4,i4/3.6/5}{\node[gate] (\n) at (\x,0) {\v};}
 \draw[->] (g1)--(o);\draw[->] (g2)--(o);
 \draw[->] (i1)--(g1);\draw[->] (i2)--(g1);\draw[->] (i3)--(g2);\draw[->] (i4)--(g2);
 \node[right] at (4.4,2.2) {$V_0=(54)$};\node[right] at (4.4,1.1) {$V_1=(6,9)$};\node[right] at (4.4,0) {$V_2=(2,3,4,5)$};
\end{tikzpicture}
\end{center}
\textbf{Layer function.} Let $V_i:\{0,1\}^{k_i}\to\FF$ give the gate values at layer $i$ and $\mathrm{add}_i,\mathrm{mult}_i:\{0,1\}^{k_i+2k_{i+1}}\to\{0,1\}$ be wiring predicates (1 iff gate $z$ at layer $i$ is an add/mult of gates $x,y$ at layer $i+1$). Then for all $z$,
\[
V_i(z)=\sum_{x,y\in\{0,1\}^{k_{i+1}}}\Big[\mathrm{add}_i(z,x,y)\big(V_{i+1}(x)+V_{i+1}(y)\big)+\mathrm{mult}_i(z,x,y)V_{i+1}(x)V_{i+1}(y)\Big].
\]
Replacing every function by its MLE extends this identity to all $z\in\FF^{k_i}$: $\tilde V_i(z)=\sum_{x,y}f_{z}(x,y)$ where $f_z$ is a low-degree polynomial in $(x,y)$.
\begin{algorithm}[label=pr:gkr]{GKR (outline)}
\begin{enumerate}
\item $\Prv$ sends claimed output $\tilde V_0$; $\Ver$ picks random $z_0$ and obtains the claim $\tilde V_0(z_0)=m_0$.
\item For $i=0,\dots,D-1$: run sum-check on $\tilde V_i(z_i)=\sum_{x,y}f_{z_i}(x,y)$. At the end the verifier needs $\tilde V_{i+1}(x^*)$ and $\tilde V_{i+1}(y^*)$ at two random points. The prover claims two values; to avoid doubling the claims at each layer, $\Ver$ sends a line $\ell(t)$ through $x^*,y^*$, $\Prv$ sends $\tilde V_{i+1}\circ\ell$ (degree $\le k_{i+1}$), $\Ver$ checks consistency at $t=0,1$, picks $t^*$ and defines $z_{i+1}=\ell(t^*)$ with claim $m_{i+1}=\tilde V_{i+1}\circ\ell(t^*)$.
\item At the input layer, $\Ver$ computes $\tilde V_D(z_D)$ directly from the inputs (an MLE evaluation, $O(n)$) and compares with $m_D$.
\end{enumerate}
\end{algorithm}
\begin{example}[One layer of GKR on the circuit above]
Output $V_0=54$. The single output gate multiplies gates $0$ and $1$ of layer 1, so $\mathrm{mult}_0(0,x,y)=(1-x)\,y$ for $x,y\in\{0,1\}$ and $\tilde V_1(x)=6(1-x)+9x=6+3x$. The identity says $54=\sum_{x,y\in\{0,1\}}(1-x)y\,(6+3x)(6+3y)$ (only $(x,y)=(0,1)$ survives: $6\cdot9=54$ \checkmark). Run sum-check on $f(x,y)=(1-x)y(6+3x)(6+3y)$:
\begin{itemize}
\item $g_1(X)=f(X,0)+f(X,1)=9(1-X)(6+3X)=54-27X-27X^2$, and $g_1(0)+g_1(1)=54+0=54$ \checkmark. $r_1=2$: $g_1(2)=-108$.
\item $g_2(Y)=f(2,Y)=(1-2)\,Y\,(12)(6+3Y)=-72Y-36Y^2$; $g_2(0)+g_2(1)=-108$ \checkmark. $r_2=3$: $g_2(3)=-216-324=-540$.
\item Final: $\Ver$ needs $\mathrm{mult}_0(0,2,3)=(1-2)\cdot3=-3$ (computable by itself), $\tilde V_1(2)=12$, $\tilde V_1(3)=15$ (claimed by the prover): $-3\cdot12\cdot15=-540$ \checkmark. The claims about $\tilde V_1$ at $2$ and $3$ are reduced (via the line trick) to a single claim about $\tilde V_1$ at one random point, which feeds layer 1's sum-check, and so on until the inputs.
\end{itemize}
\end{example}
\begin{theorem}[GKR 2008]
For every log-space uniform fan-in-two layered circuit of depth $D$, size $S$, with $n$ inputs, there is a public-coin interactive proof with soundness error $O(D\log S/|\FF|)$, $O(D\log S)$ rounds, communication $O(D\log S)$ field elements, verifier time $\tilde O(n+D)$ and prover time $\poly(S)$ (quasi-linear $O(S\log S)$ for regular circuits \cite{Tha13}).
\end{theorem}
Via Fiat--Shamir, GKR gives non-interactive proofs; combined with a polynomial commitment for the inputs and zero-knowledge masking, one obtains a zk-SNARK (Libra, Virgo, Spartan).

\section*{Exercises}
\begin{exercise}\label{ex:sumcheckrun}
Run the sum-check protocol for $g(x_1,x_2)=x_1x_2+3x_2$ over $\FF_{11}$ with $H=\sum g$, using challenges $r_1=4,r_2=7$. Then simulate a cheating prover that claims $H'=H+1$ and find which $r_1$ let it succeed.
\end{exercise}
\begin{exercise}
Prove that the MLE of the equality function $\mathrm{eq}(x,y)=\prod_i(x_iy_i+(1-x_i)(1-y_i))$ is itself; deduce $\tilde f(r)=\sum_{b}f(b)\,\mathrm{eq}(b,r)$.
\end{exercise}
\begin{exercise}
Show that the sum-check prover can be implemented in $O(2^v)$ time for $g=\tilde f\cdot\tilde h$ when $f,h$ are given as tables (hint: "bookkeeping tables" which are folded by $T[b]\leftarrow(1-r)T[0b]+rT[1b]$ after each round).
\end{exercise}
\begin{exercise}
Verify that the exact degree bound matters: why would $\FF=\FF_2$ break the soundness argument?
\end{exercise}

\chapter{Arithmetisation: From Programs to Polynomial Identities}
\label{ch:arith}
To prove statements about computation we convert them into algebraic constraints over a large finite field $\FF$ (typically $|\FF|\approx2^{254}$ or $2^{64}$ with extension fields). This process is called \emph{arithmetisation}. Its aim is to turn ``this program ran correctly'' into ``these polynomials satisfy certain identities'', which Schwartz--Zippel allows us to check at a single random point.
\begin{center}
\begin{tikzpicture}[node distance=0.55cm,every node/.style={font=\small},b/.style={draw,rounded corners,fill=blue!6,minimum height=0.8cm,align=center}]
 \node[b] (p) {Program\\(Rust, Circom,\\Cairo,\ldots)};
 \node[b,right=of p] (c) {Arithmetic\\circuit};
 \node[b,right=of c] (r) {R1CS /\\PLONKish /\\AIR};
 \node[b,right=of r] (q) {QAP /\\polynomial\\identities};
 \node[b,right=of q,fill=green!8] (s) {SNARK /\\STARK};
 \draw[->] (p)--(c);\draw[->] (c)--(r);\draw[->] (r)--(q);\draw[->] (q)--(s);
\end{tikzpicture}
\end{center}

\section{Arithmetic circuits and constraint systems}
\begin{definition}[Arithmetic circuit]
A DAG whose input nodes are labelled by variables or constants of $\FF$ and whose inner nodes are $+$ or $\times$ gates. \textsf{CircuitSAT}: given $C$ and public input $\vec x$ does there exist $\vec w$ with $C(\vec x,\vec w)=0$?
\end{definition}
Every Boolean statement is embedded by means of the constraint $b(1-b)=0$ ($b\in\{0,1\}$): $\mathrm{AND}(a,b)=ab$, $\mathrm{NOT}(a)=1-a$, $\mathrm{XOR}(a,b)=a+b-2ab$, $\mathrm{OR}(a,b)=a+b-ab$. Comparison and range proofs use bit decomposition: $v=\sum_{i<n}2^ib_i$ with $b_i(1-b_i)=0$ proves $0\le v<2^n$ (if $2^n<|\FF|$).

\begin{definition}[R1CS]\label{def:r1cs}
A \emph{Rank-1 Constraint System} over $\FF$ with $m$ variables and $d$ constraints is a triple of matrices $A,B,C\in\FF^{d\times m}$. A vector $\vec a=(1,\vec x,\vec w)\in\FF^m$ (first entry the constant $1$, then public inputs, then private witness) \emph{satisfies} it if for every row $j\in[d]$
\[
\Big(\sum_iA_{ji}a_i\Big)\cdot\Big(\sum_iB_{ji}a_i\Big)=\sum_iC_{ji}a_i,\qquad\text{i.e. }(A\vec a)\circ(B\vec a)=C\vec a.
\]
\end{definition}
Each constraint is ``(linear form) $\times$ (linear form) = (linear form)''; additions are free. R1CS is equivalent to circuits up to a constant factor.

\begin{example}[Running example: ``I know $x$ with $x^3+x+5=35$'']\label{ex:r1cs}
Flatten into gates:
\[ v_1=x\cdot x,\quad v_2=v_1\cdot x,\quad v_3=(v_2+x)\cdot1,\quad \mathrm{out}=(v_3+5)\cdot1. \]
Variable vector $\vec a=(1,\mathrm{out},x,v_1,v_2,v_3)$; public inputs $(1,\mathrm{out})$, private $(x,v_1,v_2,v_3)$. The witness $x=3$ gives $\vec a=(1,35,3,9,27,30)$. The four constraints, rows of $A,B,C$ (columns in the order of $\vec a$):
\begin{center}\small
\begin{tabular}{c|c|c|c|l}
gate & $A_j$ & $B_j$ & $C_j$ & constraint\\\hline
1 & $(0,0,1,0,0,0)$ & $(0,0,1,0,0,0)$ & $(0,0,0,1,0,0)$ & $x\cdot x=v_1$\\
2 & $(0,0,0,1,0,0)$ & $(0,0,1,0,0,0)$ & $(0,0,0,0,1,0)$ & $v_1\cdot x=v_2$\\
3 & $(0,0,1,0,1,0)$ & $(1,0,0,0,0,0)$ & $(0,0,0,0,0,1)$ & $(x+v_2)\cdot1=v_3$\\
4 & $(5,0,0,0,0,1)$ & $(1,0,0,0,0,0)$ & $(0,1,0,0,0,0)$ & $(5+v_3)\cdot1=\mathrm{out}$
\end{tabular}
\end{center}
Check: gate 1: $3\cdot3=9$; gate 2: $9\cdot3=27$; gate 3: $(3+27)\cdot1=30$; gate 4: $(5+30)\cdot1=35$. \checkmark
\end{example}

\section{Quadratic Arithmetic Programs}
R1CS has $d$ constraints that a verifier would check one by one---linear work. QAPs (Gennaro--Gentry--Parno--Raykova 2013) encode all $d$ constraints in a \emph{single} polynomial divisibility.
\begin{definition}[QAP]\label{def:qap}
Fix distinct points $\omega_1,\dots,\omega_d\in\FF$ and $t(X)=\prod_j(X-\omega_j)$. For each variable index $i\in\{0..m-1\}$ let $u_i(X),v_i(X),w_i(X)$ be the unique polynomials of degree $<d$ with
\[u_i(\omega_j)=A_{ji},\quad v_i(\omega_j)=B_{ji},\quad w_i(\omega_j)=C_{ji}\qquad(j\in[d]).\]
For a vector $\vec a$ define $A(X)=\sum_ia_iu_i(X)$, $B(X)=\sum a_iv_i(X)$, $C(X)=\sum a_iw_i(X)$.
\end{definition}
\begin{theorem}[R1CS $\Leftrightarrow$ QAP]\label{thm:qap}
$\vec a$ satisfies the R1CS if and only if $t(X)$ divides $P(X):=A(X)B(X)-C(X)$. Equivalently, there is $h(X)$ of degree $\le d-2$ with $A(X)B(X)-C(X)=h(X)t(X)$.
\end{theorem}
\begin{proof}
For each $j$, $P(\omega_j)=\big(\sum_ia_iA_{ji}\big)\big(\sum_ia_iB_{ji}\big)-\sum_ia_iC_{ji}$, which is zero exactly when constraint $j$ holds. A polynomial vanishes on all $\omega_j$ iff it is divisible by $\prod_j(X-\omega_j)=t(X)$ (Lemma~\ref{lem:factor} repeatedly). Degrees: $\deg A,\deg B\le d-1$, so $\deg P\le2d-2$ and $\deg h\le d-2$.
\end{proof}
\begin{intuition}
Think of $\omega_1,\dots,\omega_d$ as ``time stamps'' of the gates. $A(X)$ is the left operand wire of the gate active at ``time'' $X$, etc. $A\cdot B-C$ is the ``error polynomial''; it vanishes on all time stamps iff all gates are correct.
\end{intuition}
\begin{example}[QAP for Example~\ref{ex:r1cs}]
Take $\omega=(1,2,3,4)$, $t(X)=(X-1)(X-2)(X-3)(X-4)=X^4-10X^3+35X^2-50X+24$. The column of $A$ for variable $x$ (index 2) is $(1,0,1,0)^\top$, so $u_2$ is the cubic through $(1,1),(2,0),(3,1),(4,0)$:
\[u_2(X)=-\tfrac23X^3+5X^2-\tfrac{34}3X+8.\]
Similarly $v_2(X)=\tfrac13X^3-\tfrac52X^2+\tfrac{31}6X-2$ (column $(1,1,0,0)$) and $w_1(X)=\tfrac16X^3-X^2+\tfrac{11}6X-1$ (column $(0,0,0,1)$). Summing with the witness $(1,35,3,9,27,30)$:
\begin{align*}
A(X)&=-\tfrac{31}{6}X^3+\tfrac{77}{2}X^2-\tfrac{220}{3}X+43, & B(X)&=\tfrac23X^3-5X^2+\tfrac{31}3X-3,\\
C(X)&=\tfrac{17}{6}X^3-\tfrac{49}{2}X^2+\tfrac{215}3X-41.
\end{align*}
Spot checks: $A(1)=3,B(1)=3,C(1)=9$ (gate~1: $3\cdot3=9$). Polynomial division gives remainder $0$ and
\[h(X)=\frac{A(X)B(X)-C(X)}{t(X)}=-\tfrac{31}{9}X^2+\tfrac{307}{18}X-\tfrac{11}{3}.\]
A wrong witness, e.g.\ $v_2=28$, makes $A(3)B(3)-C(3)\ne0$ so the division leaves a nonzero remainder.
\end{example}
\paragraph{Probabilistic check.} A prover who claims $\vec a$ satisfies the R1CS gives $h$; the verifier picks random $\tau\in\FF$ and checks $A(\tau)B(\tau)-C(\tau)=h(\tau)t(\tau)$. If the claim were false, $P-ht$ would be a nonzero polynomial of degree $\le2d-2$, so by Schwartz--Zippel (Lemma~\ref{lem:sz}) the check passes with probability $\le(2d-2)/|\FF|$. The chapters that follow differ in \emph{how} the evaluation at the hidden point $\tau$ is made succinct, non-interactive and zero-knowledge.

\section{PLONKish arithmetisation}\label{sec:plonkish}
PLONK replaces R1CS by a single universal gate type plus a \emph{permutation} argument for wiring:
\[q_L\,a+q_R\,b+q_O\,c+q_M\,ab+q_C=0.\]
Each row $i$ of the \emph{execution trace} has wire values $(a_i,b_i,c_i)$ and fixed selectors $(q_L,q_R,q_O,q_M,q_C)_i$ (circuit-specific, public). Addition gate: $q_L=q_R=1,q_O=-1$. Multiplication gate: $q_M=1,q_O=-1$. Constant check: $q_L=1,q_C=-k,\ldots$
\begin{example}[Example~\ref{ex:r1cs} in PLONK form]
\begin{center}\small
\begin{tabular}{c|ccc|ccccc|l}
row & $a$ & $b$ & $c$ & $q_L$ & $q_R$ & $q_O$ & $q_M$ & $q_C$ & meaning\\\hline
1 & $x$ & $x$ & $v_1$ & 0&0&$-1$&1&0 & $x\cdot x-v_1=0$\\
2 & $v_1$ & $x$ & $v_2$ & 0&0&$-1$&1&0 & $v_1x-v_2=0$\\
3 & $v_2$ & $x$ & $v_3$ & 1&1&$-1$&0&0 & $v_2+x-v_3=0$\\
4 & $v_3$ & 0 & $\mathrm{out}$ & 1&0&$-1$&0&5 & $v_3+5-\mathrm{out}=0$
\end{tabular}
\end{center}
With the witness: $(a,b,c)$ rows are $(3,3,9)$, $(9,3,27)$, $(27,3,30)$, $(30,0,35)$. Row-wise constraints hold. What is \emph{not} yet enforced is that the same variable is used in different cells: $x$ appears in $a_1,b_1,b_2,b_3$, and $c_1=a_2$, $c_2=a_3$, $c_3=a_4$. These are \emph{copy constraints}, enforced by a permutation $\sigma$ on cells (Chapter~\ref{ch:snark}).
\end{example}

\section{Algebraic Intermediate Representation (AIR)}
For long, uniform computations (a CPU or VM running $T$ steps), STARKs use an \emph{execution trace}: a matrix with $T$ rows (time) and $w$ columns (registers), plus
\begin{itemize}
\item \textbf{transition constraints}: low-degree polynomials $\mathcal P_k(\vec s_i,\vec s_{i+1})=0$ relating a row to the next;
\item \textbf{boundary constraints}: initial and final values.
\end{itemize}
\begin{example}[Fibonacci]
Single column $s_0..s_7=1,1,2,3,5,8,13,21$. Claim: ``the 8th Fibonacci number is $21$''. Transition: $s_{i+2}-s_{i+1}-s_i=0$ for $i=0..5$. Boundary: $s_0=1$, $s_1=1$, $s_7=21$. Let $\omega$ be a generator of the multiplicative subgroup of order $8$, $f(\omega^i)=s_i$ (interpolate). The transition constraint becomes the polynomial identity $f(\omega^2X)-f(\omega X)-f(X)=0$ for $X\in\{\omega^0,\dots,\omega^5\}$, i.e.\ the polynomial
\[Q_{\mathrm{tr}}(X)=\frac{f(\omega^2X)-f(\omega X)-f(X)}{\prod_{i=0}^{5}(X-\omega^i)}\]
must be a \emph{polynomial} (of low degree). Similarly $Q_{\mathrm{bd}}(X)=(f(X)-1)/(X-\omega^0)$ etc. The prover commits to $f$ and to the quotients; the verifier checks that each is of low degree via FRI (Chapter~\ref{ch:pcs}).
\end{example}
\begin{remark}[Comparison]
R1CS/QAP: unstructured, flexible, constant-size proofs with pairings. PLONKish: custom gates and lookups, universal SRS. AIR: uniform/repeating computations, transparent (hash-based). Hybrids exist: Halo2, Plonky2/3 (AIR over small fields + FRI), Nova/HyperNova (folding R1CS instances).
\end{remark}

\section*{Exercises}
\begin{exercise}\label{ex:r1csex}
Write an R1CS (give $A,B,C$) for ``$y=x^4$'' and for the Boolean constraint ``$b\in\{0,1\}$ and $z=b\cdot c_1+(1-b)\cdot c_0$'' (a multiplexer).
\end{exercise}
\begin{exercise}
In the toy QAP above, compute the polynomial $A(X)B(X)-C(X)$ modulo $t(X)$ for the wrong witness $(1,35,3,9,28,30)$ and find which gates it violates.
\end{exercise}
\begin{exercise}
Prove that for $d$ points and $\omega_j$ the $d$-th roots of unity, $t(X)=X^d-1$, and that Lagrange interpolation can be done by an inverse FFT in $O(d\log d)$.
\end{exercise}
\begin{exercise}
Show that R1CS satisfiability is $\NP$-complete by reducing from \textsf{CircuitSAT}.
\end{exercise}

\chapter{Polynomial Commitment Schemes}
\label{ch:pcs}
A polynomial commitment scheme (PCS) lets a prover commit to a polynomial $f$ with a short string $C$ and later prove, with a short proof, that $f(z)=y$ for a point $z$ chosen by the verifier. PCS are the cryptographic core of modern proof systems: an \emph{information-theoretic} proof system (a ``polynomial IOP'') in which the verifier queries oracles to polynomials becomes a real argument by replacing the oracles with commitments (and Fiat--Shamir).
\begin{definition}[Polynomial commitment scheme]
A tuple $(\Setup,\mathsf{Commit},\mathsf{Open},\mathsf{Verify})$:
\begin{itemize}
\item $\Setup(1^\lambda,D)\to\mathsf{pp}$ for degree bound $D$;
\item $\mathsf{Commit}(\mathsf{pp},f)\to C$;
\item $\mathsf{Open}(\mathsf{pp},f,z)\to(y,\pi)$ with $y=f(z)$;
\item $\mathsf{Verify}(\mathsf{pp},C,z,y,\pi)\to\{0,1\}$.
\end{itemize}
\textbf{Correctness}; \textbf{evaluation binding} (no PPT adversary opens one $C$ to $(z,y)$ and $(z,y')$, $y\ne y'$, both accepted); \textbf{knowledge soundness} (an extractor recovers an $f$ of degree $\le D$ consistent with all accepted openings); optionally \textbf{hiding} (commitment and openings reveal nothing about $f$ beyond evaluated points).
\end{definition}
The main PCS families:
\begin{center}\small
\begin{tabular}{@{}p{3.1cm}p{2.6cm}p{3.2cm}p{2.8cm}p{2.4cm}@{}}\toprule
Scheme & Setup & Proof size & Verifier & Assumption\\\midrule
KZG & trusted (universal) & $O(1)$ (1 group element) & $O(1)$ (2 pairings) & $q$-SDH / AGM\\
IPA (Bulletproofs/Halo) & transparent & $O(\log n)$ & $O(n)$ (amortisable) & DL\\
FRI & transparent & $O(\log^2n)$ & $O(\log^2n)$ & collision-resistant hash (PQ)\\
Dory / Hyrax / Brakedown & transparent & varies & varies & DLIN / hash / codes\\\bottomrule
\end{tabular}
\end{center}

\section{KZG commitments (Kate--Zaverucha--Goldberg 2010)}
\subsection{The construction}
Let $e:\GG_1\times\GG_2\to\GG_T$ be a pairing of prime order $q$ and generators $g_1,g_2$; bracket notation $[a]_i=g_i^a$.
\begin{algorithm}[label=pr:kzg]{KZG polynomial commitment}
\textbf{Setup} ($D$): sample the \emph{toxic waste} $\tau\getsr\FF_q$ and publish the structured reference string
\[\mathsf{srs}=\big([1]_1,[\tau]_1,[\tau^2]_1,\dots,[\tau^D]_1;\ [1]_2,[\tau]_2\big),\]
then \textbf{delete} $\tau$.\\
\textbf{Commit}$(f=\sum_{i\le D}f_iX^i)$: $C:=\sum_if_i[\tau^i]_1=[f(\tau)]_1$ (a multi-scalar multiplication, no knowledge of $\tau$ needed).\\
\textbf{Open}$(f,z)$: $y=f(z)$; quotient $\psi(X)=\dfrac{f(X)-y}{X-z}$ (a polynomial of degree $\le D-1$ by Lemma~\ref{lem:factor}); $\pi:=[\psi(\tau)]_1$.\\
\textbf{Verify}$(C,z,y,\pi)$: accept iff
\[e\big(C-[y]_1,\,[1]_2\big)=e\big(\pi,\,[\tau]_2-[z]_2\big).\]
\end{algorithm}
\begin{theorem}\label{thm:kzg}
KZG is correct. It is evaluation-binding under the $q$-strong Diffie--Hellman assumption ($q$-SDH): given $([\tau^i]_1,[\tau^i]_2)_{i\le D}$ no PPT algorithm outputs $(z,[1/(\tau-z)]_1)$. Knowledge soundness holds in the algebraic group model.
\end{theorem}
\begin{proof}
\emph{Correctness.} $f(X)-y=(X-z)\psi(X)$, evaluate at $\tau$: $f(\tau)-y=(\tau-z)\psi(\tau)$. Taking pairings, $e([f(\tau)-y]_1,[1]_2)=e([\psi(\tau)]_1,[\tau-z]_2)$, which is the verification equation.

\emph{Binding.} Suppose $\mathcal A$ outputs $C,z$ and two valid openings $(y,\pi)$, $(y',\pi')$ with $y\ne y'$. Then
$e(C-[y]_1,[1]_2)=e(\pi,[\tau-z]_2)$ and $e(C-[y']_1,[1]_2)=e(\pi',[\tau-z]_2)$. Dividing, $e([y'-y]_1,[1]_2)=e(\pi-\pi',[\tau-z]_2)$; so $\pi-\pi'=\big[\tfrac{y'-y}{\tau-z}\big]_1$, hence $\Big[\tfrac1{\tau-z}\Big]_1=(y'-y)^{-1}(\pi-\pi')$ is computed, breaking $q$-SDH.
\end{proof}
\begin{intuition}
The commitment is a \emph{point evaluation at a secret place} $\tau$. Since the prover does not know $\tau$, the only way she can produce $[\psi(\tau)]_1$ with the property $\psi(X)(X-z)=f(X)-y$ is to know an actual quotient polynomial, which exists only if $f(z)=y$ really holds. The pairing checks one multiplicative relation between three hidden numbers.
\end{intuition}
\begin{example}[KZG ``in the exponent'']
Work with scalars modulo $q=11$ (we show the exponents; $[a]_1$ corresponds to $a$). Let $f(X)=X^2+1$, $z=2$, $y=f(2)=5$. Quotient: $(X^2+1-5)/(X-2)=(X^2-4)/(X-2)=X+2=:\psi(X)$. With secret $\tau=7$: commitment exponent $f(\tau)=50\equiv6$; proof exponent $\psi(\tau)=9$. The pairing equation in the exponent reads $f(\tau)-y=(\tau-z)\psi(\tau)$: LHS $6-5=1$; RHS $(7-2)\cdot9=45\equiv1\pmod{11}$. \checkmark
A false claim $y'=4$ would require a proof element $\pi'$ with $(\tau-z)\pi'=f(\tau)-y'=2$, i.e.\ $\pi'=2/(\tau-z)$. The element $[1/(\tau-z)]_1$ is not a polynomial expression of degree $\le D$ in $\tau$, so it cannot be formed from the SRS; computing it is exactly the $q$-SDH problem.
\end{example}

\subsection{Extensions}
\begin{itemize}
\item \textbf{Batch opening at many points.} To prove $f(z_i)=y_i$ for $i\in[k]$ let $I(X)$ interpolate $(z_i,y_i)$ and $Z(X)=\prod(X-z_i)$; then $\psi=(f-I)/Z$, $\pi=[\psi(\tau)]_1$, check $e(C-[I(\tau)]_1,[1]_2)=e(\pi,[Z(\tau)]_2)$.
\item \textbf{Several polynomials at the same point.} Random linear combination $\sum_j\gamma^jf_j$ (challenge $\gamma$): by Schwartz--Zippel an incorrect combination is caught except with probability $k/|\FF|$. One proof for all.
\item \textbf{Hiding.} Add a random multiple of a blinding polynomial: $C=[f(\tau)+\rho\,\eta(\tau)]_1$ with an additional SRS element.
\item \textbf{Universal setup.} The SRS depends only on a degree bound $D$, not on the circuit; the ``Powers of Tau'' ceremony has thousands of contributors, and security needs only \emph{one} honest participant.
\end{itemize}

\section{FRI: a hash-based low-degree test}
\subsection{Reed--Solomon codes and proximity}
Evaluating a polynomial of degree $<k$ on a domain $D\subseteq\FF$ with $|D|=n=k/\rho$ ($\rho$ = \emph{rate}, e.g.\ $\rho=1/4$) gives a codeword of the Reed--Solomon code $\mathrm{RS}[\FF,D,\rho]$. A function $f:D\to\FF$ is $\delta$-\emph{far} if every codeword differs in at least $\delta n$ positions. FRI (Fast Reed--Solomon IOP of Proximity, Ben-Sasson--Bentov--Horesh--Riabzev) allows a verifier to distinguish ``$f$ is a codeword'' from ``$f$ is $\delta$-far'' with $O(\log n)$ queries to a Merkle commitment.

\subsection{Folding}
Any polynomial splits into even and odd parts: $f(X)=f_e(X^2)+X\,f_o(X^2)$ with $\deg f_e,\deg f_o<k/2$. For a verifier challenge $\beta$ define
\[f^{(1)}(Y)=f_e(Y)+\beta f_o(Y).\]
Then $\deg f^{(1)}<k/2$ and, at $Y=x^2$ for $x\in D$,
\[f^{(1)}(x^2)=\frac{f(x)+f(-x)}2+\beta\,\frac{f(x)-f(-x)}{2x}.\]
So the folded evaluation on $D^2=\{x^2\}$ (half the size) is computable from just two values of $f$. If $f$ is far from low degree, then so is $f^{(1)}$ for most $\beta$ (the ``proximity gap'' lemma).
\begin{algorithm}[label=pr:fri]{FRI (commit and query phases)}
\textbf{Commit phase.} $f^{(0)}=f$ on $D_0=D$; Merkle-commit to $f^{(0)}$. For $r=0..R-1$: $\Ver$ sends $\beta_r$; $\Prv$ computes $f^{(r+1)}=\mathrm{Fold}(f^{(r)},\beta_r)$ on $D_{r+1}=D_r^2$ and commits. Finally $\Prv$ sends the (constant or small) last polynomial in the clear.\\
\textbf{Query phase} (repeated $s$ times): $\Ver$ picks a random $x\in D_0$. For each $r$ $\Prv$ opens (with Merkle paths) $f^{(r)}(x_r)$ and $f^{(r)}(-x_r)$ ($x_r=x^{2^r}$), and $\Ver$ checks the folding equation $f^{(r+1)}(x_r^2)=\frac{f^{(r)}(x_r)+f^{(r)}(-x_r)}2+\beta_r\frac{f^{(r)}(x_r)-f^{(r)}(-x_r)}{2x_r}$, and finally the last-layer value.
\end{algorithm}
\begin{center}
\begin{tikzpicture}[every node/.style={font=\small}]
 \foreach \r/\n/\c in {0/8/blue!15,1/4/blue!25,2/2/blue!35,3/1/blue!45}{
   \foreach \i in {1,...,\n}{
     \pgfmathsetmacro{\w}{8/\n}
     \draw[fill=\c] ({(\i-1)*\w},-\r*1.1) rectangle ({\i*\w},{-\r*1.1-0.6});}
   \node[left] at (-0.1,{-\r*1.1-0.3}) {$f^{(\r)}$ ($|D_\r|=\n$)};}
 \draw[->,thick] (4.2,-0.65)--(4.2,-1.05) node[midway,right] {\scriptsize fold, $\beta_0$};
 \draw[->,thick] (4.2,-1.75)--(4.2,-2.15) node[midway,right] {\scriptsize fold, $\beta_1$};
 \draw[->,thick] (4.2,-2.85)--(4.2,-3.25) node[midway,right] {\scriptsize fold, $\beta_2$};
\end{tikzpicture}
\end{center}
\begin{example}[One folding step, $\FF_{17}$]
Let $\FF=\FF_{17}$, $D_0=\{1,2,4,8,9,13,15,16\}$ (the quadratic residues, order 8), $D_1=D_0^2=\{1,4,13,16\}$. Take $f(X)=X^3+2X^2+3X+4$ ($k=4$, rate $1/2$). Then $f_e(Y)=2Y+4$ (even-degree terms $2X^2+4$ with $Y=X^2$) and $f_o(Y)=Y+3$ (as $X^3+3X=X(X^2+3)$). With $\beta=3$: $f^{(1)}(Y)=2Y+4+3(Y+3)=5Y+13$, degree $1<k/2=2$.\\
Consistency check at $x=2$ (so $x^2=4$): $f(2)=8+8+6+4=26\equiv9$, $f(-2)=f(15)=-8+8-6+4=-2\equiv15$. Then $\frac{9+15}2=12$ and $\frac{9-15}{2\cdot2}=\frac{-6}{4}\equiv-6\cdot13=-78\equiv7$ (as $4\cdot13=52\equiv1$), so $f^{(1)}(4)=12+3\cdot7=33\equiv16$. Directly: $5\cdot4+13=33\equiv16$. \checkmark\\
After $\log_2k$ folds the polynomial is constant and the verifier can check it directly.
\end{example}
\begin{theorem}[FRI soundness, informal]
If $f$ is $\delta$-far from $\mathrm{RS}[\FF,D,\rho]$ then, for $\delta$ below the Johnson bound, the verifier rejects with probability $\ge1-(1-\delta)^s-\varepsilon_{\mathrm{fold}}$ where $s$ is the number of queries and $\varepsilon_{\mathrm{fold}}=O(n/|\FF|)$ accounts for unlucky folding challenges. Querying $s\approx\lambda/\log_2(1/(1-\delta))$ locations gives $\lambda$ bits of security; smaller rate $\rho$ means larger $\delta$, i.e.\ fewer queries but a larger domain.
\end{theorem}
FRI turns a Merkle commitment into a polynomial commitment: to open $f(z)=y$, run FRI on the quotient $(f(X)-y)/(X-z)$; the check that this is low-degree convinces the verifier $f(z)=y$ (otherwise the quotient has a pole at $z$).

\section{Inner-product arguments and Bulletproofs}
\subsection{Pedersen vector commitment}
For generators $\vec g=(g_1..g_n)\in\GG^n$ with unknown relations, $\mathsf{Com}(\vec a;r)=\prod_ig_i^{a_i}\cdot h^r$ commits to a vector; binding under DL; perfectly hiding.
\subsection{The core: proving an inner product in $O(\log n)$}
Statement: given $P=\vec g^{\vec a}\vec h^{\vec b}u^{\langle\vec a,\vec b\rangle}$ ($\vec g^{\vec a}:=\prod g_i^{a_i}$), prove knowledge of $\vec a,\vec b\in\ZZ_q^n$ ($n$ a power of two). Naively send $\vec a,\vec b$ ($2n$ elements); the trick is \emph{recursive folding}.
\begin{protocol}[label=pr:ipa]{Inner-product argument (one folding round)}
Write $\vec a=(\vec a_L\|\vec a_R)$ and likewise $\vec b,\vec g,\vec h$ (halves of size $n/2$). 
\begin{flow}{Prover}{Verifier}
$L=\vec g_R^{\,\vec a_L}\,\vec h_L^{\,\vec b_R}\,u^{\langle\vec a_L,\vec b_R\rangle}$ & & \\
$R=\vec g_L^{\,\vec a_R}\,\vec h_R^{\,\vec b_L}\,u^{\langle\vec a_R,\vec b_L\rangle}$ & $\ra{L,R}$ & \\
 & $\la{x}$ & $x\getsr\ZZ_q^*$\\
$\vec a'=x\vec a_L+x\inv\vec a_R$ & & $\vec g'=\vec g_L^{\,x\inv}\circ\vec g_R^{\,x}$\\
$\vec b'=x\inv\vec b_L+x\vec b_R$ & & $\vec h'=\vec h_L^{\,x}\circ\vec h_R^{\,x\inv}$\\
 & & $P'=L^{x^2}\,P\,R^{x^{-2}}$
\end{flow}
Now the same statement is claimed for $(\vec g',\vec h',P',\vec a',\vec b')$ with length $n/2$. Recurse until $n=1$, then the prover sends the two scalars $a,b$ and the verifier checks $P=g^ah^bu^{ab}$.
\end{protocol}
\begin{lemma}[Invariant]
If $P=\vec g^{\vec a}\vec h^{\vec b}u^{\langle\vec a,\vec b\rangle}$ then $P'=\vec g'^{\vec a'}\vec h'^{\vec b'}u^{\langle\vec a',\vec b'\rangle}$.
\end{lemma}
\begin{proof}
Expand:
\begin{align*}
\vec g'^{\vec a'}&=\vec g_L^{\,\vec a_L}\ \vec g_R^{\,\vec a_R}\ \vec g_L^{\,x^{-2}\vec a_R}\ \vec g_R^{\,x^2\vec a_L},\\
\vec h'^{\vec b'}&=\vec h_L^{\,\vec b_L}\ \vec h_R^{\,\vec b_R}\ \vec h_L^{\,x^2\vec b_R}\ \vec h_R^{\,x^{-2}\vec b_L},\\
u^{\langle\vec a',\vec b'\rangle}&=u^{\langle\vec a,\vec b\rangle+x^2\langle\vec a_L,\vec b_R\rangle+x^{-2}\langle\vec a_R,\vec b_L\rangle}.
\end{align*}
Collect: $\vec g^{\vec a}\vec h^{\vec b}u^{\langle\vec a,\vec b\rangle}\cdot\big(\vec g_R^{\vec a_L}\vec h_L^{\vec b_R}u^{\langle\vec a_L,\vec b_R\rangle}\big)^{x^2}\cdot\big(\vec g_L^{\vec a_R}\vec h_R^{\vec b_L}u^{\langle\vec a_R,\vec b_L\rangle}\big)^{x^{-2}}=P\,L^{x^2}R^{x^{-2}}=P'$.
\end{proof}
\begin{example}[Scalar folding in $\ZZ_{11}$]
$\vec a=(1,2,3,4)$, $\vec b=(5,6,7,8)$, $\langle\vec a,\vec b\rangle=70\equiv4$. Halves: $\vec a_L=(1,2)$, $\vec a_R=(3,4)$, $\vec b_L=(5,6)$, $\vec b_R=(7,8)$; cross terms $\langle\vec a_L,\vec b_R\rangle=23$, $\langle\vec a_R,\vec b_L\rangle=39$. Challenge $x=2$, $x\inv=6$, $x^2=4$, $x^{-2}=3$.\\
$\vec a'=2(1,2)+6(3,4)=(20,28)\equiv(9,6)$; $\vec b'=6(5,6)+2(7,8)=(44,52)\equiv(0,8)$; $\langle\vec a',\vec b'\rangle=48\equiv4$. Predicted by the invariant: $\langle\vec a,\vec b\rangle+4\cdot23+3\cdot39=70+92+117=279=25\cdot11+4\equiv4$. \checkmark\\
Two vectors of length 4 became vectors of length 2 at the cost of two group elements $L,R$.
\end{example}
\begin{theorem}
The protocol has $\log_2n$ rounds, proof size $2\log_2n$ group elements plus $2$ scalars, and is an argument of knowledge under DL (extraction by $3$-special soundness per round: transcripts for three challenges $x_1,x_2,x_3$ yield a ``witness'' of length $n$ via a linear system, forming a tree of $3^{\log n}=n^{\log3}$ transcripts---polynomial).
\end{theorem}
\subsection{Application: range proofs (Bulletproofs, B\"unz et al.\ 2018)}
To show that a Pedersen commitment $V=g^vh^\gamma$ hides $v\in[0,2^n)$:
\begin{enumerate}
\item Let $\vec a_L\in\{0,1\}^n$ be the bits of $v$ and $\vec a_R=\vec a_L-\vec1^n$. The statement is equivalent to three conditions: $\langle\vec a_L,2^n\rangle=v$, $\ \vec a_L\circ\vec a_R=\vec0$ (each bit is $0$ or $1$), $\ \vec a_L-\vec a_R=\vec1^n$. Here $2^n=(1,2,4,\dots,2^{n-1})$.
\item The verifier sends random $y,z$. By Schwartz--Zippel the three conditions hold iff (w.h.p.) the single scalar identity
\[z^2v+\delta(y,z)=\big\langle\vec a_L-z\vec1^n,\ \vec y^{\,n}\circ(\vec a_R+z\vec1^n)+z^2 2^n\big\rangle,\qquad\delta(y,z)=(z-z^2)\langle\vec1^n,\vec y^{\,n}\rangle-z^3\langle\vec1^n,2^n\rangle\]
holds, where $\vec y^{\,n}=(1,y,\dots,y^{n-1})$.
\item To hide the vectors, the prover adds blinding vectors $\vec s_L,\vec s_R$ and defines the degree-one vector polynomials
\[\vec l(X)=(\vec a_L-z\vec1^n)+\vec s_LX,\qquad \vec r(X)=\vec y^{\,n}\circ(\vec a_R+z\vec1^n+\vec s_RX)+z^22^n,\]
with $t(X)=\langle\vec l(X),\vec r(X)\rangle=t_0+t_1X+t_2X^2$ and $t_0=z^2v+\delta(y,z)$ exactly when the statement is true.
\item The prover commits to $t_1,t_2$ (Pedersen commitments $T_1,T_2$), receives a challenge $x$, and reveals $\vec l=\vec l(x)$, $\vec r=\vec r(x)$, $\hat t=t(x)$ with blinding factors. The verifier checks $g^{\hat t}h^{\tau_x}=V^{z^2}g^{\delta(y,z)}T_1^xT_2^{x^2}$ (which checks $t_0$ and the form of $t$ at a random point), and that $\vec l,\vec r$ are consistent with the commitments to $\vec a_L,\vec a_R,\vec s_L,\vec s_R$ and satisfy $\hat t=\langle\vec l,\vec r\rangle$---the last by the inner-product argument, $O(\log n)$ communication.
\end{enumerate}
\begin{example}[Checking the identity in step 2, $n=2$, over $\ZZ$]
Let $v=2$ (bits $\vec a_L=(0,1)$, $\vec a_R=(-1,0)$), $y=2$, $z=3$. Then $\vec y^{\,2}=(1,2)$, $2^2=(1,2)$, $\langle\vec1,\vec y^{\,2}\rangle=3$, $\langle\vec1,2^2\rangle=3$, so $\delta=(3-9)\cdot3-27\cdot3=-18-81=-99$ and LHS $=9\cdot2-99=-81$. RHS: $\vec a_L-z\vec1=(-3,-2)$; $\vec a_R+z\vec1=(2,3)$, $\vec y\circ(\cdot)=(2,6)$, plus $z^2 2^2=(9,18)$ gives $(11,24)$; inner product $-33-48=-81$. \checkmark
\end{example}
A 64-bit range proof is $\approx700$ bytes; aggregating $m$ proofs adds only $2\log_2m$ group elements. Used in Monero, Mimblewimble and confidential assets.

\section*{Exercises}
\begin{exercise}
Prove that if the prover commits to a polynomial of degree $>D$ the KZG verification fails for most $z$ or the prover cannot form $\pi$ (hint: the SRS only contains powers up to $\tau^D$; discuss degree enforcement through a second commitment shifted by $\tau^{D-d}$).
\end{exercise}
\begin{exercise}
Derive the KZG multi-point opening verification equation, and the batched version for several polynomials at the same point.
\end{exercise}
\begin{exercise}\label{ex:foldcheck}
Check, for $f(X)=X^3+2X^2+3X+4$ over $\FF_{17}$ with $\beta=3$, that the folding relation holds at $x=4$ as well (compute $f(4),f(13)$, then $f^{(1)}(16)$).
\end{exercise}
\begin{exercise}
Verify the final round of the inner-product argument for $n=2$: show that with $\vec a=(a_1,a_2)$, $\vec b=(b_1,b_2)$ and challenge $x$, $P'$ equals $g'^{a'}h'^{b'}u^{a'b'}$.
\end{exercise}
\begin{exercise}
Explain why IPA verification is $O(n)$ (hint: the folded generators $\vec g'$ are the verifier's burden) and how a single multi-exponentiation with scalars $s_i=\prod_jx_j^{\pm1}$ resolves it.
\end{exercise}

\chapter{SNARKs and STARKs}
\label{ch:snark}
\section{Definitions}
\begin{definition}[zk-SNARK]\label{def:snark}
A \emph{zk-SNARK} (zero-knowledge Succinct Non-interactive ARgument of Knowledge) for an $\NP$ relation $R$ is a triple $(\mathsf{Setup},\mathsf{Prove},\mathsf{Verify})$ with $\mathsf{Setup}(R)\to(\mathsf{pk},\mathsf{vk})$ (the CRS), $\mathsf{Prove}(\mathsf{pk},x,w)\to\pi$, $\mathsf{Verify}(\mathsf{vk},x,\pi)\to\{0,1\}$ such that
\begin{description}
\item[Completeness] honest proofs verify;
\item[Knowledge soundness] for every PPT prover $\mathcal A$ there is a PPT extractor $\Ext_{\mathcal A}$ (non-black-box; given $\mathcal A$'s coins) such that whenever $\mathcal A$ outputs an accepting $(x,\pi)$ then $\Ext_{\mathcal A}$ outputs $w$ with $(x,w)\in R$, except with negligible probability;
\item[Zero-knowledge] there is a simulator with trapdoor that outputs proofs indistinguishable from real ones (for the same $x$);
\item[Succinctness] $|\pi|=\poly(\lambda)$ independent of $|w|$ (typically $O(1)$ or $O(\log|C|)$) and verification time $\poly(\lambda,|x|)$ (+ $\polylog|C|$).
\end{description}
\end{definition}
\begin{remark}[Why non-black-box extraction?]
Gentry--Wichs (2011) show that succinct non-interactive arguments with black-box (falsifiable) soundness for $\NP$ cannot be proven secure under falsifiable assumptions. So SNARK security rests on knowledge assumptions (KEA, AGM, generic group model) or the random oracle model.
\end{remark}

\section{Groth16}
Groth16 \cite{Gro16} is the pairing-based zk-SNARK for QAPs with the smallest known proof size (3 group elements: $\approx128$ bytes on BN254 / 192 bytes on BLS12-381) and cheapest verification (3 pairings), and is deployed in Zcash Sapling, Tornado Cash and many rollups.

\subsection{Algorithm}
Let the QAP be $(\{u_i,v_i,w_i\}_{i=0}^{m-1},t)$ over $\FF_q$ for the relation, with the first $\ell$ variables public ($a_0=1$ and the statement). Let $\pi=(A,B,C)$.
\begin{algorithm}[label=pr:groth16]{Groth16}
\textbf{Setup.} Sample toxic waste $\tau,\alpha,\beta,\gamma,\delta\getsr\FF_q^*$. Set
\begin{itemize}[nosep]
\item $\mathsf{pk}\supseteq[\alpha]_1,[\beta]_1,[\delta]_1,[\beta]_2,[\delta]_2,\ \{[\tau^k]_1\}_{k<d},\ \{[\tau^k]_2\}_{k<d}$,
\item $\mathsf{pk}\supseteq\Big\{\Big[\tfrac{\beta u_i(\tau)+\alpha v_i(\tau)+w_i(\tau)}{\delta}\Big]_1\Big\}_{i\ge\ell}$ (private variables), $\ \Big\{\Big[\tfrac{\tau^kt(\tau)}{\delta}\Big]_1\Big\}_{k\le d-2}$,
\item $\mathsf{vk}=\Big([\alpha]_1,[\beta]_2,[\gamma]_2,[\delta]_2,\ \Big\{\Big[\tfrac{\beta u_i(\tau)+\alpha v_i(\tau)+w_i(\tau)}{\gamma}\Big]_1\Big\}_{i<\ell}\Big)$.
\end{itemize}
\textbf{Prove}$(\mathsf{pk},\vec a)$: compute $h(X)=(A(X)B(X)-C(X))/t(X)$ for the witness polynomials. Sample $r,s\getsr\FF_q$ (\emph{zero-knowledge randomisers}). Output
\begin{align*}
A&=\big[\alpha+\textstyle\sum_ia_iu_i(\tau)+r\delta\big]_1,\qquad B=\big[\beta+\textstyle\sum_ia_iv_i(\tau)+s\delta\big]_2,\\
C&=\Big[\tfrac{\sum_{i\ge\ell}a_i(\beta u_i(\tau)+\alpha v_i(\tau)+w_i(\tau))+h(\tau)t(\tau)}{\delta}+sA_\bullet+rB_\bullet-rs\delta\Big]_1
\end{align*}
where $A_\bullet=\alpha+\sum a_iu_i(\tau)+r\delta$ and $B_\bullet=\beta+\sum a_iv_i(\tau)+s\delta$ are the exponents of $A$ and $B$; $B$ is also needed in $\GG_1$ (publish $[\beta]_1,[v_i(\tau)]_1$ to form it).\\
\textbf{Verify}$(\mathsf{vk},\vec x,\pi)$: with $\mathsf{IC}=\sum_{i<\ell}a_i\Big[\tfrac{\beta u_i+\alpha v_i+w_i}{\gamma}\Big]_1$, accept iff
\[e(A,B)=e([\alpha]_1,[\beta]_2)\cdot e(\mathsf{IC},[\gamma]_2)\cdot e(C,[\delta]_2).\]
\end{algorithm}
\subsection{Correctness}
Write $K_i=\beta u_i(\tau)+\alpha v_i(\tau)+w_i(\tau)$, $\bar A=\sum_ia_iu_i(\tau)$, $\bar B=\sum_ia_iv_i(\tau)$, $\bar C=\sum_ia_iw_i(\tau)$, so that
\[\sum_ia_iK_i=\beta\bar A+\alpha\bar B+\bar C,\qquad h(\tau)t(\tau)=\bar A\bar B-\bar C.\]
Let $\mathcal A=\alpha+\bar A+r\delta$ and $\mathcal B=\beta+\bar B+s\delta$ be the exponents of $A$ and $B$, and $\mathcal C=\frac{\sum_{i\ge\ell}a_iK_i+h(\tau)t(\tau)}{\delta}+s\mathcal A+r\mathcal B-rs\delta$ the exponent of $C$. The right-hand side of the verification equation has exponent (the $\mathsf{IC}$ term contributes $\sum_{i<\ell}a_iK_i$ since it was divided by $\gamma$ and is paired with $\gamma$):
\begin{align*}
\alpha\beta+\sum_{i<\ell}a_iK_i+\delta\mathcal C
&=\alpha\beta+\sum_ia_iK_i+h(\tau)t(\tau)+\delta\big(s\mathcal A+r\mathcal B-rs\delta\big)\\
&=\alpha\beta+\beta\bar A+\alpha\bar B+\bar C+\bar A\bar B-\bar C+\delta\big(s\mathcal A+r\mathcal B-rs\delta\big)\\
&=(\alpha+\bar A)(\beta+\bar B)+\delta\big(s\mathcal A+r\mathcal B-rs\delta\big).
\end{align*}
On the other hand
\[\mathcal A\mathcal B=(\alpha+\bar A)(\beta+\bar B)+s\delta(\alpha+\bar A)+r\delta(\beta+\bar B)+rs\delta^2,\]
and $\delta(s\mathcal A+r\mathcal B-rs\delta)=s\delta(\alpha+\bar A)+rs\delta^2+r\delta(\beta+\bar B)+rs\delta^2-rs\delta^2=s\delta(\alpha+\bar A)+r\delta(\beta+\bar B)+rs\delta^2$. The two expressions coincide, so $e(A,B)$ equals the right-hand side. \hfill$\square$

\begin{example}[Numerical Groth16 ``in the exponent'']\label{ex:groth16}
Take the QAP of Example~\ref{ex:r1cs} over $\FF_{101}$, points $\omega=(1,2,3,4)$ and the witness $\vec a=(1,35,3,9,27,30)$. Public variables: $\ell=2$ ($a_0=1$, $a_1=35$). Setup randomness $\tau=7,\alpha=2,\beta=3,\gamma=5,\delta=11$. Prover randomness $r=4$, $s=9$. We display exponents (a real implementation hides them in $\GG_1,\GG_2$).
\begin{center}\small
\begin{tabular}{c|cccccc}
$i$ & 0 & 1 & 2 & 3 & 4 & 5\\\hline
$u_i(\tau)$ & 100 & 0 & 46 & 36 & 56 & 20\\
$v_i(\tau)$ & 76 & 0 & 26 & 0 & 0 & 0\\
$w_i(\tau)$ & 0 & 20 & 0 & 91 & 36 & 56
\end{tabular}
\end{center}
$t(\tau)=(6)(5)(4)(3)=360\equiv57$. Then $\bar A=48$, $\bar B=53$, $\bar C=30$, $\bar A\bar B-\bar C=2514\equiv90$ and $h(\tau)=90\cdot57^{-1}\equiv76$ (indeed $57\cdot76=4332\equiv90$).\\
\textbf{Proof exponents:} $A=\alpha+\bar A+r\delta=2+48+44=94$; $B=\beta+\bar B+s\delta=3+53+99=155\equiv54$; $C=44$.\\
\textbf{Verifier:} $\mathsf{IC}\cdot\gamma=\sum_{i<2}a_i(\beta u_i+\alpha v_i+w_i)=(300+152)+35\cdot20=452+700\equiv41$. Then
\[AB=94\cdot54=5076\equiv26,\qquad \alpha\beta+\mathsf{IC}\gamma+C\delta=6+41+44\cdot11=6+41+484\equiv26\pmod{101}.\ \checkmark\]
A forged witness e.g.\ $v_2=28$ makes $h$ non-polynomial; the prover cannot produce $C$ without the exact (hidden) $\tau$ powers.
\end{example}

\subsection{Zero-knowledge and soundness}
\begin{theorem}[Groth 2016]
In the generic bilinear group model, the above is a perfectly zero-knowledge, knowledge-sound SNARK for QAPs.
\end{theorem}
\begin{proof}[Proof idea]
\emph{ZK.} Given the trapdoor $(\tau,\alpha,\beta,\gamma,\delta)$ the simulator picks $A,B$ uniformly and sets $C=\frac{AB-\alpha\beta-\mathsf{IC}\gamma}{\delta}$, satisfying the verification equation by definition. A real proof has $A=\alpha+\bar A+r\delta$ and $B=\beta+\bar B+s\delta$ with $r,s$ uniform, so $(A,B)$ is uniform, and $C$ is determined by the verification equation: same distribution.
\emph{Soundness.} Any proof is a generic-group computation; the adversary's $A$, $B$, $C$ are linear combinations of the SRS elements. Matching the formal polynomial identity in the indeterminates $(\tau,\alpha,\beta,\gamma,\delta)$ shows that the coefficient of $\alpha\beta$ forces the shape $A=\alpha+\sum a_iu_i+\dots$; the $\gamma$ and $\delta$ divisions separate public from private parts; finally the polynomial identity holds over the indeterminate $\tau$, so $(\sum a_iu_i)(\sum a_iv_i)-\sum a_iw_i$ is divisible by $t(X)$, i.e.\ $\vec a$ is a QAP witness (Theorem~\ref{thm:qap}), and the extractor reads off $\vec a$ from the coefficients.
\end{proof}
\begin{warning}[Trusted setup]
Whoever knows $\tau,\alpha,\beta,\gamma,\delta$ can forge proofs of false statements (the simulator above). Groth16 requires a \emph{circuit-specific} setup (phase 2 of an MPC ceremony after a universal phase 1 ``Powers of Tau''); as long as one participant destroys her share, the CRS is secure. Another limitation: Groth16 proofs are malleable (re-randomisable) unless bound to a context.
\end{warning}

\section{PLONK}
PLONK \cite{GWC19} uses the PLONKish arithmetisation (Section~\ref{sec:plonkish}) and KZG to obtain a SNARK with \emph{universal and updatable} setup: one SRS for all circuits of size $\le D$.
\subsection{Setup of the argument}
Let $H=\{\omega^0,\dots,\omega^{n-1}\}$ be the subgroup of $n$-th roots of unity in $\FF$ (so $Z_H(X)=X^n-1$). The prover interpolates wire polynomials $a(X),b(X),c(X)$ with $a(\omega^i)=a_i$, etc.; selector polynomials $q_L,q_R,q_O,q_M,q_C$ are fixed by the circuit. The \emph{gate constraint}
\[G(X)=q_La+q_Rb+q_Oc+q_Mab+q_C+\mathrm{PI}(X)\ \equiv\ 0\ \text{ on }H,\]
($\mathrm{PI}$ encodes public inputs) holds iff $Z_H(X)$ divides $G(X)$.
\subsection{Copy constraints: the permutation argument}
Let the $3n$ cells be indexed by $\{1..3n\}$ and $\sigma$ a permutation mapping each cell to the \emph{next} cell containing the same variable (cycles = equivalence classes). Define $f:\{1..3n\}\to\FF$ the cell values. Copy constraints hold iff $f(i)=f(\sigma(i))$ for all $i$.
\begin{lemma}[Multiset equality via grand product]\label{lem:perm}
For random $\beta,\gamma$: $f(i)=f(\sigma(i))\ \forall i$ iff (w.h.p.) $\prod_i\big(f(i)+\beta i+\gamma\big)=\prod_i\big(f(i)+\beta\sigma(i)+\gamma\big)$.
\end{lemma}
\begin{proof}
($\Rightarrow$) If $f=f\circ\sigma$ then the right product is a reordering of the left. ($\Leftarrow$) Consider the pairs $(i,f(i))$ and $(\sigma(i),f(i))$ as multisets. The identity as polynomials in $\beta,\gamma$ is $\prod_i(\gamma+f(i)+\beta i)=\prod_i(\gamma+f(i)+\beta\sigma(i))$; by unique factorisation of polynomials in $\FF[\beta,\gamma]$ the multisets of linear factors $\{(f(i),i)\}$ and $\{(f(i),\sigma(i))\}$ coincide, i.e.\ $\{(i,f(i))\}=\{(\sigma(i),f(i))\}$, which means $f(i)=f(\sigma(i))$ exactly (as $\sigma$ is a permutation). If they are unequal as polynomials, Schwartz--Zippel shows the identity holds for random $(\beta,\gamma)$ with probability $\le3n/|\FF|$.
\end{proof}
\begin{example}
Three cells with values $f=(3,9,3)$ and the copy constraint ``cell 1 equals cell 3'', i.e.\ $\sigma=(1\ 3)$ ($\sigma(1)=3,\sigma(3)=1,\sigma(2)=2$). Take $\beta=2,\gamma=1$. LHS $=(3+2+1)(9+4+1)(3+6+1)=6\cdot14\cdot10=840$; RHS $=(3+6+1)(9+4+1)(3+2+1)=10\cdot14\cdot6=840$. Equal. If the prover cheats with $f=(3,9,4)$: LHS $=6\cdot14\cdot11=924$, RHS $=(3+6+1)(14)(4+2+1)=980\ne924$ --- detected.
\end{example}
To verify a product over $H$ with a low-degree identity, the prover defines the accumulator $Z(\omega^0)=1$ and
\[Z(\omega^{i+1})=Z(\omega^i)\prod_{\mathsf{col}\in\{a,b,c\}}\frac{w_{\mathsf{col},i}+\beta\,\mathrm{id}_{\mathsf{col}}(\omega^i)+\gamma}{w_{\mathsf{col},i}+\beta\,S_{\sigma,\mathsf{col}}(\omega^i)+\gamma},\]
where $w_{\mathsf{col},i}$ is the wire value in column $\mathsf{col}$ of row $i$, the \emph{identity labels} are $\mathrm{id}_a(\omega^i)=\omega^i$, $\mathrm{id}_b(\omega^i)=k_1\omega^i$, $\mathrm{id}_c(\omega^i)=k_2\omega^i$ with $k_1,k_2$ chosen so that $H,k_1H,k_2H$ are disjoint cosets (they give each of the $3n$ cells a distinct label), and $S_{\sigma,\mathsf{col}}$ are the (public, circuit-specific) permutation polynomials that evaluate to the label of the cell $\sigma$ maps to. The prover requires $Z(\omega^{n})=Z(\omega^0)=1$ and checks the recurrence as polynomial identity on $H$.
\subsection{Putting it together}
\begin{algorithm}[label=pr:plonk]{PLONK prover (simplified)}
\begin{enumerate}
\item \textbf{Round 1.} Commit (KZG) to $a(X),b(X),c(X)$ (blinded).
\item \textbf{Round 2.} Challenges $\beta,\gamma$; commit to the accumulator $Z(X)$.
\item \textbf{Round 3.} Challenge $\alpha$. Form the \emph{quotient polynomial} $T(X)=\frac{G(X)+\alpha\,\text{(perm. constraint)}+\alpha^2\,(Z(X)-1)L_0(X)}{Z_H(X)}$. Split $T$ into pieces of degree $<n$ and commit.
\item \textbf{Round 4.} Challenge $\mathfrak z$ (evaluation point). Send the evaluations $\bar a=a(\mathfrak z),\bar b,\bar c,\bar s_{\sigma1},\bar s_{\sigma2},\overline{Z\omega}=Z(\mathfrak z\omega)$.
\item \textbf{Round 5.} Challenge $v$. Send two KZG opening proofs (batched with powers of $v$) for the linearisation polynomial at $\mathfrak z$ and for $Z$ at $\mathfrak z\omega$.
\end{enumerate}
Make non-interactive by Fiat--Shamir.
\end{algorithm}
The verifier recomputes $T(\mathfrak z)$ from the evaluations, checks the quotient identity at the single point $\mathfrak z$ and verifies the two KZG openings with one pairing equation (batched). By Schwartz--Zippel a polynomial identity checked at a random $\mathfrak z\notin H$ implies the identity holds, except with probability $O(n/|\FF|)$. Proof size: 9 group elements + 6 field elements ($\approx600$ bytes); verifier: 2 pairings + $O(\ell)$ group ops.
\begin{center}
\begin{tikzpicture}[>=Stealth,every node/.style={font=\small}]
 \draw[thick] (0,0)--(0,-5.2);\draw[thick] (7,0)--(7,-5.2);
 \node at (0,0.3) {Prover};\node at (7,0.3) {Verifier (or hash)};
 \draw[->] (0,-0.6)--(7,-0.6) node[midway,above] {$[a],[b],[c]$};
 \draw[<-] (0,-1.3)--(7,-1.3) node[midway,above] {$\beta,\gamma$};
 \draw[->] (0,-2.0)--(7,-2.0) node[midway,above] {$[Z]$};
 \draw[<-] (0,-2.7)--(7,-2.7) node[midway,above] {$\alpha$};
 \draw[->] (0,-3.4)--(7,-3.4) node[midway,above] {$[T_{lo}],[T_{mid}],[T_{hi}]$};
 \draw[<-] (0,-4.1)--(7,-4.1) node[midway,above] {$\mathfrak z$};
 \draw[->] (0,-4.8)--(7,-4.8) node[midway,above] {evals at $\mathfrak z$, KZG proofs};
\end{tikzpicture}
\end{center}
\begin{example}[PLONK gate and permutation on Example~\ref{ex:r1cs}]
The rows $(a,b,c)$ are $(3,3,9)$, $(9,3,27)$, $(27,3,30)$, $(30,0,35)$; take $n=4$ and let $H=\{1,\omega,\omega^2,\omega^3\}$ with $\omega$ of order $4$. The gate polynomial at $\omega^2$ (row 3, where $q_L=q_R=1$, $q_O=-1$) gives $27+3-30=0$. \checkmark

The copy constraints partition the twelve cells into cycles:
\begin{itemize}[nosep]
\item $\{a_1,b_1,b_2,b_3\}$ (all equal to $x=3$),
\item $\{c_1,a_2\}$ (value $9$), $\{c_2,a_3\}$ (value $27$), $\{c_3,a_4\}$ (value $30$).
\end{itemize}
The grand product equals $1$ iff every cycle carries a single value. Changing $b_2$ to $4$ would break the gate ($9\cdot4\ne27$) and the cycle of $x$.
\end{example}
\paragraph{Extensions.} \emph{Custom gates} (e.g.\ a degree-5 gate for Poseidon S-boxes), \emph{lookup arguments} (plookup, Halo2 lookups, Lasso/LogUp: prove $f\subseteq T$ for a table $T$, replacing bit decompositions for range checks and bitwise ops), \emph{recursion} (Halo: accumulate IPA openings so a proof can verify a proof).

\section{STARKs}
STARKs (Ben-Sasson, Bentov, Horesh, Riabzev 2018; ``Scalable Transparent ARguments of Knowledge'') eliminate trusted setup and pairings: only hash functions (collision-resistant / random oracle), hence plausibly post-quantum.
\begin{algorithm}[label=pr:stark]{STARK for an AIR}
\textbf{Input:} AIR over trace length $T=2^k$, trace polynomials $f_1..f_w$ over the subgroup $G=\langle g\rangle$, $|G|=T$.
\begin{enumerate}
\item \textbf{Trace commitment.} Prover interpolates each column, evaluates on a larger \emph{LDE domain} $D\supset G$ ($|D|=\rho^{-1}T$ — blow-up e.g.\ $\times4$ or $\times8$), Merkle-commits.
\item \textbf{Constraint composition.} Verifier sends random $\alpha_j$. Prover forms the composition polynomial
\[C(X)=\sum_j\alpha_j\,\frac{P_j\big(f_1(X),..,f_w(X),f_1(gX),..,f_w(gX)\big)}{Z_j(X)}\]
where $Z_j$ is the vanishing polynomial of the set of points where the $j$-th constraint must hold. If the trace is valid, each quotient is a polynomial, hence $C$ has low degree; otherwise, some quotient has a pole and $C$ is far from low degree. Commit to $C$ on $D$.
\item \textbf{OOD/DEEP sampling} (optional) fix consistency between $C$ and $f_i$ at a random point $\mathfrak z$ outside $D$.
\item \textbf{FRI.} Run FRI on $C$ (and the DEEP quotient) to prove low degree (Algorithm~\ref{pr:fri}).
\item \textbf{Queries.} Verifier checks, at random positions $x\in D$, that $C(x)$ equals the constraint combination computed from the opened trace values $f_i(x),f_i(gx)$ (Merkle openings).
\end{enumerate}
\end{algorithm}
\begin{example}[Fibonacci STARK]
Trace $s=(1,1,2,3,5,8,13,21)$, $T=8$. Let $g$ generate a multiplicative subgroup $G$ of order $8$ of the field and let $f$ be the degree-$\le7$ polynomial with $f(g^i)=s_i$. The transition polynomial $R(X)=f(g^2X)-f(gX)-f(X)$ must vanish at $g^0,\dots,g^5$, so
\[Q_{\mathrm{tr}}(X)=\frac{R(X)}{\prod_{i=0}^{5}(X-g^i)}\]
is a polynomial of degree $\le7-6=1$. The boundary quotients $\frac{f(X)-1}{X-g^0}$, $\frac{f(X)-1}{X-g^1}$, $\frac{f(X)-21}{X-g^7}$ have degree $\le6$. The prover Merkle-commits to $f$ on an evaluation domain of size $32$ (blow-up $4$), commits to the random combination of the quotients, and runs FRI. If the prover cheats by changing one entry (say $s_4=6$), then $R$ is nonzero at $g^2$ and $g^3$ (because $s_4$ occurs in two transitions), so the corresponding quotient has a pole there and is far from a low-degree polynomial; FRI rejects with high probability.
\end{example}
\paragraph{Properties.} Transparent (public randomness only); proof size $O(\log^2 T)$ (50--200 KB in practice); prover quasi-linear $O(T\log T)$ (FFTs and hashing); verifier $O(\log^2T)$; plausibly post-quantum. ZK is obtained by randomising the trace (adding random rows and masking polynomials) so that the $O(\lambda)$ opened values reveal nothing.

\section{Comparison}\label{sec:snarkcomp}
\begin{center}\small
\begin{tabular}{@{}p{1.8cm}p{2.4cm}p{2.3cm}p{2.6cm}p{2.5cm}p{2.5cm}@{}}\toprule
 & Groth16 & PLONK(KZG) & Bulletproofs & STARK & Spartan/Libra\\\midrule
Setup & circuit-specific & universal & none & none & none/universal\\
Proof size & $\approx 128$ B & $\approx 0.5$ KB & $O(\log n)$ (1--2 KB) & $\approx 50$--$200$ KB & $O(\log^2)$ \\
Verifier & 3 pairings & 2 pairings & $O(n)$ & polylog & polylog / $\sqrt n$\\
Prover & $O(n\log n)$ FFT+MSM & $O(n\log n)$ & $O(n)$ MSM & $O(n\log n)$ hashes & $O(n)$\\
Assumption & GGM/AGM (pairings) & AGM, $q$-SDH & DL & Hash (PQ) & DL / hash\\\bottomrule
\end{tabular}
\end{center}

\section{Recursion and folding}
A SNARK verifier is itself a small circuit, so a prover can prove ``I verified the previous proof''. This \emph{proof-carrying data} / recursion approach yields incrementally verifiable computation (IVC): constant-size proof for an unbounded computation. Folding schemes (Nova, SuperNova, HyperNova, ProtoStar) avoid verifying a full SNARK at each step: two R1CS instances $(u_1,w_1)$, $(u_2,w_2)$ are linearly combined, $u=u_1+ru_2$ with random $r$, into one relaxed instance; the cross-term $E$ is committed; only the final folded instance is proven.

\section*{Exercises}
\begin{exercise}
Re-derive the Groth16 verification equation starting from the requirement that the verifier tests $AB=\alpha\beta+\ldots$; explain the role of $\gamma$ and $\delta$ (they separate public from private parts and prevent mix-and-match attacks).
\end{exercise}
\begin{exercise}
Recompute Example~\ref{ex:groth16} with a different $\tau$ and randomisers; check the pairing equation. What happens if you change $a_5=30$ to $31$ (and keep $h$ as computed from the honest witness)?
\end{exercise}
\begin{exercise}
Prove Lemma~\ref{lem:perm} carefully, including the soundness error $\le 3n/|\FF|$.
\end{exercise}
\begin{exercise}\label{ex:starkq}
Compute, for a STARK over a field with $\lambda=100$ bits of security, blow-up factor $8$ (rate $1/8$) and proximity parameter $\delta\approx0.4$, the number of FRI queries needed (use $s\approx\lambda/\log_2\frac1{1-\delta}$).
\end{exercise}
\begin{exercise}
Show that Groth16's $(A,B)$ is uniform (zero-knowledge) when $r,s$ are uniform, and explain what leaks if $r=s=0$.
\end{exercise}

\chapter{Complexity-Theoretic Foundations and Limits}
\label{ch:theory}
\section{Classes of zero-knowledge languages}
Let $\PZK,\SZK,\CZK$ be the classes of languages with perfect, statistical, and computational zero-knowledge \emph{proof} systems (against malicious verifiers); $\HVSZK$ restricts to honest verifiers.
\begin{theorem}[Landscape]\label{thm:landscape}
\begin{enumerate}
\item $\BPP\subseteq\PZK\subseteq\SZK\subseteq\CZK\subseteq\IP=\PSPACE$.
\item $\GI,\ \GNI,\ \QR,\ \QNR\in\PZK$/$\SZK$ (Sections~\ref{sec:gi}--\ref{sec:qr}, GMW, GMR).
\item (Goldreich--Micali--Wigderson; Impagliazzo--Yung; Ben-Or et al.) If one-way functions exist then $\CZK=\IP=\PSPACE$. Conversely, if $\CZK$ contains a language outside $\BPP$ then (Ostrovsky--Wigderson) a weak form of one-way function exists.
\item $\HVSZK=\SZK$ (Goldreich--Sahai--Vadhan).
\item $\SZK\subseteq\AM\cap\coAM$ (Fortnow; Aiello--H{\aa}stad). Hence an $\NP$-complete language is not in $\SZK$ unless the polynomial hierarchy collapses.
\item $\SZK$ is closed under complement (Okamoto) and has the complete problems \textsc{Statistical Difference} and \textsc{Entropy Difference} (Sahai--Vadhan).
\end{enumerate}
\end{theorem}
\begin{center}
\begin{tikzpicture}[every node/.style={font=\small}]
 \draw[rounded corners,fill=gray!6] (0,0) rectangle (11,5.2);\node at (5.5,4.9) {$\IP=\PSPACE=\CZK$ (under OWF)};
 \draw[rounded corners,fill=blue!6] (0.4,0.4) rectangle (7.2,4.4);\node at (3.8,4.1) {$\AM$};
 \draw[rounded corners,fill=green!8] (0.8,0.8) rectangle (6.4,3.4);\node at (3.6,3.1) {$\SZK\subseteq\AM\cap\coAM$};
 \draw[rounded corners,fill=yellow!15] (1.2,1.2) rectangle (4.2,2.6);\node at (2.7,2.3) {$\PZK$};
 \node at (2.7,1.6) {$\GI,\QR$};
 \node at (5.6,1.7) {\textsc{SD},\textsc{ED}};
 \node[align=center] at (8.9,2.6) {3-COL, SAT\\ (only under\\ assumptions)};
\end{tikzpicture}
\end{center}
\begin{theorem}[$\SZK\subseteq\AM\cap\coAM$, sketch]\label{thm:szkam}
$\SZK\subseteq\AM$: it is enough to put \textsc{Entropy Difference} in $\AM$. Statistical zero-knowledge implies that the simulator's output on $x\in L$ has some specific entropy relation with the output on $x\notin L$. The $\AM$ verifier uses the Goldwasser--Sipser set lower bound protocol to certify that the entropy of the simulator's output distribution is above a threshold. Since $\SZK$ is closed under complement, $\SZK\subseteq\coAM$ as well.
\end{theorem}
Interpretation: statistical ZK proofs are for problems of ``intermediate'' complexity (factoring-like, lattice-like), not for $\NP$-hard problems. Statistical ZK \emph{arguments} for $\NP$ do exist (from collision-resistant hash functions: Brassard--Chaum--Cr\'epeau; Naor--Ostrovsky--Venkatesan--Yung), since computationally bounded provers cannot exploit binding.

\section{Sahai--Vadhan problems: the SD problem}
\begin{definition}[Statistical Difference, $\mathrm{SD}_{\alpha,\beta}$]
Input: circuits $C_0,C_1$ sampling distributions $D_0,D_1$. YES if $\SD(D_0,D_1)\ge\beta$; NO if $\SD(D_0,D_1)\le\alpha$, with $\alpha^2<\beta$ (e.g.\ $\alpha=1/3$, $\beta=2/3$); the condition $\alpha^2<\beta$ is exactly what Sahai--Vadhan's \emph{polarization lemma} needs in order to amplify the gap.
\end{definition}
\begin{theorem}[Sahai--Vadhan]
$\mathrm{SD}$ is $\SZK$-complete; for $\GI$, $\GNI$, $\QR$ this is seen by taking $D_b$ = random isomorphic copies of $G_b$.
\end{theorem}
\begin{example}[$\GNI\le\mathrm{SD}$]
Let $D_b$ be the distribution of $\sigma(G_b)$, $\sigma\getsr S_n$. If $G_0\cong G_1$ then $D_0=D_1$ ($\SD=0$); if not then their supports are disjoint ($\SD=1$). So $(G_0,G_1)\in\GNI$ iff $\SD(D_0,D_1)=1$. Conversely every $\SD$ instance has a GNI-like ZK proof: the verifier flips a coin $b$ and sends a sample from $D_b$; the prover guesses $b$; its success probability is $\frac12+\frac12\SD(D_0,D_1)$. This is exactly the protocol of Section~\ref{sec:gni}, generalised.
\end{example}

\section{Impossibility results}
\begin{theorem}[Goldreich--Oren 1994]
If $L$ has a ZK proof system (with auxiliary input) then the verifier must be randomized and interactive: non-interactive ZK in the plain model is possible only for $L\in\BPP$.
\end{theorem}
\begin{proof}
In a one-message protocol the verifier's decision is a function of $(x,\pi)$ and its coins. Run $\Sim(x)$ (the simulator for the honest verifier): by ZK, on $x\in L$ the simulator outputs views indistinguishable from accepting ones; on $x\notin L$, soundness says no $\pi$ is accepted except w.p. $s$, so the simulator's output (which can be checked by running $\Ver$) is rejected w.h.p. Hence $L\in\BPP$ is decided by ``simulate and verify''.
\end{proof}
\begin{theorem}[Goldreich--Krawczyk 1996]
If a language has a 3-round (auxiliary-input) black-box ZK proof system, then $L\in\BPP$. If it has a constant-round public-coin black-box ZK proof system, then $L\in\BPP$.
\end{theorem}
\begin{intuition}
A black-box simulator can only rewind. With a verifier that answers with $c=\mathrm{PRF}(a)$, a pseudorandom function of the first message, rewinding gains nothing: the simulator must predict the challenge. A simulator that produces a convincing transcript for a (pseudo)random challenge must, by an extractor-like argument, have been able to answer \emph{two} challenges, whence it decides $L$.
\end{intuition}
\begin{theorem}[Barak 2001]
Non-black-box simulation circumvents these: there is a constant-round public-coin ZK argument for $\NP$ with negligible soundness error. The simulator uses the \emph{code} of $\Ver^*$: the protocol includes a branch that is provable only by knowing a short commitment to a program that predicts the verifier's next message; the simulator commits to the verifier's own code.
\end{theorem}
\begin{theorem}[Gentry--Wichs 2011]
No SNARG for $\NP$ has a security reduction to a falsifiable assumption in a black-box way; knowledge assumptions are unavoidable for succinct non-interactive \emph{arguments} (hence SNARK's extractability assumptions and ROM-based STARKs).
\end{theorem}
\begin{theorem}[Composition: parallel and concurrent]
The sequential composition theorem (Thm.~\ref{thm:seqcomp}) holds, but parallel composition of ZK may fail (Goldreich--Krawczyk construct, under standard assumptions, a ZK proof system whose two-fold parallel repetition is not ZK). Concurrent ZK requires $\tilde\Omega(\log\lambda)$ rounds for black-box simulation (Canetti--Kilian--Petrank--Rosen) and is achievable (Richardson--Kilian, Prabhakaran--Rosen--Sahai).
\end{theorem}

\section{PCPs, IOPs and the origins of succinctness}
The sum-check protocol (Ch.~\ref{ch:sumcheck}) led to the PCP theorem: an $\NP$ statement has a proof that can be verified by reading only $O(1)$ bits of it.
\begin{definition}[PCP]
$L\in\PCP[r(n),q(n)]$ if there is a polynomial-time verifier that tosses $r(n)$ coins, queries $q(n)$ positions of an oracle $\pi$, accepts $x\in L$ with some $\pi$ with probability $1$ and rejects $x\notin L$ with probability $\ge1/2$ for all $\pi$.
\end{definition}
\begin{theorem}[ALMSS 1992]
$\NP=\PCP[O(\log n),O(1)]$.
\end{theorem}
\begin{center}
\begin{tikzpicture}[>=Stealth,every node/.style={font=\small}]
 \node[draw,rounded corners,fill=blue!6,text width=2.9cm,align=center] (a) at (0,0) {PCP\\(long proof $\pi$,\\ query a few bits)};
 \node[draw,rounded corners,fill=blue!6,text width=2.9cm,align=center] (b) at (4.2,0) {Merkle-commit $\pi$,\\ open queried bits};
 \node[draw,rounded corners,fill=blue!6,text width=2.9cm,align=center] (c) at (8.4,0) {Fiat--Shamir:\\ queries = hash\\of root};
 \node[draw,rounded corners,fill=green!8,text width=2.0cm,align=center] (d) at (12.0,0) {Kilian / Micali\\ SNARG};
 \draw[->] (a)--(b);\draw[->] (b)--(c);\draw[->] (c)--(d);
\end{tikzpicture}
\end{center}
\begin{theorem}[Kilian 1992; Micali 2000]
PCP + Merkle commitments (collision-resistant hash) gives a 4-message succinct argument for $\NP$ (Kilian), which Fiat--Shamir makes non-interactive in the random oracle model (Micali's computationally sound proofs).
\end{theorem}
\emph{Interactive oracle proofs (IOPs)} (Ben-Sasson--Chiesa--Spooner; Reingold--Rothblum--Rothblum) generalise PCPs and interactive proofs: in each round the prover sends an oracle (a long string) and the verifier sends randomness; the verifier queries only a few symbols. STARKs, Aurora, Fractal, Plonky2, Ligero are IOP + hash-based compilations; Marlin/PLONK are \emph{polynomial IOPs} compiled using a polynomial commitment scheme.

\begin{center}
\begin{tikzpicture}[>=Stealth,every node/.style={font=\small}]
 \node[draw,rounded corners,fill=yellow!10,text width=3.2cm,align=center] (io) at (0,0) {\textbf{Information-theoretic}\\ polynomial IOP / IOP\\ (oracles)};
 \node[draw,rounded corners,fill=blue!6,text width=3.2cm,align=center] (pc) at (5.2,0) {\textbf{Cryptographic compiler}\\ KZG / FRI / IPA\\ (commitments)};
 \node[draw,rounded corners,fill=green!8,text width=3.2cm,align=center] (fs) at (10.4,0) {\textbf{Fiat--Shamir}\\ $\Rightarrow$ non-interactive\\ SNARK/STARK};
 \draw[->] (io)--(pc);\draw[->] (pc)--(fs);
\end{tikzpicture}
\end{center}
This modular ``IOP $\to$ PCS $\to$ FS'' pipeline explains the modern landscape: choosing a different PCS changes setup, proof size and assumptions (see the comparison table in Section~\ref{sec:snarkcomp}) without touching the arithmetisation.

\section{Quantum considerations}
\begin{itemize}
\item \emph{Soundness against quantum provers.} Rewinding is problematic for quantum adversaries (no-cloning): classical proofs of knowledge via rewinding do not carry over (Unruh's quantum rewinding, quantum PoK).
\item \emph{Fiat--Shamir in the quantum random oracle model} (QROM): Don et al., Liu--Zhandry give tight-ish reductions for some protocols; Dilithium/Falcon-like signatures (Fiat--Shamir with aborts) use lattice $\Sigma$-protocols.
\item \emph{Post-quantum ZK:} lattice-based (Lyubashevsky; LaBRADOR), hash-based (STARKs, Ligero++, Brakedown), MPC-in-the-head (Ishai--Kushilevitz--Ostrovsky--Sahai: simulate an MPC protocol among $n$ virtual parties, open $t$ views; used in Picnic and Banquet).
\item \emph{Quantum advantage:} $\mathsf{MIP}^*=\mathsf{RE}$ (Ji--Natarajan--Vidick--Wright--Yuen 2020) shows entangled provers can prove astonishingly hard statements; ZK analogues exist.
\end{itemize}

\section*{Exercises}
\begin{exercise}
Prove that $\PZK\subseteq\SZK\subseteq\CZK$ and that $\BPP\subseteq\PZK$.
\end{exercise}
\begin{exercise}
Show that $\GNI\le\mathrm{SD}$ and $\QNR\le\mathrm{SD}$ by exhibiting the sampling circuits.
\end{exercise}
\begin{exercise}
Give the details of the argument that a one-message ZK proof system in the plain model implies $L\in\BPP$ (Goldreich--Oren). Where does the argument fail in the CRS model?
\end{exercise}
\begin{exercise}
State precisely why the Kilian PCP-to-argument transformation needs a \emph{commitment with local opening} and construct one from a Merkle tree; analyse the soundness error when the verifier makes $q$ queries.
\end{exercise}

\chapter{Applications, Practical Considerations and Further Reading}
\label{ch:apps}
\section{Applications}
\begin{description}
\item[Identification and authentication.] Schnorr/FFS identification; passwordless authentication (OPAQUE with ZK password proofs); proving age or residency from a government credential without revealing the credential (ISO 18013-5 mDL + zk).
\item[Signatures.] Schnorr/EdDSA, BLS (not FS-based), Dilithium; ring signatures and group signatures (CDS OR-proofs, linkable ring signatures in Monero); blind signatures.
\item[Anonymous credentials.] Camenisch--Lysyanskaya, U-Prove, BBS+ signatures with selective disclosure: the holder proves in ZK possession of a signed credential with attribute predicates.
\item[Private payments.] Zerocash/Zcash: a note commitment $\mathsf{cm}=\mathsf{Com}(v,\mathsf{pk};\rho)$ is inserted into a Merkle tree; spending proves (in ZK) knowledge of the opening, Merkle membership, a correctly derived nullifier $\mathsf{nf}$, and balance $\sum v_{in}=\sum v_{out}$. Confidential transactions in Monero use Pedersen commitments and Bulletproofs for range.
\item[Verifiable computation and blockchain scaling.] zk-rollups (zkSync, StarkNet, Polygon zkEVM, Scroll, Linea): a prover executes thousands of transactions off-chain and posts a SNARK/STARK of the state transition; L1 verifies one proof. zkVMs (RISC Zero, SP1, Jolt) prove arbitrary RISC-V execution.
\item[Verifiable machine learning, data and databases.] zkML proofs of inference; proofs of SQL queries; proofs of solvency of an exchange (Merkle sum tree + ZK range proofs).
\item[Secure multi-party computation.] GMW compiler: each party attaches a ZK proof of correct behaviour. Also MPC-in-the-head signatures and ZK proofs.
\item[Voting and auctions.] Verifiable shuffles (Neff/Groth), proofs of correct ElGamal decryption (Chaum--Pedersen), sealed-bid auctions.
\end{description}

\section{Choosing a proof system}
\begin{center}\small
\begin{tabular}{@{}p{4.2cm}p{9.5cm}@{}}\toprule
Requirement & Suggested tool\\\midrule
Prove discrete-log relations, few statements & $\Sigma$-protocols + Fiat--Shamir\\
Range proofs, no trusted setup, small statement & Bulletproofs\\
Smallest proof, on-chain verification, fixed circuit & Groth16\\
Evolving circuits, universal setup & PLONK/Halo2 with KZG\\
No trusted setup, large computation, post-quantum & STARKs (FRI)\\
Proving long iterative computations & Folding (Nova) / recursion\\
Maximum prover speed, structured computations & sum-check based (Spartan, Jolt, Lasso) \\\bottomrule
\end{tabular}
\end{center}

\section{Implementation pitfalls (a checklist)}
\begin{enumerate}
\item Hash \emph{everything} in Fiat--Shamir (statement, public parameters, all prior messages); use domain separation.
\item Check group membership/subgroup of every received point (small-subgroup and invalid-curve attacks); reject the identity where inappropriate.
\item Use constant-time scalar operations; generate nonces deterministically or from a CSPRNG.
\item Under-constrained circuits are the dominant real-world bug: every witness variable must be pinned by constraints (range checks, boolean checks, non-zero checks via inverse witnesses $x\cdot x^{-1}=1$). Audit with formal tools (Circomspect, Picus, ZKAP).
\item Trusted setup hygiene: verify transcripts of the ceremony; sample $\tau$ in an MPC; use a beacon.
\item Soundness bits: with $|\FF|\approx2^{64}$, Schwartz--Zippel error $d/|\FF|$ is not negligible; sample challenges from an extension field.
\item Malleability and replay: bind proofs to application context (contract address, nonce) by including them as public inputs.
\end{enumerate}

\section{Where to go from here}
\begin{itemize}
\item Foundational texts: Goldreich, \emph{Foundations of Cryptography Vol.\ 1} (Ch.\ 4); Goldwasser--Micali--Rackoff; Arora--Barak Ch.\ 8--11; Vadhan's thesis on SZK.
\item $\Sigma$-protocols and applications: Damgård's lecture notes ``On $\Sigma$-protocols''; Boneh--Shoup, \emph{A Graduate Course in Applied Cryptography} (Ch.\ 19, 20).
\item SNARKs/STARKs: Thaler, \emph{Proofs, Arguments, and Zero-Knowledge}; Ben-Sasson et al.\ STARK papers; Justin Thaler's survey; Vitalik Buterin's QAP/PLONK blog posts; Groth16 paper; the PLONK paper; ZK Hack lectures.
\item Open problems: tight Fiat--Shamir security in the standard model; transparent SNARKs with $O(1)$ verification; post-quantum SNARKs with short proofs; ZK with optimal prover overhead (linear time, small constant); concurrent security of recursion; formal verification of circuits.
\end{itemize}

\section*{Capstone exercises}
\begin{exercise}
Design a Schnorr-based protocol that lets a prover show that a Pedersen commitment $C=g^vh^r$ hides a value $v\in\{3,5,7\}$. Give completeness, special soundness and SHVZK, and run it with the toy group.
\end{exercise}
\begin{exercise}
Implement (in Python, using the toy group and a prime field of 17 elements) Schnorr, Chaum--Pedersen, the OR-proof, sum-check (Example~\ref{ex:sumcheck}), the KZG ``in the exponent'' check and one FRI fold; test the soundness of each by simulating cheating provers.
\end{exercise}
\begin{exercise}
Write an R1CS for the statement ``I know a preimage $s$ with $s^2+3s+2=0$ over $\FF_{101}$'', convert to QAP (with points $1,2,\ldots$), and produce the Groth16 ``exponent-level'' proof with your own toxic waste.
\end{exercise}
\begin{exercise}
Prove that Protocol~\ref{pr:or} (OR-proof) is not zero-knowledge against a malicious verifier unless the challenge is committed in advance, by exhibiting the verifier that biases challenges.
\end{exercise}
\begin{exercise}
Survey: pick a deployed system (Zcash Sapling/Orchard, a zk-rollup, or a zkVM) and identify: the arithmetisation, the PCS, the setup, the Fiat--Shamir hashes, and the assumptions; compare with the comparison table in Section~\ref{sec:snarkcomp}.
\end{exercise}


\appendix
\chapter{Notation and Toy-Parameter Reference}
\section{Notation}
\begin{center}\small
\begin{tabular}{@{}p{3.3cm}p{10.5cm}@{}}\toprule
Symbol & Meaning\\\midrule
$\lambda$, $\negl$, $\poly$ & security parameter; negligible / polynomial functions\\
$x\getsr S$ & sample $x$ uniformly from the finite set $S$\\
$\pair{\Prv(w)}{\Ver}(x)$ & output of the verifier in the interaction on common input $x$, prover witness $w$\\
$\view^{\Prv(w)}_{\Ver^*(z)}(x)$ & view of $\Ver^*$ (random tape and received messages), auxiliary input $z$\\
$\Sim$, $\Ext$ & simulator, knowledge extractor\\
$\equiv,\ \approx_s,\ \approx_c$ & perfectly / statistically / computationally indistinguishable\\
$R$, $L_R$ & $\NP$ relation and its language\\
$\GG,g,q$ & cyclic group, generator, prime order\\
$[a]_1,[a]_2,[a]_T$ & $g_1^a,g_2^a,e(g_1,g_2)^a$\\
$\tilde f$ & multilinear extension of $f:\{0,1\}^v\to\FF$\\
$u_i,v_i,w_i,t,h$ & QAP polynomials, vanishing polynomial, quotient\\
$\ra{m}$, $\la{m}$ & message $m$ sent prover$\to$verifier / verifier$\to$prover\\\bottomrule
\end{tabular}
\end{center}
\section{Toy group and fields used in the examples}
\begin{itemize}
\item $\GG=\langle2\rangle\subset\Zs{23}$, order $q=11$, $g=2$, $h=3=2^8$. Discrete-log table: $2^k$ for $k=0..10$ is $1,2,4,8,16,9,18,13,3,6,12$.
\item Inverses mod $11$: $2\inv=6,\ 3\inv=4,\ 4\inv=3,\ 5\inv=9,\ 6\inv=2,\ 7\inv=8,\ 8\inv=7,\ 9\inv=5,\ 10\inv=10$.
\item Inverses mod $23$ used: $18\inv=9$, $4\inv=6$, $13\inv=16$.
\item $N=77=7\cdot11$: $23\inv=67$; $10^2=23$; secret roots $s=10,4$ with $v=23,16$.
\item Prime fields: $\FF_{17}$ (FRI, Schwartz--Zippel), $\FF_{101}$ (Groth16 example), $\ZZ_{11}$ (IPA, KZG exponents).
\end{itemize}

\chapter{Hints and Answers to Selected Exercises}
\begin{description}
\item[Exercise~\ref{ex:detv}] If $\Ver$ is deterministic, a prover's best strategy against it is a fixed sequence of answers; the witness is the full list of prover messages, and the $\NP$ verifier replays $\Ver$ on it.
\item[Exercise~\ref{ex:extractor} / Theorem~\ref{thm:ss2pok}] Mimic the proof of the extractor theorem with $N=2$: sample until accept, rewind once and sample the other challenge; expected cost $O(1/\varepsilon)$.
\item[Exercise~\ref{ex:rounds3col}] Need $(1-1/1000)^N\le2^{-80}$, i.e.\ $N\ge80\ln2/(-\ln0.999)\approx5.54\times10^4$ sequential rounds of Protocol~\ref{pr:col3} --- which illustrates why Blum's protocol (error $1/2$ per round, $80$ rounds) is far more efficient.
\item[Exercise~\ref{ex:nonce} (nonce reuse)] Two Schnorr transcripts $(t,c,s)$, $(t',c',s')$ with the same $r$ have $t=t'$ and $s-s'=(c-c')x$, so $x=(s-s')(c-c')^{-1}$.
\item[Exercise~\ref{ex:weakfs} (weak FS)] The attacker picks $t$, computes $c=H(t)$, picks any $s$ and sets $y=(g^st^{-1})^{1/c}$; then $(t,s)$ is accepted for the statement $y$, whose discrete logarithm the attacker does not know.
\item[Exercise~\ref{ex:forkloss} (forking loss)] A forger with running time $T_F$ and success probability $\varepsilon$ gives a DL solver with time $\approx2T_F$ and success $\approx\varepsilon^2/Q$. If solving DL needs $2^{128}$ time at success probability $\approx1$, this only yields $T_F/\varepsilon^2\gtrsim2^{127}/Q$; for $Q=2^{60}$ and $\varepsilon\approx1$ the reduction guarantees only about $67$ bits.
\item[Exercise~\ref{ex:sumcheckrun} (sum-check)] $g=x_1x_2+3x_2$, $H=0+3+0+4=7$. $g_1(X)=X+3$ with $g_1(0)+g_1(1)=7$; $r_1=4$, $g_1(4)=7$. $g_2(Y)=g(4,Y)=7Y$, $g_2(0)+g_2(1)=7$; $r_2=7$, $g_2(7)=49\equiv5$; final $g(4,7)=28+21=49\equiv5$. A prover claiming $H'=8$ must send $g_1^*=g_1+d$ with $d(0)+d(1)=1$; if $d(X)=aX+b$ with $a\ne0$ it passes round~1 only when $r_1=-b/a$, probability $1/11$.
\item[Exercise~\ref{ex:r1csex}] For $y=x^4$: $x\cdot x=v_1$, $v_1\cdot v_1=y$. Multiplexer: $b(1-b)=0$ is the constraint $b\cdot(1-b)=0$, and $z-c_0=b\cdot(c_1-c_0)$ is $b\cdot(c_1-c_0)=z-c_0$.
\item[Exercise~\ref{ex:foldcheck}] $f(4)=64+32+12+4=112\equiv10$, $f(13)=f(-4)=-64+32-12+4=-40\equiv11$, so $\frac{10+11}2=\frac{21}2\equiv4\cdot9=36\equiv2$ ($2\inv=9$) and $\frac{10-11}{2\cdot4}=\frac{-1}8\equiv-15\equiv2$ ($8\inv=15$); $f^{(1)}(16)=2+3\cdot2=8$. Directly: $5\cdot16+13=93\equiv8$. \checkmark
\item[Exercise~\ref{ex:starkq}] With $\delta=0.4$: $\log_2(1/0.6)\approx0.737$, so $s\approx100/0.737\approx136$ queries.
\end{description}

\begin{thebibliography}{99}
\bibitem{GMR89} S.\ Goldwasser, S.\ Micali, C.\ Rackoff. The knowledge complexity of interactive proof systems. \emph{SIAM J.\ Comput.} 18(1):186--208, 1989 (STOC 1985).
\bibitem{GMW91} O.\ Goldreich, S.\ Micali, A.\ Wigderson. Proofs that yield nothing but their validity or all languages in NP have zero-knowledge proof systems. \emph{J.\ ACM} 38(3):690--728, 1991.
\bibitem{FFS88} U.\ Feige, A.\ Fiat, A.\ Shamir. Zero-knowledge proofs of identity. \emph{J.\ Cryptology} 1(2):77--94, 1988.
\bibitem{FS86} A.\ Fiat, A.\ Shamir. How to prove yourself. CRYPTO 1986.
\bibitem{Sch91} C.-P.\ Schnorr. Efficient signature generation by smart cards. \emph{J.\ Cryptology} 4(3):161--174, 1991.
\bibitem{CP92} D.\ Chaum, T.\ Pedersen. Wallet databases with observers. CRYPTO 1992.
\bibitem{CDS94} R.\ Cramer, I.\ Damgård, B.\ Schoenmakers. Proofs of partial knowledge and simplified design of witness hiding protocols. CRYPTO 1994.
\bibitem{Ped91} T.\ Pedersen. Non-interactive and information-theoretic secure verifiable secret sharing. CRYPTO 1991.
\bibitem{PS96} D.\ Pointcheval, J.\ Stern. Security proofs for signature schemes. EUROCRYPT 1996.
\bibitem{BN06} M.\ Bellare, G.\ Neven. Multi-signatures in the plain public-key model and a general forking lemma. CCS 2006.
\bibitem{Sha92} A.\ Shamir. IP $=$ PSPACE. \emph{J.\ ACM} 39(4):869--877, 1992.
\bibitem{LFKN92} C.\ Lund, L.\ Fortnow, H.\ Karloff, N.\ Nisan. Algebraic methods for interactive proof systems. \emph{J.\ ACM} 39(4):859--868, 1992.
\bibitem{GKR15} S.\ Goldwasser, Y.\ Kalai, G.\ Rothblum. Delegating computation: interactive proofs for muggles. \emph{J.\ ACM} 62(4), 2015.
\bibitem{Tha13} J.\ Thaler. Time-optimal interactive proofs for circuit evaluation. CRYPTO 2013.
\bibitem{GGPR13} R.\ Gennaro, C.\ Gentry, B.\ Parno, M.\ Raykova. Quadratic span programs and succinct NIZKs without PCPs. EUROCRYPT 2013.
\bibitem{Gro16} J.\ Groth. On the size of pairing-based non-interactive arguments. EUROCRYPT 2016.
\bibitem{GWC19} A.\ Gabizon, Z.\ Williamson, O.\ Ciobotaru. PLONK: Permutations over Lagrange-bases for oecumenical noninteractive arguments of knowledge. ePrint 2019/953.
\bibitem{KZG10} A.\ Kate, G.\ Zaverucha, I.\ Goldberg. Constant-size commitments to polynomials and their applications. ASIACRYPT 2010.
\bibitem{BBHR18} E.\ Ben-Sasson, I.\ Bentov, Y.\ Horesh, M.\ Riabzev. Fast Reed--Solomon interactive oracle proofs of proximity. ICALP 2018.
\bibitem{BBHR19} E.\ Ben-Sasson, I.\ Bentov, Y.\ Horesh, M.\ Riabzev. Scalable zero knowledge with no trusted setup. CRYPTO 2019.
\bibitem{BBBPWM18} B.\ B\"unz, J.\ Bootle, D.\ Boneh, A.\ Poelstra, P.\ Wuille, G.\ Maxwell. Bulletproofs: short proofs for confidential transactions and more. IEEE S\&P 2018.
\bibitem{BCS16} E.\ Ben-Sasson, A.\ Chiesa, N.\ Spooner. Interactive oracle proofs. TCC 2016.
\bibitem{Kil92} J.\ Kilian. A note on efficient zero-knowledge proofs and arguments. STOC 1992.
\bibitem{Mic00} S.\ Micali. Computationally sound proofs. \emph{SIAM J.\ Comput.} 30(4), 2000.
\bibitem{BFM88} M.\ Blum, P.\ Feldman, S.\ Micali. Non-interactive zero-knowledge and its applications. STOC 1988.
\bibitem{FLS99} U.\ Feige, D.\ Lapidot, A.\ Shamir. Multiple non-interactive zero knowledge proofs under general assumptions. \emph{SIAM J.\ Comput.} 29(1), 1999.
\bibitem{GOS06} J.\ Groth, R.\ Ostrovsky, A.\ Sahai. Perfect non-interactive zero knowledge for NP. EUROCRYPT 2006.
\bibitem{GS08} J.\ Groth, A.\ Sahai. Efficient non-interactive proof systems for bilinear groups. EUROCRYPT 2008.
\bibitem{Nao91} M.\ Naor. Bit commitment using pseudorandomness. \emph{J.\ Cryptology} 4(2), 1991.
\bibitem{HILL99} J.\ H{\aa}stad, R.\ Impagliazzo, L.\ Levin, M.\ Luby. A pseudorandom generator from any one-way function. \emph{SIAM J.\ Comput.} 28(4), 1999.
\bibitem{GK96} O.\ Goldreich, H.\ Krawczyk. On the composition of zero-knowledge proof systems. \emph{SIAM J.\ Comput.} 25(1), 1996.
\bibitem{GO94} O.\ Goldreich, Y.\ Oren. Definitions and properties of zero-knowledge proof systems. \emph{J.\ Cryptology} 7(1), 1994.
\bibitem{Bar01} B.\ Barak. How to go beyond the black-box simulation barrier. FOCS 2001.
\bibitem{BIN97} M.\ Bellare, R.\ Impagliazzo, M.\ Naor. Does parallel repetition lower the error in computationally sound protocols? FOCS 1997.
\bibitem{GK03} S.\ Goldwasser, Y.\ Kalai. On the (in)security of the Fiat--Shamir paradigm. FOCS 2003.
\bibitem{CGH04} R.\ Canetti, O.\ Goldreich, S.\ Halevi. The random oracle methodology, revisited. \emph{J.\ ACM} 51(4), 2004.
\bibitem{SV03} A.\ Sahai, S.\ Vadhan. A complete problem for statistical zero knowledge. \emph{J.\ ACM} 50(2), 2003.
\bibitem{GW11} C.\ Gentry, D.\ Wichs. Separating succinct non-interactive arguments from all falsifiable assumptions. STOC 2011.
\bibitem{CCHLRR19} R.\ Canetti et al. Fiat--Shamir: from practice to theory. STOC 2019.
\bibitem{PS19} C.\ Peikert, S.\ Shiehian. Noninteractive zero knowledge for NP from (plain) LWE. CRYPTO 2019.
\bibitem{Mau09} U.\ Maurer. Unifying zero-knowledge proofs of knowledge. AFRICACRYPT 2009.
\bibitem{Tha22} J.\ Thaler. \emph{Proofs, Arguments, and Zero-Knowledge}. Foundations and Trends in Privacy and Security, 2022.
\bibitem{Gol01} O.\ Goldreich. \emph{Foundations of Cryptography, Vol.\ 1: Basic Tools}. Cambridge University Press, 2001.
\bibitem{BS23} D.\ Boneh, V.\ Shoup. \emph{A Graduate Course in Applied Cryptography}. 2023.
\bibitem{ALMSS98} S.\ Arora, C.\ Lund, R.\ Motwani, M.\ Sudan, M.\ Szegedy. Proof verification and the hardness of approximation problems. \emph{J.\ ACM} 45(3), 1998.
\bibitem{BCCT} N.\ Bitansky, R.\ Canetti, A.\ Chiesa, E.\ Tromer. From extractable collision resistance to succinct non-interactive arguments of knowledge. ITCS 2012.
\bibitem{Kothapalli} A.\ Kothapalli, S.\ Setty, I.\ Tzialla. Nova: recursive zero-knowledge arguments from folding schemes. CRYPTO 2022.
\end{thebibliography}
\end{document}
